% ===========================================================================
% informe.tex - el informe de la campaña de calibración.
%
% ESTE ARCHIVO ES PARA EDITAR A MANO. Nada del proyecto lo reescribe.
%
% Los números y las tablas NO se tipean: salen de generado/, que reescribe
%     python analisis/informe.py
% cada vez que se recalculan las calibraciones. Por eso en el texto se lee
% \NCorridas y no un 25 que mañana es otro. La lista completa de macros
% disponibles está al principio de generado/numeros.tex.
%
% Para compilar a mano:  pdflatex informe.tex  (dos veces, por el índice)
% ===========================================================================
\documentclass[11pt,a4paper]{article}
% ---------------------------------------------------------------------------
% preambulo.tex - paquetes y estilo del informe.
%
% Este archivo es EDITABLE: nada del proyecto lo reescribe. Lo que si se
% regenera solo es la carpeta generado/ (tablas y numeros), que la escribe
% `python analisis/informe.py`.
% ---------------------------------------------------------------------------
\usepackage[utf8]{inputenc}
% lmodern va ANTES de fontenc y no es decorativo: sin una fuente escalable,
% microtype aborta la compilacion con "auto expansion is only possible with
% scalable fonts".
\usepackage{lmodern}
\usepackage[T1]{fontenc}
% `es-nodecimaldot` deja el punto decimal como esta escrito. Sin esa opcion,
% babel convierte a coma SOLO dentro de modo matematico, asi que las tablas
% dirian 1,6561 y el texto -donde los numeros entran como macros- 1.6561.
\usepackage[spanish,es-noquoting,es-tabla,es-nodecimaldot]{babel}
\usepackage[a4paper,margin=2.4cm]{geometry}

\usepackage{amsmath}
\usepackage{amssymb}          % \gtrsim y compania
\usepackage{graphicx}
% `float` da el especificador [H]: pone la figura AHI, sin dejarla flotar.
% Se usa en el apéndice, donde las figuras son el contenido y no un apoyo.
\usepackage{float}
\usepackage{booktabs}
\usepackage{longtable}
\usepackage{array}
\usepackage{xcolor}
\usepackage[font=small,labelfont=bf]{caption}
\usepackage{microtype}
\usepackage[hidelinks]{hyperref}

\graphicspath{{figuras/}}

\definecolor{azul}{HTML}{1F4E79}
\definecolor{gris}{HTML}{555F6A}
\definecolor{naranjo}{HTML}{C0661A}

% Diagramas (el flujo del metodo y la cadena de senal). pgfmath, que viene con
% tikz, es ademas lo que calcula k = b/12 en informe.tex.
\usepackage{tikz}
\usetikzlibrary{arrows.meta,positioning,decorations.pathreplacing}

% Etiquetas de procedencia. Van al comienzo del cuerpo de una seccion de
% metodo (no en el titulo, para no ensuciar el indice) y dicen de donde sale
% lo que sigue: del paper de Kong et al., del paper modificado, o propio.
\newcommand{\etiqueta}[2]{\leavevmode{\setlength{\fboxsep}{1.5pt}%
   \colorbox{#1}{\color{white}\scriptsize\sffamily\bfseries\strut #2}}}
\newcommand{\enpaper}{\etiqueta{azul}{DEL PAPER}}
\newcommand{\adaptado}{\etiqueta{gris}{ADAPTADO}}
\newcommand{\propio}{\etiqueta{naranjo}{PROPIO}}

\hypersetup{colorlinks=false,
            pdftitle={Calibracion por borde Compton de cuatro centelladores},
            pdfsubject={Calibracion en energia}}

% Un par de atajos para no repetir unidades en el texto.
\newcommand{\mVns}{mV$\cdot$ns}
\newcommand{\DkeV}{\mVns/keV}
\newcommand{\keVdep}{keV\textsubscript{dep}}

% Caja para las advertencias, que en este informe son la mitad del contenido.
\newenvironment{aviso}%
  {\par\medskip\noindent\begin{minipage}{\linewidth}%
   \color{gris}\small\rule{\linewidth}{0.4pt}\par\medskip}%
  {\par\medskip\rule{\linewidth}{0.4pt}\end{minipage}\par\medskip}

% Los numeros del informe salen de aca, no se tipean.
% Numeros del informe. Los escribe analisis/informe.py: no editar.
% Cada uno se usa en el texto como \NombreDeLaMacro.
\newcommand{\NCorridas}{28}
\newcommand{\NDetectores}{4}
\newcommand{\NPuntosTestA}{12}
\newcommand{\NGruposTestB}{8}
\newcommand{\Rama}{Edep\_keV}
\newcommand{\PostFijo}{12}
\newcommand{\CocEsperado}{0.980}
\newcommand{\QuenchCo}{0.9541}
\newcommand{\QuenchCs}{0.9349}
\newcommand{\BMin}{7.73}
\newcommand{\BMax}{10.42}
\newcommand{\DesacMin}{2.2}
\newcommand{\DesacMax}{19.3}
\newcommand{\ErrRelMediano}{0.19}
\newcommand{\RepMin}{22}
\newcommand{\RepMax}{27}
\newcommand{\ControlMax}{1.24}
\newcommand{\ChiMin}{1.18}
\newcommand{\ChiMax}{2.70}
\newcommand{\FwhmMin}{23}
\newcommand{\FwhmMax}{37}
\newcommand{\CargaMin}{72}
\newcommand{\CargaMax}{82}
\newcommand{\NSinFondo}{26}
\newcommand{\RhoMin}{0.892}
\newcommand{\RhoMax}{0.983}
\newcommand{\MasaElectron}{510.999}
\newcommand{\EgammaCs}{661.657}
\newcommand{\EgammaCoA}{1173.228}
\newcommand{\EgammaCoB}{1332.492}
\newcommand{\BordeCs}{477.3}
\newcommand{\BordeCoA}{963.4}
\newcommand{\BordeCoB}{1118.1}
\newcommand{\NArranques}{8}
\newcommand{\NBins}{500}
\newcommand{\DispersionMax}{0.19}
\newcommand{\DispersionMediana}{menos de 0.001}
\newcommand{\CorrMin}{-0.69}
\newcommand{\CorrMax}{-0.33}
\newcommand{\NFigurasDiag}{168}
\newcommand{\CuentasBin}{272}
\newcommand{\InicioRel}{0.45}
\newcommand{\FechaInforme}{22-09-2026}
\newcommand{\SigmaRelNCC}{0.09}
\newcommand{\NBarrido}{240}
\newcommand{\CotaDMin}{0.6}
\newcommand{\CotaDMax}{1.6}
\newcommand{\CotaSigmaMin}{0.02}
\newcommand{\CotaSigmaMax}{0.30}
\newcommand{\CotaLambdaMin}{10}
\newcommand{\LambdaSemillaRel}{15}
\newcommand{\InicioAjusteRel}{0.25}
\newcommand{\NMarkov}{100}
\newcommand{\Alfa}{0.95}
\newcommand{\EjCorrida}{detector A a 1600 V con Cs-137}
\newcommand{\EjB}{649}
\newcommand{\EjRho}{0.9808}
\newcommand{\EjVentanaLo}{292}
\newcommand{\EjVentanaHi}{824}
\newcommand{\EjSemillaD}{1.3597}
\newcommand{\EjCotaDMin}{0.816}
\newcommand{\EjCotaDMax}{2.175}
\newcommand{\EjD}{1.3563}
\newcommand{\EjDesvSemilla}{-0.2}
\newcommand{\NSemillas}{28}
\newcommand{\SemillaMediana}{0.51}
\newcommand{\SemillaMax}{4.9}
% --- la seccion 'Del pulso al espectro' ---
\newcommand{\Muestras}{1024}
\newcommand{\Registro}{200}
\newcommand{\Muestreo}{5.12}
\newcommand{\DtMedio}{0.195}
\newcommand{\DtMin}{0.09}
\newcommand{\DtMax}{0.30}
\newcommand{\EventosTipico}{80\,000}
\newcommand{\RecorteBordes}{2}
\newcommand{\UmbralRuido}{3}
\newcommand{\FracOnset}{10}
\newcommand{\GuardaBase}{4}
\newcommand{\MinBase}{10}
\newcommand{\VSatura}{495}
\newcommand{\NVentana}{600}
\newcommand{\AmpMinimaForma}{20}
\newcommand{\KConv}{3}
\newcommand{\TramoConv}{4}
\newcommand{\MargenRegistro}{5}
\newcommand{\PreMin}{2.0}
\newcommand{\PreMax}{8.8}
\newcommand{\TComun}{40}
\newcommand{\EjEventos}{80\,322}
\newcommand{\EjBuenos}{79\,755}
\newcommand{\EjDescartados}{0.71}
\newcommand{\EjOnsetMediano}{69}
\newcommand{\EjOnsetCinco}{67}
\newcommand{\EjOnsetNoventaycinco}{103}
\newcommand{\EjAmpMediana}{57}
\newcommand{\EjAmpNoventaynueve}{115}
\newcommand{\EjAreaMediana}{416}
\newcommand{\EjNForma}{570}
\newcommand{\EjPico}{4.1}
\newcommand{\EjFwhmPulso}{7.1}
\newcommand{\EjCarga}{71.8}
\newcommand{\EjCargaComun}{81.7}
\newcommand{\EjPostMedido}{98}
\newcommand{\EjCargaPostMedido}{97.3}
\newcommand{\EjCargaPico}{25.2}
\newcommand{\EjRegistro}{130}
\newcommand{\EjVariacionForma}{4.9}
\newcommand{\EjAnchoBin}{2.7}
\newcommand{\EjRangoMax}{1365}
\newcommand{\EjUmbral}{152}
\newcommand{\EjTPared}{1373}
\newcommand{\EjTVivo}{1.7}
\newcommand{\EjViva}{0.12}
\newcommand{\EjScaler}{46}
\newcommand{\EjTasaGuardada}{58}
\newcommand{\EjKf}{0.00064}
\newcommand{\EjPre}{2.0}
\newcommand{\EjNFondo}{18\,794}
\newcommand{\EjTVivoFondo}{2669}
\newcommand{\EjKfPared}{0.51}
\newcommand{\EjFactorVivo}{806}
\newcommand{\RangoFactor}{1.6}
\newcommand{\RangoPercentil}{99.5}
\newcommand{\VivaCsMin}{0.08}
\newcommand{\VivaCsMax}{0.58}
\newcommand{\VivaCoMin}{28}
\newcommand{\VivaCoMax}{52}
\newcommand{\RegistroMin}{47}
\newcommand{\RegistroMax}{158}
\newcommand{\CargaComunMin}{81}
\newcommand{\CargaComunMax}{83}
\newcommand{\VariacionFormaMax}{8.3}
\newcommand{\VariacionFormaMediana}{6.7}
\newcommand{\VariacionAltaMax}{4.3}
\newcommand{\VariacionAltaMediana}{2.5}
\newcommand{\EjVariacionAlta}{3.0}
\newcommand{\DescartadosMax}{6.2}
\newcommand{\DescartadosMediana}{1.2}
\newcommand{\SaturadosMax}{0.88}
\newcommand{\TruncadosMax}{5.8}
\newcommand{\SinPreMax}{6.4}
\newcommand{\DmilseisA}{1.3563}
\newcommand{\DmilseisB}{1.0717}
\newcommand{\DmilseisC}{1.2055}
\newcommand{\DmilseisD}{0.4934}
\newcommand{\RazonGanancias}{2.7}
\newcommand{\EjemploKeV}{737}
\newcommand{\SDNuestroMin}{1}
\newcommand{\SDNuestroMax}{53}
\newcommand{\CosenoNuestro}{1.000000}
\newcommand{\RdosNuestroMin}{0.936}
\newcommand{\SDPaperMin}{1.6}
\newcommand{\SDPaperMax}{7.3}
\newcommand{\RdosPaperMin}{0.999055}
\newcommand{\RdosPaperMax}{0.999974}
\newcommand{\CosenoPaperMin}{0.999993}
\newcommand{\CosenoPaperMax}{0.999999}
\newcommand{\RobustezPaper}{1.2}
\newcommand{\RobustezPaperSinAngosta}{0.3}
\newcommand{\ReduccionDif}{80}
\newcommand{\ReduccionGauss}{85}
\newcommand{\RobustezNuestraMin}{0.5}
\newcommand{\RobustezNuestraMax}{1.6}
\newcommand{\SubidaCincuenta}{34}
\newcommand{\EjErrCurv}{0.0015}
\newcommand{\EjErrPerfil}{0.0016}
\newcommand{\EjErrEscalado}{0.0022}
\newcommand{\EjEscala}{1.43}
\newcommand{\EjChiNdf}{2.05}
\newcommand{\CierreN}{30}
\newcommand{\CierreCuentasBin}{292}
\newcommand{\CierreD}{1.3563}
\newcommand{\CierreSesgoNeyman}{-0.04}
\newcommand{\CierreErrSesgoNeyman}{0.03}
\newcommand{\CierreDispNeyman}{0.17}
\newcommand{\CierreErrRepNeyman}{0.15}
\newcommand{\CierreCobNeyman}{18}
\newcommand{\CierreSesgoSigmaNeyman}{-0.4}
\newcommand{\CierreSemillaNeyman}{0.57}
\newcommand{\CierreSesgoPearson}{+0.00}
\newcommand{\CierreErrSesgoPearson}{0.03}
\newcommand{\CierreDispPearson}{0.17}
\newcommand{\CierreErrRepPearson}{0.15}
\newcommand{\CierreCobPearson}{21}
\newcommand{\CierreSesgoSigmaPearson}{+0.4}
\newcommand{\CierreSemillaPearson}{0.57}
\newcommand{\EjNptsAuto}{163}
\newcommand{\EjRhoRecalc}{0.9808}
\newcommand{\NptsMin}{12}
\newcommand{\NptsMax}{400}
\newcommand{\NptsRhoPocos}{0.9828}
\newcommand{\NptsRhoAuto}{0.9808}
\newcommand{\NptsBDesvMax}{0.3}
\newcommand{\ClasSigmaMin}{0.03}
\newcommand{\ClasSigmaMax}{0.18}
\newcommand{\ClasEgamma}{661.7}
\newcommand{\ClasEmax}{477.3}
\newcommand{\ClasMitadLo}{+0.7}
\newcommand{\ClasMitadHi}{+7.5}
\newcommand{\ClasMitadSinLo}{+0.7}
\newcommand{\ClasMitadSinHi}{+7.5}
\newcommand{\ClasDifLo}{+0.7}
\newcommand{\ClasDifHi}{+7.1}
\newcommand{\ClasDifSinLo}{+0.3}
\newcommand{\ClasDifSinHi}{+5.4}
\newcommand{\ClasGaussLo}{+2.3}
\newcommand{\ClasGaussHi}{+7.1}
\newcommand{\ClasGaussSinLo}{+2.3}
\newcommand{\ClasGaussSinHi}{+7.1}
\newcommand{\ClasNReales}{28}
\newcommand{\ClasRealMitadMed}{+11.9}
\newcommand{\ClasRealMitadMin}{+9.4}
\newcommand{\ClasRealMitadMax}{+21.4}
\newcommand{\ClasRealDifMed}{+10.3}
\newcommand{\ClasRealDifMin}{+4.3}
\newcommand{\ClasRealDifMax}{+30.0}
\newcommand{\ClasRealGaussMed}{+12.8}
\newcommand{\ClasRealGaussMin}{+11.8}
\newcommand{\ClasRealGaussMax}{+35.5}
\newcommand{\FotoSobreComptonCs}{4\times10^{-5}}
\newcommand{\FotoSobreComptonNaICs}{0.13}
\newcommand{\LibreCaminoCs}{11.6}
\newcommand{\ProbInteraccionCs}{35}
\newcommand{\FotoSobreComptonCoB}{1\times10^{-5}}
\newcommand{\FotoSobreComptonNaICoB}{0.04}
\newcommand{\LibreCaminoCoB}{16.2}
\newcommand{\ProbInteraccionCoB}{27}
\newcommand{\CruceFotoComptonNaI}{266}
\newcommand{\CruceFotoComptonPlastico}{21}
\newcommand{\DensidadBC}{1.032}
\newcommand{\LadoDetector}{5.08}
\newcommand{\SigmaThomson}{0.665}
\newcommand{\SigmaKNCs}{0.256}
\newcommand{\ElectronesPorGramo}{3.26\times10^{23}}
\newcommand{\XcomFraccionH}{0.0853}
\newcommand{\XcomFraccionC}{0.9147}
\newcommand{\BirksSim}{0.126}
\newcommand{\KNvsXcom}{0.9}
\newcommand{\KNvsXcomBajo}{35}
\newcommand{\EcUnoImpresaCs}{184.3}
\newcommand{\EcUnoImpresaCoA}{209.8}
\newcommand{\SimRadio}{2.54}
\newcommand{\SimSemialtura}{2.54}
\newcommand{\SimPared}{10}
\newcommand{\SimTapa}{6}
\newcommand{\SimBirks}{0.126}
\newcommand{\SimDistancia}{10}
\newcommand{\SimPmtSemilargo}{1.5}
\newcommand{\SimCorte}{0.1}
\newcommand{\SimDensidad}{1.032}
\newcommand{\SimHaz}{\texttt{pencil}}
\newcommand{\SimMaterial}{\texttt{G4\_PLASTIC\_SC\_VINYLTOLUENE}}
\newcommand{\SimFisica}{\texttt{G4EmStandardPhysics\_option4}}
\newcommand{\SimVersion}{\texttt{11.4}}
\newcommand{\SimVentanaPmt}{\texttt{G4\_Pyrex\_Glass}}
\newcommand{\SimLargo}{5.08}
\newcommand{\SimBinsPlantilla}{350}
\newcommand{\SimPrimariosCo}{3}
\newcommand{\SimFilasCo}{852079}
\newcommand{\SimFracDepositaCo}{28.4}
\newcommand{\SimFracSimpleCo}{75.1}
\newcommand{\SimFracMultipleCo}{22.8}
\newcommand{\SimFracFotoCo}{2.9}
\newcommand{\SimNcMaxCo}{11}
\newcommand{\SimQfBajoCo}{0.863}
\newcommand{\SimQfAjusteCo}{0.954}
\newcommand{\SimEdepMaxCo}{1333}
\newcommand{\SimBordeLuzCo}{921}
\newcommand{\SimQfAltoCo}{0.954}
\newcommand{\SimRazonLineasCo}{1.049}
\newcommand{\SimPrimariosCs}{10}
\newcommand{\SimFilasCs}{3531670}
\newcommand{\SimFracDepositaCs}{35.3}
\newcommand{\SimFracSimpleCs}{73.0}
\newcommand{\SimFracMultipleCs}{26.3}
\newcommand{\SimFracFotoCs}{1.3}
\newcommand{\SimNcMaxCs}{13}
\newcommand{\SimQfBajoCs}{0.859}
\newcommand{\SimQfAjusteCs}{0.935}
\newcommand{\SimEdepMaxCs}{662}
\newcommand{\SimBordeLuzCs}{448}
\newcommand{\SimBordeLuzCoB}{1069}
\newcommand{\SimExcesoMax}{1.7}
\newcommand{\BinsMin}{100}
\newcommand{\BinsMax}{700}
\newcommand{\BinsAnchoMin}{0.9}
\newcommand{\BinsAnchoMax}{6.6}
\newcommand{\BinsDesvMax}{0.06}
\newcommand{\BinsSigmaE}{51}
\newcommand{\BinsSigmaSobreBin}{54}
\newcommand{\BinsErrD}{0.0025}
\newcommand{\BinsNa}{100}
\newcommand{\BinsDa}{1.3565}
\newcommand{\BinsNb}{175}
\newcommand{\BinsDb}{1.3565}
\newcommand{\BinsNc}{350}
\newcommand{\BinsDc}{1.3563}
\newcommand{\BinsNd}{700}
\newcommand{\BinsDd}{1.3555}



\title{\color{azul}\bfseries Calibración en energía de cuatro centelladores
       plásticos por borde Compton\\[2mm]
       \normalsize\mdseries\color{black}
       Campaña de \NCorridas{} corridas: Co-60 y Cs-137 a 1550--1700 V}
\author{Mateo\thanks{Física Nuclear, Pontificia Universidad Católica de Chile.
        Completar autoría y filiación antes de circular el documento.}}
\date{\FechaInforme}

\begin{document}
\maketitle

\begin{abstract}
\noindent
Se calibraron \NCorridas{} corridas de los \NDetectores{} detectores (A, B, C
y D) tomadas entre junio y septiembre de 2026, convirtiendo el área del pulso
en \mVns{} a energía \emph{depositada} en keV mediante ajuste del borde Compton
contra la simulación de Geant4 de cada fuente. Todas se integran en la misma
compuerta de \PostFijo{} ns tras el onset y se ajustan con el mismo modelo, la
plantilla convolucionada sin más componentes. La ganancia sigue en los cuatro
detectores una ley de potencia con el alto voltaje, $D \propto V^{b}$, con
$b$ entre \BMin{} y \BMax{}, compatible con lo que se espera de un
fotomultiplicador y muy preferida sobre una recta. En cambio la
\emph{linealidad del eje} no cierra: en los \NPuntosTestA{} puntos de operación
donde hay dos fuentes, las calibraciones de Cs-137 y Co-60 difieren entre
\DesacMin\,\% y \DesacMax\,\% una vez descontado el quenching de Birks, con
signo que cambia de un detector a otro, mientras que el error estadístico de
cada calibración es del orden del \ErrRelMediano\,\%. El informe documenta el
método paso a paso, las decisiones de análisis comunes a todas las corridas,
cómo se lee cada tabla y cada figura, los controles que se hicieron para
descartar que el desacuerdo sea un artefacto, y qué mediciones faltan para
cerrarlo.
\end{abstract}

\tableofcontents
\bigskip

% ===========================================================================
\section{Cómo leer este informe}
\label{sec:comoleer}

El informe entrega un número por corrida, la ganancia $D$, y dos análisis
construidos sobre esos números. Para no tener que leerlo entero antes de
encontrar lo que se busca, esto es el mapa.

\begin{center}
\begin{tabular}{lll}
\toprule
\textbf{Si la pregunta es\ldots} & \textbf{está en} & \textbf{lo entrega} \\
\midrule
Qué se midió y cuándo & \S\ref{sec:medido} & texto \\
Qué es un pulso y cómo se integra & \S\ref{sec:pulsos} &
   texto, figuras~\ref{fig:integracion} a~\ref{fig:area-espectro},
   tabla~\ref{tab:pulsos} \\
Por qué un plástico no tiene fotopico & \S\ref{sec:interaccion} &
   texto, figura~\ref{fig:teoria-interaccion} \\
De dónde sale la energía del borde & \S\ref{sec:cinematica} &
   texto, tabla~\ref{tab:cinematica}, figura~\ref{fig:teoria-cinematica} \\
Cómo funciona el NCC & \S\ref{sec:ncc} &
   texto, figura~\ref{fig:teoria-ncc} \\
Cómo funciona el recocido simulado & \S\ref{sec:sa} &
   texto, figura~\ref{fig:teoria-sa} \\
Cuánto se equivocan los métodos clásicos & \S\ref{sec:clasicos} &
   tabla~\ref{tab:clasicos}, figura~\ref{fig:teoria-clasicos} \\
De dónde sale la barra de error & \S\ref{sec:estadistica} &
   tabla~\ref{tab:cierre}, figura~\ref{fig:teoria-estadistica} \\
Qué hay en la simulación y qué le falta & \S\ref{sec:geant4} &
   tabla~\ref{tab:simulacion}, figura~\ref{fig:simulacion} \\
Cómo se obtiene $D$, paso a paso & \S\ref{sec:cadena} & texto \\
Cuánto dio cada corrida & \S\ref{sec:calibraciones} & tabla~\ref{tab:calib} \\
Si la calibración de un detector es lineal & \S\ref{sec:testA} &
   tabla~\ref{tab:testA}, figuras~\ref{fig:cociente} y~\ref{fig:porpunto} \\
Cómo se compara un detector con otro & \S\ref{sec:detectores} &
   tablas~\ref{tab:testB} y~\ref{tab:detectores} \\
Qué pasó con el detector D & \S\ref{sec:repro} &
   tabla~\ref{tab:repro} \\
Cómo se compara esto con el paper & \S\ref{sec:paper} &
   tablas~\ref{tab:robustez} y~\ref{tab:paper} \\
Cómo convertir un pulso medido a energía & \S\ref{sec:uso} & receta \\
Qué NO se puede concluir todavía & \S\ref{sec:limites} & lista \\
Qué significa cada símbolo & apéndice~\ref{ap:glosario} & glosario \\
Qué archivo es cada cosa & apéndice~\ref{ap:archivos} & tabla \\
\bottomrule
\end{tabular}
\end{center}

\bigskip
\noindent
Tres convenciones que valen para todo el documento:

\begin{itemize}
\item \textbf{Las unidades.} El área del pulso está en \mVns{} (milivolt por
      nanosegundo: es la integral del pulso, no su altura). La energía está en
      keV \emph{depositados} en el centellador. Entonces $D$ va en \DkeV{} y su
      inverso $1/D$ en keV por \mVns, que es el que se usa para convertir.
\item \textbf{Ningún número de este texto está tipeado a mano.} Todos salen de
      los datos al generar el documento (apéndice~\ref{ap:reproducir}). Si se
      recalculan las calibraciones, el informe se rehace con los valores
      nuevos y no hay forma de que el texto quede desactualizado respecto de
      las tablas.
\item \textbf{Los avisos en gris} marcan las cosas que hay que tener presentes
      para no leer de más un resultado. No son notas al pie decorativas: cada
      uno está porque el número de al lado se puede malinterpretar.
\end{itemize}

% ===========================================================================
\section{Qué se midió}
\label{sec:medido}

\subsection{La campaña}

Cuatro detectores (A, B, C y D), cada uno un centellador plástico leído por un
fotomultiplicador, digitalizados con una placa DRS4. Tres voltajes nominales
(1600, 1650 y 1700 V) y dos fuentes gamma, más una corrida de junio del
detector A a 1550 V. En total \NCorridas{} corridas repartidas en cuatro
sesiones:

\begin{itemize}
\item \textbf{23 de junio}: Co-60 del detector A a 1550 V, en XML, con su
      fondo. Es la única corrida de la campaña con fondo propio además del
      Cs-137 del detector A a 1600 V.
\item \textbf{26 de agosto}: Cs-137 del detector A a los tres voltajes, en XML.
\item \textbf{1 de septiembre}: Cs-137 de los cuatro detectores y Co-60 del A,
      en formato binario del DRS4. Las tres corridas del detector D de esa
      sesión están \textbf{descartadas}: el detector tenía la tierra mal
      conectada (sección~\ref{sec:repro}).
\item \textbf{2 y 8 de septiembre}: Co-60 del detector C (2-sep) y Co-60 más
      Cs-137 del detector D (8-sep). Éstas son las corridas válidas del
      detector D, ya con la tierra arreglada, y las únicas de toda la campaña
      donde las dos fuentes se midieron el mismo día.
\item \textbf{22 de septiembre}: Co-60 del detector B a los tres voltajes, en
      binario. Era la medición que faltaba para que los cuatro detectores
      tuvieran las dos fuentes. Se verificó que los archivos no fueran copias
      de los de otro detector antes de incorporarlos, por lo que había pasado
      el 10 de septiembre (sección~\ref{sec:limites}).
\end{itemize}

La fecha de cada corrida no se tipea: se lee de la fecha del propio archivo y
aparece en todas las tablas. No es un dato decorativo. Dos fuentes medidas en
sesiones distintas pueden no estar en el mismo punto de operación real aunque
el voltaje nominal coincida, y eso es exactamente lo que hay que poder
descartar en la sección~\ref{sec:testA}.

\subsection{Qué es $D$ y por qué es lo único que se busca}

El observable es el área del pulso: la integral de la señal del
fotomultiplicador, en \mVns. Es proporcional a la carga que sale del tubo y por
lo tanto a la luz que produjo el centellador. Calibrar es encontrar el factor
que lleva esa área a energía:
\[
   \text{área [\mVns]} \;=\; D \cdot \text{energía [\keVdep]},
   \qquad\text{o al revés,}\qquad
   E \;=\; \frac{\text{área}}{D}.
\]

La recta pasa por el \textbf{origen}, sin ordenada: área cero corresponde a
energía cero. No es una simplificación cómoda, es parte del método (la relación
canal--energía del modelo es $E_{\text{eff}} = E\cdot B/E_{\max}$), y por eso
hay un solo parámetro que ajustar en vez de dos. Que esa hipótesis se sostenga
o no es justamente lo que testea la sección~\ref{sec:testA}.

\begin{aviso}
\textbf{Escala.} Todo este informe está en energía \emph{depositada}
(plantilla \texttt{\Rama} de la simulación), no en luz. La diferencia no es
cosmética: el área mide luz, y la luz por keV depositado depende de la energía
porque el quenching de Birks todavía no saturó. Las $D$ de este informe se
llevan ese factor adentro y \textbf{no} son comparables con las $D$ en keVee
(keV electrón-equivalentes) del resto del proyecto.
\end{aviso}

% ===========================================================================
\section{Del pulso al espectro}
\label{sec:pulsos}

Todo lo que el informe llama «área» sale de esta sección. Antes de que haya
un borde Compton que ajustar hay un pulso eléctrico digitalizado, y la
decisión de qué número se le asigna a ese pulso es la que fija el observable.
Acá está qué registra el digitalizador, cómo se convierte cada registro en un
área, por qué la compuerta es la que es, y cómo esas áreas se vuelven el
espectro que después se ajusta. Los pulsos que se muestran son reales, de la
corrida \EjCorrida{} (la misma que ilustra el NCC en la
sección~\ref{sec:ncc}); los números por corrida están en la
tabla~\ref{tab:pulsos}.

\subsection{Qué registra el digitalizador}

Cada evento es un registro de \Muestras{} muestras de la señal del
fotomultiplicador, \Registro{} ns a \Muestreo{} GS/s, en mV. El pulso es
negativo (es la corriente del ánodo sobre la carga del digitalizador) y la
placa dispara sobre la propia señal: cuando cruza el umbral del discriminador
se guarda el registro con el pulso en una posición fija, que en la corrida de
ejemplo es \EjOnsetMediano{} ns desde el inicio (panel (a) de la
figura~\ref{fig:integracion}). Fija en promedio, no evento a evento: el
discriminador tiene \emph{walk}, los pulsos chicos cruzan el umbral más tarde
que los grandes, y entre los percentiles 5 y 95 el onset se mueve de
\EjOnsetCinco{} a \EjOnsetNoventaycinco{} ns. Una corrida típica tiene
\EventosTipico{} registros.

Tres cosas del registro que el análisis tiene que respetar y que no son
obvias mirando un pulso:

\begin{itemize}
\item \textbf{El eje de tiempo no es uniforme.} El DRS4 muestrea con un anillo
      de celdas de capacitores y cada celda tiene su propio ancho: el paso
      entre muestras vale \DtMedio{} ns en promedio pero va de \DtMin{} a
      \DtMax{} ns (percentiles 1 y 99) de una celda a la siguiente, y cambia
      evento a evento con la celda en que cayó el trigger. El archivo trae el
      tiempo de cada muestra, y toda integral se hace sobre ese eje, nunca
      multiplicando por un paso constante. Las corridas de la sesión del 26 de
      agosto se guardaron sin esa calibración, con un eje uniforme; el análisis
      no distingue los dos casos porque integra siempre sobre el eje que viene.
\item \textbf{Las celdas de los extremos no son señal.} La primera y la última
      del anillo dan espigas de decenas de mV; se descartan \RecorteBordes{}
      muestras de cada lado antes de mirar nada. Y el tiempo de la primera
      muestra no es cero ni es el mismo entre sesiones, así que todo se mide
      desde el inicio del registro y no en el eje absoluto del archivo.
\item \textbf{Dos formatos, un solo lector.} Las sesiones de junio y agosto
      están en el XML del DRSOSC; las de septiembre, en el binario del DRS4,
      que no tiene extensión. El formato se decide por los primeros bytes del
      archivo y el canal que contiene la señal se lee del primer evento, porque
      cambió entre sesiones y leer el canal equivocado devuelve cero pulsos sin
      ningún error.
\end{itemize}

\begin{figure}[htbp]
\centering
\includegraphics[width=\linewidth]{integracion_pulso.pdf}
\caption{Del registro al área, sobre la corrida \EjCorrida. (a) Treinta
registros tal como vienen del digitalizador: el pulso cae siempre cerca del
mismo lugar, con el \emph{walk} del trigger marcado (mediana y percentiles 5 a
95 del onset). (b) Un pulso de amplitud mediana con cada paso de la
integración: la zona verde es de donde sale la línea base (mediana de las
muestras hasta \GuardaBase{} ns antes del onset), la línea punteada es el
\FracOnset\,\% de la amplitud cuyo cruce define el onset, la zona naranja es la
compuerta $[-\EjPre, +\PostFijo]$~ns respecto del onset y el área roja es lo
que se integra. (c) El pulso promedio de \EjNForma{} eventos alineados por su
onset y normalizados a amplitud 1, con la compuerta común y el punto donde el
pulso se hunde en el ruido. (d) La carga acumulada del pulso promedio, en
porcentaje de la que llega a acumularse en el registro: la compuerta recoge
el \EjCarga\,\%.}
\label{fig:integracion}
\end{figure}

\subsection{De un registro a un número}
\label{sec:area}

El panel (b) de la figura~\ref{fig:integracion} es un pulso real con cada paso
marcado. En el orden en que se ejecutan:

\begin{enumerate}
\item \textbf{Dos filtros antes de integrar.} Un registro idéntico al anterior
      se descarta (el DRSOSC a veces vuelca dos veces el mismo evento, y
      contarlo dos veces mete un pico falso en el espectro), y un registro cuya
      excursión pico a pico no llega a \UmbralRuido{} mV no tiene pulso. Los dos
      cortes se cuentan, y en la corrida de ejemplo no descartan ninguno.
\item \textbf{Localizar el pulso con una referencia provisoria.} La mediana del
      registro entero sirve de línea base preliminar, porque el pulso ocupa una
      décima parte de las muestras y la mediana no lo ve. Respecto de ella se
      busca la excursión máxima, que da el pico y la polaridad (deducida, no
      supuesta), y desde el pico se retrocede hasta el último cruce del
      \FracOnset\,\% de la amplitud: ése es el \textbf{onset}. Se retrocede
      desde el pico y no se avanza desde el inicio del registro porque una
      fluctuación de la línea base que pase el umbral dispararía un onset falso
      decenas de ns antes del pulso.
\item \textbf{La línea base, relativa al onset.} Es la mediana de las muestras
      anteriores a \GuardaBase{} ns antes del onset (zona verde). Mediana y no
      media porque cuesta lo mismo y es inmune a que se cuele el pie de un
      pulso anterior. Del mismo tramo sale el ruido, como desviación robusta
      ($1.4826$ veces la mediana de las desviaciones absolutas), que es la
      vara de todos los umbrales que siguen: ningún criterio depende de cuántos
      mV de ruido tiene este detector en particular. Si quedan menos de
      \MinBase{} muestras antes del onset, la línea base se toma de la cola
      posterior a la compuerta y el evento queda marcado, pero \emph{no} se
      descarta: rechazar por posición del onset sería un corte que depende de la
      amplitud, porque el \emph{walk} pone el onset más tarde en los pulsos
      chicos, y un corte dependiente de la amplitud deforma justo lo que se
      quiere calibrar.
\item \textbf{Volver a medir con la línea base buena.} Pico, amplitud y onset
      se recalculan respecto de la línea base definitiva. Una amplitud por
      debajo de cinco veces el ruido se marca como «sin pulso».
\item \textbf{La compuerta}, $[\text{onset} - \text{pre},\ \text{onset} +
      \text{post}]$, anclada al onset \emph{de ese evento}. El \emph{pre} se
      mide sobre cada archivo (sección~\ref{sec:compuerta}; vale entre
      \PreMin{} y \PreMax{} ns en la campaña) y es zona de línea base, así que
      aporta ruido pero no señal; el \emph{post} es \PostFijo{} ns para todas
      las corridas. Anclar al onset y no a un tiempo fijo del registro es lo
      que hace a la medición inmune al \emph{walk} y al jitter del trigger: la
      ventana viaja con el pulso.
\item \textbf{Integrar.} El área es la integral trapezoidal de la señal menos
      la línea base dentro de la compuerta, sobre el eje de tiempo real del
      registro, con el signo dado vuelta para que sea positiva:
      \[
         \text{área} \;=\; \int_{\text{onset}-\text{pre}}^{\text{onset}+\text{post}}
         \big[\,b - v(t)\,\big]\,dt \qquad [\text{mV}\cdot\text{ns}],
      \]
      con $b$ la línea base. Es la carga que entregó el fotomultiplicador en
      esa ventana, a menos de la impedancia de entrada.
\item \textbf{Marcar, no descartar.} Cada evento sale con banderas:
      \emph{saturado} (alguna muestra alcanzó los \VSatura{} mV, el tope del
      digitalizador), \emph{truncado} (la compuerta se sale del registro por
      cualquiera de los dos extremos), \emph{sin pre-pulso} (la línea base
      salió de la cola) y \emph{sin pulso}. Qué banderas descartan se decide
      al histogramar, no acá, y es la misma decisión para muestra y fondo: si
      no, los dos histogramas tendrían eficiencias distintas y la resta no
      significaría nada. Se descartan los saturados, los truncados y los sin
      pulso; los «sin pre-pulso» se conservan, por la razón del punto 3.
\end{enumerate}

En la corrida de ejemplo, de \EjEventos{} registros quedan \EjBuenos{}
(\EjDescartados\,\% descartados); en la campaña la fracción descartada tiene
mediana \DescartadosMediana\,\% y máximo \DescartadosMax\,\% (el Co-60 de
junio, donde una parte de los pulsos cae al principio del registro, antes de
la posición del trigger, y la compuerta se sale por ese extremo). La saturación no pasa de
\SaturadosMax\,\% en ninguna corrida.

\begin{aviso}
La línea base relativa al onset no es un refinamiento: es la corrección más
grande que se le hizo a esta cadena. Una versión anterior la promediaba sobre
el primer 20\,\% del registro, y en la corrida de junio el pulso arranca
apenas después de ese tramo: la ventana de línea base se comía la subida del
propio pulso, el sesgo se multiplicaba al integrar, y como el jitter del
trigger movía el pulso evento a evento, cada uno perdía una fracción distinta
de su carga. Eso ensanchaba el espectro hasta duplicar la resolución aparente,
y cambiaba la ganancia medida en un factor mayor que cualquier otro efecto de
este informe. La regla que queda es que nada del cálculo del área usa una
ventana absoluta del registro.
\end{aviso}

\subsection{La compuerta: cómo se mide y por qué es corta}
\label{sec:compuerta}

Para saber cuánto dura el pulso hay que mirarlo con menos ruido del que tiene
un evento solo. Se toman los primeros \NVentana{} registros del archivo, se
quedan los de amplitud mayor que \AmpMinimaForma{} mV (los chicos tienen el
onset ruidoso; esto es sólo para medir la forma, el área después se calcula
sobre todos), se alinea cada uno por su onset interpolando sobre una grilla
común, se lo normaliza por su amplitud y se promedia. El resultado es la forma
del pulso, panel (c) de la figura~\ref{fig:integracion}: llega al pico
\EjPico{} ns después del onset, tiene \EjFwhmPulso{} ns de ancho a media
altura y una cola larga con un pequeño rebote. Sobre esa forma promedio se aplica
el criterio de convergencia: el pulso «empezó» en el último instante antes del
pico en que la señal estuvo quieta (por debajo de \KConv{} veces el ruido del
promedio) durante \TramoConv{} ns seguidos, y «terminó» en el primer instante
después del pico en que vuelve a estarlo. Lo primero da el \emph{pre}; lo
segundo, la duración que el pulso pide: en la corrida de ejemplo,
\EjPostMedido{} ns, donde ya acumuló el \EjCargaPostMedido\,\% de su carga
(panel (d)).

Ese \emph{post} medido no es el que se usa, por dos razones que conviene
separar:

\begin{itemize}
\item \textbf{El registro no siempre lo contiene.} Después del onset
      sobran entre \RegistroMin{} y \RegistroMax{} ns según la corrida
      (columna \emph{Registro} de la tabla~\ref{tab:pulsos}), y las más
      cortas son las de septiembre, donde el trigger quedó tarde en el
      registro. Cuando el pulso
      pide más de lo que hay, la medición recorta el \emph{post} al espacio
      disponible en los eventos de amplitud sobre la mediana (los chicos tienen
      el onset más tarde por el \emph{walk} y sesgarían el percentil) menos
      \MargenRegistro{} ns de margen. Un evento cuya compuerta se sale del
      registro pierde una fracción de carga que depende de cuándo llegó, que
      es exactamente el error que la compuerta anclada al onset existe para
      evitar.
\item \textbf{Cada archivo pediría una compuerta distinta}, y dos corridas
      integradas con compuertas distintas recogen fracciones distintas de la
      carga, lo que entra derecho en $D$. Para comparar corridas la compuerta
      tiene que ser la misma.
\end{itemize}

De ahí la decisión de la sección~\ref{sec:decisiones}: \emph{post} $=
\PostFijo$ ns para las \NCorridas{} corridas. Corta, además, porque pasado el
grueso del pulso seguir integrando agrega línea base y su ruido más que señal.
El precio está en el panel (d): la compuerta recoge el \EjCarga\,\% de la carga
del pulso registrado, y en la campaña entre \CargaMin{} y \CargaMax\,\%
(columna \emph{Carga} de la tabla~\ref{tab:calib}).

\paragraph{Ese rango no es una diferencia entre detectores.} La
figura~\ref{fig:pulsos-detectores} superpone el pulso promedio de los cuatro
detectores a 1600~V con Cs-137: la subida, el pico y la caída coinciden. Lo que
varía de una corrida a otra es cuánto registro hay después del onset, y como
la fracción «de la carga registrada» se normaliza a lo que el registro
contiene, una corrida con \RegistroMin{} ns de registro parece recoger más que
una con \RegistroMax. Normalizando en cambio a la carga acumulada hasta
\TComun{} ns del onset, que todos los registros contienen, la compuerta recoge
entre \CargaComunMin{} y \CargaComunMax\,\% en las \NCorridas{} corridas
(última columna de la tabla~\ref{tab:pulsos}). O sea que la compuerta común
mide en todas la misma porción del mismo pulso, que es lo que hace comparables
las $D$ entre detectores, voltajes y sesiones aunque el trigger haya quedado en
lugares distintos.

\begin{figure}[htbp]
\centering
\includegraphics[width=\linewidth]{pulsos_detectores.pdf}
\caption{Los cuatro detectores a 1600~V con Cs-137. (a) El pulso promedio de
cada uno, alineado por onset y normalizado a amplitud 1: la forma es la misma.
(b) La carga acumulada, relativa a la acumulada hasta \TComun{} ns del onset
(línea punteada vertical), que es el tramo que todos los registros contienen;
cada curva llega hasta donde llega su registro, marcado con una barra. Los
puntos son lo que recoge la compuerta común de \PostFijo{} ns.}
\label{fig:pulsos-detectores}
\end{figure}

\subsection{Control: la forma del pulso no depende de su tamaño}
\label{sec:forma-amplitud}

Una compuerta que recoge una fracción fija de la carga sólo sirve si la forma
del pulso es la misma para todos los tamaños. Si un pulso grande fuera
proporcionalmente más largo que uno chico, la compuerta le cortaría una
fracción distinta y el eje del espectro dejaría de ser lineal, y ninguna
calibración posterior lo arreglaría. El test es barato y no necesita saber la
energía: el cociente área/amplitud de cada evento en función de la amplitud
tiene que ser plano.

El panel (a) de la figura~\ref{fig:area-espectro} lo muestra para la corrida
de ejemplo. De punta a punta del rango de amplitudes la mediana por tramo se
mueve \EjVariacionForma\,\%, y casi todo eso está en los pulsos más chicos:
por encima de la amplitud mediana la variación es \EjVariacionAlta\,\%. En la
campaña, la variación de punta a punta tiene mediana \VariacionFormaMediana\,\%
y máximo \VariacionFormaMax\,\%; restringida a la mitad de arriba, mediana
\VariacionAltaMediana\,\% y máximo \VariacionAltaMax\,\%. La caída hacia los
pulsos chicos no es necesariamente de la forma: la amplitud es una sola
muestra, con el ruido de una sola muestra, y el onset al \FracOnset\,\% se
mide sobre unos pocos mV de señal; las dos cosas empeoran justo ahí. Lo que
el test acota es cuánto puede deformar el eje una compuerta corta, y el
control de la tabla~\ref{tab:control} (75~ns contra \PostFijo) dice por otro
camino que el desacuerdo entre fuentes de la sección~\ref{sec:testA} no viene
de acá.

\begin{figure}[htbp]
\centering
\includegraphics[width=\linewidth]{area_a_espectro.pdf}
\caption{De las áreas al espectro, corrida \EjCorrida. (a) Área contra
amplitud, evento por evento, con la mediana de área/amplitud por tramo de
amplitud. (b) Los tres histogramas con el binning común: la muestra, el fondo
ya normalizado por tiempo vivo (con $k_f = \EjKf$ queda muy por debajo de la
muestra: es despreciable en esta corrida, y el panel lo muestra en vez de
afirmarlo) y el neto, con el fin de la subida del trigger y la ventana de
ajuste que después elige el NCC.}
\label{fig:area-espectro}
\end{figure}

\subsection{De las áreas al espectro neto}
\label{sec:areas-espectro}

Con un área por evento, el espectro es un histograma. Lo que hay que decidir
es el eje, la normalización y el error, y las tres decisiones se ven en el
panel (b) de la figura~\ref{fig:area-espectro}.

\paragraph{El eje.} \NBins{} bines entre cero y \RangoFactor{} veces el
percentil \RangoPercentil{} de las áreas de la muestra (\EjRangoMax{} mV$\cdot$ns
en la corrida de ejemplo, o sea bines de \EjAnchoBin{} mV$\cdot$ns). El tope
se fija así y no con el área máxima porque unos pocos eventos de área enorme
(apilamientos, cósmicos) estirarían el rango y dejarían toda la física apretada
en los primeros bines. Muestra y fondo se histograman con \emph{el mismo} eje:
si cada uno tuviera el suyo, restarlos sería restar intervalos distintos de
mV$\cdot$ns.

\paragraph{El tiempo vivo.} Para restar el fondo hay que saber cuánto tiempo
estuvo midiendo cada corrida, y el reloj de pared no lo dice: el programa de
adquisición no alcanza a guardar todos los eventos que la placa dispara. La
placa lleva un contador (\emph{scaler}) con la tasa real de disparo, que se
guarda con cada evento; integrarlo da cuántos triggers hubo de verdad, y el
cociente entre los eventos guardados y ésos es la fracción del tiempo en que
la adquisición estuvo viva. En la corrida de ejemplo el scaler marca
\EjScaler{} kHz mientras se guardan \EjTasaGuardada{} eventos por segundo:
\EjTVivo{} s vivos de \EjTPared{} de pared, el \EjViva\,\%. Las corridas de
Cs-137 están vivas entre el \VivaCsMin{} y el \VivaCsMax\,\% del tiempo; las de
Co-60, con una tasa de disparo mucho menor, entre el \VivaCoMin{} y el
\VivaCoMax\,\%. Las marcas de tiempo de cada evento no sirven para esto: el
DRSOSC las escribe en ráfagas y miles de eventos comparten el mismo
milisegundo.

\paragraph{La resta.} Los dos histogramas se llevan al menor de los dos
tiempos vivos, para que ninguno se amplifique: $k_m = t_{\min}/t_{\text{vivo},m}$ y
$k_f = t_{\min}/t_{\text{vivo},f}$, y el neto es $k_m N_m - k_f N_f$. En la
corrida de ejemplo el fondo tiene \EjNFondo{} pulsos en \EjTVivoFondo{} s
vivos contra \EjTVivo{} de la muestra, así que entra con $k_f = \EjKf$: es la
curva gris del panel (b), órdenes de magnitud por debajo. Normalizando por
tiempo de pared el mismo fondo entraría con $k_f = \EjKfPared$, \EjFactorVivo{}
veces más, y la resta dejaría el espectro negativo. Sólo dos corridas tienen fondo propio; en las demás el neto
es el crudo, y la sección~\ref{sec:decisiones} dice por qué no se les presta
el fondo de otra sesión.

\paragraph{El error.} Cada histograma crudo es Poisson, así que la varianza del
neto es $k_m^2 N_m + k_f^2 N_f$, con piso 1 para que un bin vacío no tenga
error cero. Sin esta barra no hay $\chi^2$, ni bondad de ajuste, ni error
sobre $D$: es la diferencia entre un número y un número con incertidumbre.
Los bines que quedan negativos después de restar se conservan: son
fluctuaciones legítimas y recortarlos a cero sesgaría el ajuste hacia arriba.

\paragraph{Lo que no es física.} Por debajo del umbral del discriminador el
espectro es la curva de eficiencia del trigger, que sube de cero a uno, y no
hay plantilla que la contenga. El fin de esa subida se detecta en el propio
espectro como el primer máximo del neto suavizado (la eficiencia crece y todo
lo que hay debajo del borde decrece, así que donde una deja de crecer le pasa
el mando a la otra); en la corrida de ejemplo cae en \EjUmbral{} mV$\cdot$ns y
es un piso que ninguna ventana de ajuste puede cruzar hacia abajo.

\begin{center}
\captionof{table}{Los pulsos de cada corrida: lo que pasó antes del ajuste.}
\label{tab:pulsos}
{\footnotesize\setlength{\tabcolsep}{3.5pt}\begin{longtable}{lllcrrrrrrr}
\toprule
\textbf{Det.} & \textbf{HV} & \textbf{Fuente} & \textbf{Form.} & \textbf{Eventos} & \textbf{Desc.} & \textbf{Pared} & \textbf{Vivo} & \textbf{Registro} & \textbf{Carga} & \textbf{$Q_{12}/Q_{40}$} \\
\midrule
\endfirsthead
\toprule
\textbf{Det.} & \textbf{HV} & \textbf{Fuente} & \textbf{Form.} & \textbf{Eventos} & \textbf{Desc.} & \textbf{Pared} & \textbf{Vivo} & \textbf{Registro} & \textbf{Carga} & \textbf{$Q_{12}/Q_{40}$} \\
\midrule
\endhead
A & 1550 & Co-60 & XML & 50\,000 & 6.2 & 699 & 40.42 & 158 & 82 & 83 \\
A & 1600 & Co-60 & bin. & 80\,000 & 1.1 & 614 & 43.83 & 81 & 77 & 82 \\
A & 1600 & Cs-137 & XML & 80\,322 & 0.7 & 1373 & 0.12 & 130 & 72 & 82 \\
A & 1650 & Co-60 & bin. & 80\,000 & 1.0 & 592 & 43.38 & 82 & 77 & 83 \\
A & 1650 & Cs-137 & XML & 80\,000 & 0.5 & 1549 & 0.10 & 129 & 77 & 83 \\
A & 1700 & Co-60 & bin. & 80\,000 & 1.3 & 597 & 36.65 & 82 & 76 & 82 \\
A & 1700 & Cs-137 & XML & 80\,000 & 0.5 & 1613 & 0.08 & 130 & 74 & 83 \\
\midrule
B & 1600 & Co-60 & bin. & 80\,000 & 1.2 & 683 & 38.19 & 81 & 77 & 82 \\
B & 1600 & Cs-137 & bin. & 80\,000 & 1.3 & 323 & 0.31 & 81 & 76 & 82 \\
B & 1650 & Co-60 & bin. & 80\,000 & 1.0 & 669 & 36.77 & 81 & 78 & 83 \\
B & 1650 & Cs-137 & bin. & 80\,000 & 1.1 & 305 & 0.30 & 81 & 77 & 83 \\
B & 1700 & Co-60 & bin. & 80\,000 & 1.1 & 639 & 35.23 & 82 & 78 & 83 \\
B & 1700 & Cs-137 & bin. & 80\,000 & 0.9 & 287 & 0.26 & 81 & 79 & 83 \\
\midrule
C & 1600 & Co-60 & bin. & 80\,000 & 1.1 & 928 & 29.77 & 87 & 75 & 82 \\
C & 1600 & Cs-137 & bin. & 80\,000 & 1.3 & 302 & 0.32 & 81 & 76 & 82 \\
C & 1650 & Co-60 & bin. & 80\,000 & 1.1 & 860 & 28.42 & 139 & 72 & 82 \\
C & 1650 & Cs-137 & bin. & 80\,000 & 1.1 & 306 & 0.29 & 81 & 77 & 82 \\
C & 1700 & Co-60 & bin. & 80\,000 & 1.4 & 770 & 38.61 & 139 & 73 & 83 \\
C & 1700 & Cs-137 & bin. & 80\,000 & 0.9 & 307 & 0.28 & 81 & 77 & 83 \\
\midrule
D & 1600 & Co-60 & bin. & 80\,000 & 1.6 & 699 & 51.56 & 47 & 81 & 82 \\
D & 1600 & Cs-137 & bin. & 80\,000 & 2.5 & 306 & 0.47 & 48 & 81 & 82 \\
D & 1600 & Cs-137 (desc.) & bin. & 80\,000 & 4.2 & 306 & 0.58 & 81 & 73 & 81 \\
D & 1650 & Co-60 & bin. & 80\,000 & 1.3 & 645 & 50.45 & 47 & 81 & 82 \\
D & 1650 & Cs-137 & bin. & 80\,000 & 1.7 & 285 & 0.43 & 48 & 81 & 82 \\
D & 1650 & Cs-137 (desc.) & bin. & 80\,000 & 3.0 & 300 & 0.49 & 81 & 75 & 82 \\
D & 1700 & Co-60 & bin. & 80\,000 & 1.1 & 618 & 48.27 & 47 & 81 & 82 \\
D & 1700 & Cs-137 & bin. & 80\,000 & 1.5 & 285 & 0.37 & 47 & 81 & 82 \\
D & 1700 & Cs-137 (desc.) & bin. & 80\,000 & 2.1 & 288 & 0.42 & 81 & 74 & 81 \\
\bottomrule
\end{longtable}
\par\smallskip
\begin{minipage}{0.95\linewidth}\footnotesize HV en V. «Form.»: formato del archivo (XML del DRSOSC o binario del DRS4). «Eventos»: registros en el archivo. «Desc.»: porcentaje descartado por las banderas de calidad (saturacion, compuerta fuera del registro, sin pulso). «Pared»: duracion de la corrida en s; «Vivo»: fraccion del tiempo de pared en que la adquisicion estuvo viva, en \%, medida con el scaler de la placa. «Registro»: cuanto registro queda despues del onset, en ns (mediana). «Carga»: fraccion de la carga del pulso \emph{registrado} que entra en la compuerta de 12 ns, en \% (la misma columna de la tabla~\ref{tab:calib}); $Q_{12}/Q_{40}$: la misma compuerta, relativa a la carga acumulada hasta 40 ns del onset, que todos los registros contienen.\end{minipage}
}
\end{center}

% ===========================================================================
\section{Marco teórico}
\label{sec:fisica}

Esta sección es el fundamento de lo que hace la cadena: cómo interactúa un
fotón con un plástico y por qué eso deja un continuo sin fotopico, de dónde
sale la energía del borde, qué forma tiene el continuo (Klein-Nishina) y por
qué el borde medido no es un escalón; después, los dos algoritmos que el
método usa para encontrarlo sin intervención, la correlación cruzada
normalizada y el recocido simulado, y la estadística con que se le pone una
barra al resultado. La implementación va después
(sección~\ref{sec:cadena}); acá está el porqué.

\subsection{Cómo interactúa un fotón con la materia}
\label{sec:interaccion}

Un fotón de las energías que emiten las fuentes de esta campaña (entre 0.6 y
1.3 MeV) puede hacer tres cosas al atravesar un material, y cuál de las tres
domina depende del número atómico $Z$ del material y de la energía
$E_\gamma$~\cite{knoll2010}:

\begin{itemize}
\item \textbf{Absorción fotoeléctrica.} El fotón desaparece y le entrega toda
      su energía a un electrón ligado (menos su energía de ligadura, que en
      carbono es de unos cientos de eV y acá no cuenta). Es lo que produce el
      \emph{fotopico}. La probabilidad crece brutalmente con $Z$ (como
      $Z^4$ a $Z^5$) y cae rápido con la energía (como $E_\gamma^{-3}$ a
      $E_\gamma^{-3.5}$): es el proceso de los materiales pesados y de las
      energías bajas.
\item \textbf{Dispersión Compton.} El fotón choca con un electrón, le cede
      una parte de su energía y sigue en otra dirección, con menos energía.
      Cuánta parte cede depende del ángulo, y ese reparto es el tema de la
      sección siguiente. La probabilidad por átomo va como $Z$ (un electrón
      es un electrón), así que por gramo es casi la misma en cualquier
      material, y decrece suavemente con la energía.
\item \textbf{Producción de pares.} Por encima de $2 m_e c^2 = 1022$ keV el
      fotón puede convertirse en un par electrón-positrón en el campo de un
      núcleo. Va como $Z^2$ y recién se vuelve importante a varios MeV. Para
      el Co-60 está apenas por encima del umbral y es despreciable.
\end{itemize}

\begin{figure}[htbp]
\centering
\includegraphics[width=\linewidth]{teoria_interaccion.pdf}
\caption{Coeficientes másicos de atenuación por proceso, de la base de datos
XCOM de NIST~\cite{xcom}: (a) el plástico BC-404 (C$_9$H$_{10}$, que XCOM
reparte en \XcomFraccionH{} de hidrógeno y \XcomFraccionC{} de carbono en
peso, las mismas fracciones que usa la simulación), (b) NaI, un centellador
inorgánico de alto $Z$, para comparar. Las líneas punteadas verticales son
las tres líneas gamma de la campaña. En (a), la línea de trazos es la sección
eficaz de Klein-Nishina de un electrón libre multiplicada por los electrones
por gramo del plástico: por encima de 100 keV es indistinguible del
``incoherente'' de la base de datos.}
\label{fig:teoria-interaccion}
\end{figure}

La figura~\ref{fig:teoria-interaccion} lo muestra con datos y no con
argumentos. En el plástico, a la energía del Cs-137 la absorción
fotoeléctrica es $\FotoSobreComptonCs$ veces la Compton, y a la del Co-60,
$\FotoSobreComptonCoB$: el fotoeléctrico deja de dominar en el plástico a
\CruceFotoComptonPlastico{} keV, mientras que en el NaI sigue mandando hasta
\CruceFotoComptonNaI{} keV y a 662 keV todavía vale \FotoSobreComptonNaICs{}
veces la Compton. Un fotón de nuestras fuentes, dentro de un plástico, hace
Compton o no hace nada.

\paragraph{El plástico como bolsa de electrones libres.} Que la dispersión
Compton dependa sólo del número de electrones se puede verificar en la misma
figura: la sección eficaz de Klein-Nishina para \emph{un electrón libre}
(sección~\ref{sec:kn}), multiplicada por los $\ElectronesPorGramo$ electrones
por gramo del C$_9$H$_{10}$, reproduce el ``incoherente'' de XCOM a menos del
\KNvsXcom\,\% entre 100 keV y 3 MeV. Por debajo de 100 keV se separan (hasta
\KNvsXcomBajo\,\% a 10 keV), porque ahí la energía de ligadura de los
electrones ya no es despreciable frente a la que el fotón les puede
transferir. En el rango que importa acá, el electrón libre es una descripción
exacta.

\paragraph{Cuánto camina un fotón.} Con la densidad del BC-404
(\DensidadBC{} g/cm$^3$) el camino libre medio es de \LibreCaminoCs{} cm a
662 keV y \LibreCaminoCoB{} cm a 1332 keV. En los \LadoDetector{} cm del
cilindro simulado (sección~\ref{sec:geant4}), un fotón del Cs-137 tiene un
\ProbInteraccionCs\,\% de probabilidad de interactuar y uno del Co-60 un
\ProbInteraccionCoB\,\%. Dos consecuencias: la mayoría de los fotones
atraviesan el detector sin dejar nada, y el que sí dispersó sale con menos
energía y con una probabilidad del mismo orden de volver a dispersar antes de
escapar. Eso segundo es el \emph{scattering múltiple}, que llena la zona por
encima del borde y que la sección~\ref{sec:lineshape} discute.

\paragraph{Por qué no hay fotopico.} En un NaI, un fotón de 662 keV que
entra puede absorberse de una vez (la absorción es \FotoSobreComptonNaICs{}
veces la dispersión a esa energía) o dispersar
una o varias veces y terminar absorbido: en los dos casos deposita toda su
energía, y esos eventos se apilan en un pico angosto a $E_\gamma$, el
fotopico, que es el punto de calibración de la espectrometría gamma clásica.
En el plástico la absorción es $\FotoSobreComptonCs$ veces menos probable que
la dispersión: el fotón cede una fracción de su energía, variable con el
ángulo, y se va. Lo que se mide es un \emph{continuo}, y no hay ningún evento
en $E_\gamma$. Ese continuo, sin embargo, tiene un extremo bien definido, y
ese extremo es lo que se calibra.

\subsection{Cinemática Compton}
\label{sec:cinematica}

Un fotón de energía $E_\gamma$ choca con un electrón en reposo y sale con
energía $E'$ formando un ángulo $\theta$ con su dirección original. Con la
conservación de la energía y de las dos componentes del momento, en unidades
en que $p = E/c$ para el fotón y usando la relación relativista
$E_e^2 = p_e^2c^2 + m_e^2c^4$ para el electrón, el resultado es una sola
relación entre $E'$ y $\theta$:
\[
   E' \;=\; \frac{E_\gamma}{1 + \alpha\,(1 - \cos\theta)},
   \qquad \alpha = \frac{E_\gamma}{m_e c^2},
   \qquad m_e c^2 = \MasaElectron\text{ keV}.
\]
Es la fórmula de Compton~\cite{compton1923}. Lo que queda en el detector es
lo que perdió el fotón, la energía cinética del electrón de retroceso:
\[
   T \;=\; E_\gamma - E' \;=\; E_\gamma\,
   \frac{\alpha\,(1-\cos\theta)}{1 + \alpha\,(1-\cos\theta)}.
\]
$T$ crece con $\theta$ de manera monótona (panel (a) de la
figura~\ref{fig:teoria-cinematica}) y alcanza su máximo cuando el fotón
rebota hacia atrás, $\theta = 180°$, $\cos\theta = -1$:
\[
   E_{\max} \;=\; T(180°) \;=\; \frac{2\alpha}{1 + 2\alpha}\,E_\gamma .
\]
Ése es el \textbf{borde Compton}: por encima de $E_{\max}$ un solo choque no
puede dejar energía. Su posición depende sólo de $E_\gamma$ y de $m_e c^2$,
no del detector, y por eso sirve para calibrar.

\begin{figure}[htbp]
\centering
\includegraphics[width=\linewidth]{teoria_cinematica.pdf}
\caption{(a) La energía del electrón de retroceso contra el ángulo del fotón
dispersado para las tres líneas de la campaña; el punto a 180° es el borde
Compton de cada una. (b) La sección eficaz de Klein-Nishina en ángulo, en
unidades de $r_0^2$, y el límite de Thomson ($E_\gamma \to 0$): a más
energía, más hacia adelante. (c) La sección eficaz total por electrón contra
la energía: la fórmula analítica, la integral numérica de la expresión
diferencial, y el ``incoherente'' de XCOM para el plástico dividido por sus
electrones por gramo; la línea punteada es la sección eficaz de Thomson,
\SigmaThomson{} barn.}
\label{fig:teoria-cinematica}
\end{figure}

\begin{table}[htbp]
\centering
\caption{Las tres líneas de la campaña: $\alpha$, el borde Compton, qué
fracción de la energía del fotón es, la energía del fotón retrodispersado
(que es lo que la Ec.~1 del paper imprime por error como borde) y la del
fotón dispersado a 90°.}
\label{tab:cinematica}
\small
\begin{tabular}{lcccccc}
\toprule
\textbf{fuente} & \textbf{$E_\gamma$ [keV]} & \textbf{$\alpha$} & \textbf{$E_{\max}$ [keV]} & \textbf{$E_{\max}/E_\gamma$ [\%]} & \textbf{retrodispersión [keV]} & \textbf{$E'$ a 90° [keV]} \\
\midrule
Cs-137 & 661.7 & 1.295 & 477.3 & 72.1 & 184.3 & 288.3 \\
Co-60 & 1173.2 & 2.296 & 963.4 & 82.1 & 209.8 & 356.0 \\
Co-60 & 1332.5 & 2.608 & 1118.1 & 83.9 & 214.4 & 369.4 \\
\bottomrule
\end{tabular}

\end{table}

\paragraph{Dos energías más que conviene tener a mano.} La
tabla~\ref{tab:cinematica} lista, además del borde, la energía del fotón que
rebotó a 180°, $E'(180°) = E_\gamma/(1+2\alpha)$. Es lo que le queda al fotón
después de dejar el borde, y aparece por dos motivos. El primero es una
errata del paper que usamos: su Ec.~1 imprime esa expresión como si fuera
$E_{\max}$, o sea \EcUnoImpresaCs{} keV para el Cs-137 en vez de \BordeCs;
sus propias figuras usan la correcta. El segundo es físico: los fotones que
dispersan a ángulos grandes en lo que rodea al detector (la mesa, el
portafuente, las paredes) vuelven hacia él con energías cerca de
$E'(180°)$, y depositan en el plástico un continuo Compton propio, con su
borde en $T(180°)$ de esa energía menor. Es el \emph{pico de
retrodispersión} de la espectrometría gamma, y es uno de los candidatos al
exceso de cuentas a bajo canal que las corridas reales muestran
(sección~\ref{sec:decisiones}).

\subsection{Dónde cae el borde}

Con $E_{\max} = 2\alpha E_\gamma/(1+2\alpha)$, las fuentes de esta campaña
tienen su borde en \BordeCs{} keV (Cs-137) y en \BordeCoA{} y \BordeCoB{} keV
(las dos líneas del Co-60), tabla~\ref{tab:cinematica}. El Cs-137 aporta un
borde solo; el Co-60, dos.

\begin{aviso}
El Co-60 emite dos líneas gamma de intensidad casi igual, así que su espectro
tiene \emph{dos} bordes Compton solapados. No son dos puntos independientes de
la recta de calibración: los dos salen del mismo $D$, y su separación la fija
la física, no el ajuste. Por eso el Co-60 aporta un punto, no dos, y por eso
hace falta una segunda fuente para testear linealidad.
\end{aviso}

\subsection{Klein-Nishina: cuántos electrones salen a cada ángulo}
\label{sec:kn}

La cinemática dice qué energía tiene el electrón si el fotón salió a
$\theta$; no dice cuán probable es cada $\theta$. Eso lo da la sección eficaz
diferencial de Klein y Nishina~\cite{klein1929}, calculada en 1929 con la
ecuación de Dirac para un electrón \emph{libre y en reposo}:
\[
   \frac{d\sigma}{d\Omega} \;=\; r_0^2\, f^2\,
   \frac{1+\cos^2\theta}{2}
   \left[1 + \frac{\alpha^2 (1-\cos\theta)^2}
                  {(1+\cos^2\theta)\,[1+\alpha(1-\cos\theta)]}\right],
   \qquad f = \frac{1}{1 + \alpha(1-\cos\theta)},
\]
con $r_0$ el radio clásico del electrón. Tres cosas que conviene saber de
ella:

\begin{itemize}
\item \textbf{El límite de baja energía es Thomson.} Con $\alpha \to 0$,
      $f \to 1$ y el corchete $\to 1$, y queda
      $d\sigma/d\Omega = r_0^2(1+\cos^2\theta)/2$, la dispersión clásica de
      una onda electromagnética por una carga libre, simétrica entre adelante
      y atrás. Integrada da $\sigma_T = 8\pi r_0^2/3 = \SigmaThomson$ barn.
      A medida que $\alpha$ crece, el factor $f^2$ suprime los ángulos
      grandes: la dispersión se vuelve hacia adelante (panel (b) de la
      figura~\ref{fig:teoria-cinematica}) y la sección eficaz total baja (a
      662 keV vale \SigmaKNCs{} barn, panel (c)).
\item \textbf{``Libre y en reposo'' es una aproximación, y acá es buena.} Un
      electrón ligado tiene energía de ligadura y momento; la primera corre
      el umbral de la dispersión y el segundo borronea la relación entre
      $\theta$ y $T$ (el \emph{ensanchamiento Doppler} del perfil Compton).
      En carbono e hidrógeno las energías de ligadura son de cientos de eV o
      menos, contra las decenas de keV de resolución del detector: el efecto
      está por debajo de lo que se puede ver, y la
      figura~\ref{fig:teoria-interaccion} lo confirma comparando con la base
      de datos, que sí incluye la ligadura.
\item \textbf{De ángulo a energía: el jacobiano.} Lo que se mide es $T$, no
      $\theta$, así que hay que reescribir la sección eficaz por unidad de
      energía del electrón. Como $T$ es función monótona de $\theta$ y el
      problema tiene simetría azimutal, $d\Omega = 2\pi\,\sin\theta\,d\theta =
      2\pi\,d(\cos\theta)$, y
      \[
         \frac{d\sigma}{dT} \;=\; \frac{d\sigma}{d\Omega}\,
         2\pi\,\left|\frac{d\cos\theta}{dT}\right|,
         \qquad
         \frac{dT}{d\cos\theta} \;=\; -\,\frac{E_\gamma\,\alpha}
                                            {[1 + \alpha(1-\cos\theta)]^2},
      \]
      de modo que $d\sigma/dT = 2\pi\,(d\sigma/d\Omega)\,
      [1+\alpha(1-\cos\theta)]^2/(E_\gamma\alpha)$, con $\cos\theta$
      despejado de $T$. Sustituyendo y ordenando se llega a la expresión en
      energía que usa el paper (sección~\ref{sec:lineshape}). El factor
      $[1+\alpha(1-\cos\theta)]^2$ crece hacia $\theta = 180°$: es el que
      convierte una distribución angular que decrece hacia atrás en una
      distribución en energía que \emph{crece} hacia el borde, la
      acumulación que después va a importar.
\end{itemize}

\subsection{Por qué el borde medido no es un escalón}

La teoría (la sección eficaz de Klein-Nishina) predice un corte abrupto en
$E_{\max}$. Lo que se mide es una caída gradual, por tres razones que se suman:

\begin{itemize}
\item \textbf{La resolución del detector.} El número de fotones de centelleo,
      el número de fotoelectrones que producen y la ganancia de cada etapa del
      fotomultiplicador fluctúan evento a evento. El efecto neto es un
      borroneado aproximadamente gaussiano, y es el parámetro
      $\sigma_{\text{rel}}$ del ajuste.
\item \textbf{El scattering múltiple.} Un fotón puede dispersar dos o tres
      veces adentro del centellador antes de escapar, y entonces deposita más
      energía que la del borde. Eso llena la zona por encima de $E_{\max}$, que
      la teoría de un solo scattering deja vacía. No es gaussiano y no se puede
      absorber subiendo $\sigma$: es la razón principal de usar una plantilla
      de Geant4 en lugar de la fórmula analítica.
\item \textbf{El quenching de Birks.} La luz producida por unidad de energía
      depositada no es constante: baja a energías bajas. Está en la simulación,
      y es lo que separa la escala de energía depositada de la de luz.
\end{itemize}

Ajustar la forma entera contra un modelo que contiene esas tres cosas es lo que
distingue este método de los clásicos (tomar el 50\,\% de la altura del borde,
o el mínimo de la derivada), que sobreestiman sistemáticamente la posición del
borde porque la lineshape real no es un escalón.

\subsection{La lineshape y el modelo directo}
\label{sec:lineshape}

La forma del continuo la da la sección eficaz de Klein-Nishina de la
sección~\ref{sec:kn}, escrita en ángulo sólido para un fotón de energía
$h\nu$ (el paper usa $h\nu$ donde acá se escribió $E_\gamma$).
Lo que se mide, sin embargo, no es el ángulo sino la energía que queda en el
centellador, que es la del electrón de retroceso. Cambiando de variable con el
jacobiano de la sección~\ref{sec:kn} se llega a la forma que usa el paper (su
Ec. 10):
\[
   \frac{d\sigma}{dE} \;=\; \frac{A\,\pi r_0^2}{\alpha^2\, m_0c^2}
   \left[\,2 + \left(\frac{E}{h\nu - E}\right)^{\!2}
   \left(\frac{1}{\alpha^2} - \frac{2\,(h\nu-E)}{\alpha E}
         + \frac{h\nu-E}{h\nu}\right)\right],
   \qquad 0 < E \le E_{\max},
\]
y cero fuera de ese intervalo. Ese corte en $E_{\max}$ \textbf{es} el borde
Compton. La constante $A$ absorbe el prefactor: para ajustar una forma sólo
importa la forma, no la normalización.

\begin{aviso}
Las dos expresiones describen el mismo proceso, así que integrar la primera
sobre ángulo sólido y la segunda sobre energía tiene que dar la misma sección
eficaz total. Se verificó: coinciden a seis decimales para 500, 662, 1000 y
2000 keV. Es la única forma de estar seguro de que el jacobiano está bien.
\end{aviso}

El espectro medido no es esa curva sino su versión deformada por el detector.
El modelo directo tiene dos pasos:

\begin{enumerate}
\item \textbf{Del eje de energía al eje medido.} Con el mapa por el origen,
      $ch = D\,E$, la lineshape evaluada en $ch/D$ es la misma curva
      reexpresada en \mVns. Ése es el único parámetro que relaciona los dos
      ejes, y es el que se busca.
\item \textbf{La resolución.} Se convoluciona con una gaussiana normalizada,
      \[
         G(E) \;=\; \frac{1}{\sigma\sqrt{2\pi}}\,
                    e^{-E^2/(2\sigma^2)},
      \]
      lo que da el modelo completo
      $S(ch) = A\,[\,P(ch/D) \otimes G(\sigma)\,]$. El kernel se normaliza a
      suma 1 y no a integral 1 porque se aplica sobre un histograma: así la
      convolución conserva las cuentas y $A$ y $\sigma$ quedan desacoplados
      (ensanchar no cambia el número total de eventos).
\end{enumerate}

En este informe $P$ no es la fórmula analítica sino el histograma de la
simulación de Geant4, que ya trae adentro el scattering múltiple, el escape y
el quenching. La fórmula de arriba es lo que la simulación reproduce cuando
sólo hay un scattering, y sirve para entender qué forma tiene el borde y por
qué el corte es abrupto antes de la convolución.

\begin{figure}[htbp]
\centering
\includegraphics[width=\linewidth]{teoria_ncc.pdf}
\caption{(a) La lineshape de Klein-Nishina para las dos fuentes: plateau,
acumulación creciente y corte abrupto en el borde (líneas punteadas). El Co-60
muestra sus dos bordes. (b) La misma curva convolucionada con la resolución
medida en una de nuestras corridas: el corte se convierte en una caída gradual
y el segundo borde del Co-60 casi desaparece como estructura separada. (c) El
barrido del NCC sobre esa misma corrida: $\rho$ contra el canal de prueba del
borde, con un máximo claro, y en línea de trazos el borde que devuelve el
ajuste completo.}
\label{fig:teoria-ncc}
\end{figure}

\subsection{Los métodos clásicos y su sesgo}
\label{sec:clasicos}

Antes de que existiera la costumbre de ajustar la lineshape entera, la posición
del borde se leía directamente del espectro. Conviene conocer esas recetas por
dos razones: son la vara contra la que el paper mide su método (su Fig.~4), y
lo que falla en ellas es exactamente lo que el ajuste de la sección
anterior corrige. Tres, de la más ingenua a la menos:

\begin{itemize}
\item \textbf{El 50\,\% de altura.} Se toma el canal en que el flanco
      descendente cae a la mitad del máximo. Es lo que se hace a ojo, y fue
      la práctica habitual hasta que Knox y Miller~\cite{knox1972} midieron,
      con electrones de energía conocida en un NE-213, que el borde real no
      cae a la mitad sino cerca del 89\,\% de la altura. Dietze y
      Klein~\cite{dietze1982} cerraron la discusión comparando el espectro con
      una simulación Monte Carlo convolucionada con la resolución, que es en
      esencia lo que hace este informe.
\item \textbf{El mínimo de la derivada} (el ``método de la diferencia'' del
      paper). El borde es el canal donde el espectro cae más rápido, o sea el
      mínimo de $dN/d\text{canal}$. Sobre datos con ruido la derivada hay que
      calcularla sobre el espectro suavizado, y cuánto se suaviza pasa a ser
      un parámetro del método: poco, y el mínimo lo elige el ruido; mucho, y el
      flanco se corre.
\item \textbf{El escalón gaussiano} (el ``método FWHM gaussiano'' del paper).
      Se ajusta al flanco un escalón borroneado por una gaussiana, una función
      error complementaria más un fondo lineal, y se toma el punto medio del
      escalón ajustado. De paso el ajuste entrega el ancho, que es de donde
      sale el nombre.
\end{itemize}

\paragraph{Por qué los tres se equivocan, y en la misma dirección.} Si el
borde fuera un escalón perfecto borroneado por una gaussiana, los tres darían
exactamente lo mismo: la derivada de ese escalón es una gaussiana centrada en
el escalón, y su punto de media altura cae ahí también. Pero la lineshape de
Klein-Nishina no es un escalón (figura~\ref{fig:teoria-ncc}, panel a): tiene
una \emph{acumulación creciente} antes del borde. Al borronearla, esa pendiente
de fondo corre el punto de media altura, el mínimo de la derivada y el centro
del escalón ajustado hacia arriba en canal, y tanto más cuanto más ancha es la
resolución. Ninguna de las tres recetas contiene esa forma, así que ninguna
puede corregirla.

\begin{table}[htbp]
\centering
\caption{El sesgo de los tres métodos clásicos contra un borde de posición
conocida, en función de la resolución. El borde es la lineshape de
Klein-Nishina de la línea del Cs-137 (\ClasEgamma{} keV, borde en
\ClasEmax{} keV) convolucionada con una gaussiana de ancho $\sigma$. Las tres
primeras columnas llevan ruido de Poisson; las tres marcadas ``s/r'' son el
mismo cálculo sin ruido. Generada con \texttt{codigo/etapas/clasicos.py}.}
\label{tab:clasicos}
\footnotesize
\begin{tabular}{lcccccc}
\toprule
\textbf{$\sigma/B$} & \textbf{50\,\%} & \textbf{derivada} & \textbf{escalón} & \textbf{50\,\% s/r} & \textbf{derivada s/r} & \textbf{escalón s/r} \\
\midrule
0.03 & $+0.7$ & $+0.7$ & $+2.3$ & $+0.7$ & $+0.3$ & $+2.3$ \\
0.06 & $+2.0$ & $+1.5$ & $+4.0$ & $+2.0$ & $+1.2$ & $+4.0$ \\
0.09 & $+3.5$ & $+2.8$ & $+5.2$ & $+3.5$ & $+2.2$ & $+5.2$ \\
0.12 & $+4.9$ & $+3.4$ & $+6.0$ & $+4.9$ & $+3.3$ & $+6.0$ \\
0.15 & $+6.2$ & $+6.1$ & $+6.6$ & $+6.2$ & $+4.4$ & $+6.6$ \\
0.18 & $+7.5$ & $+7.1$ & $+7.1$ & $+7.5$ & $+5.4$ & $+7.1$ \\
\bottomrule
\end{tabular}

\end{table}

La tabla~\ref{tab:clasicos} lo mide. A $\sigma/B = \ClasSigmaMin$ los tres
están a menos del 3\,\%; a $\sigma/B = \ClasSigmaMax$ el 50\,\% de altura
sobreestima el borde en \ClasMitadHi\,\%, el mínimo de la derivada en
\ClasDifHi\,\% y el escalón gaussiano en \ClasGaussHi\,\%. Y las tres
columnas de la derecha son el mismo cálculo \emph{sin ruido}: el 50\,\% de
altura y el escalón gaussiano dan lo mismo con y sin ruido, y el mínimo de la
derivada sigue en \ClasDifSinHi\,\%. \textbf{El sesgo es del método, no de la
estadística}: más eventos no lo achican. La figura~\ref{fig:teoria-clasicos}
muestra el caso con la resolución de nuestra corrida de ejemplo, y la banda
gris del panel (b) es el rango de resoluciones de esta campaña, donde los tres
métodos están entre el 3 y el 7\,\% arriba.

\begin{figure}[htbp]
\centering
\includegraphics[width=\linewidth]{teoria_clasicos.pdf}
\caption{(a) Un borde sintético con la resolución de la corrida de ejemplo
($\sigma/B$ como en \EjCorrida): el borde verdadero y dónde lo pone cada
método clásico. (b) El sesgo de cada método contra $\sigma/B$, con ruido de
Poisson (líneas llenas) y sin ruido (trazos); la banda gris es el rango de
resoluciones de esta campaña.}
\label{fig:teoria-clasicos}
\end{figure}

\paragraph{Sobre las corridas reales.} Corriendo los tres métodos sobre las
\ClasNReales{} corridas de la campaña, con la región centrada en el borde que
estimó el NCC, y comparando con la $D$ del ajuste completo: el 50\,\% de
altura queda \ClasRealMitadMed\,\% arriba (mediana; entre \ClasRealMitadMin{}
y \ClasRealMitadMax\,\%), el mínimo de la derivada \ClasRealDifMed\,\% (entre
\ClasRealDifMin{} y \ClasRealDifMax\,\%) y el escalón gaussiano
\ClasRealGaussMed\,\% (entre \ClasRealGaussMin{} y \ClasRealGaussMax\,\%). Son
sesgos mayores que los de la tabla a la misma resolución, y la diferencia
tiene una explicación que la tabla no puede tener: el espectro real lleva
scattering múltiple, que llena la zona por encima del borde (sección
siguiente), y eso corre los tres estimadores todavía más hacia arriba. La
plantilla de Geant4 contiene ese scattering; la lineshape de Klein-Nishina
del sintético, no.

\begin{aviso}
Los tres métodos necesitan que se les diga \emph{dónde} está el borde, y la
región no puede ser la ventana de ajuste del modelo, que en esta campaña
arranca mucho más abajo. El exceso de cuentas a bajo canal cae más
empinado que el propio borde: si la región lo incluye, el mínimo de la
derivada se va ahí y devuelve un número sin relación con el borde. Con la
región anclada al borde del NCC eso no pasa. Es un recordatorio de que
``automático'' no quiere decir ``sin supuestos'': también estos métodos
tienen una ventana, sólo que nadie la escribe.
\end{aviso}

\subsection{Correlación cruzada normalizada (NCC)}
\label{sec:ncc}

\paragraph{El problema.} Para ajustar el borde hay que decir en qué tramo del
espectro se ajusta. Elegirlo a mano es subjetivo y el resultado depende de la
elección; peor todavía, es una decisión que se toma mirando el mismo dato que
después se ajusta. El NCC la automatiza.

\paragraph{La idea.} Comparar la \emph{forma} del espectro medido con la de la
plantilla, sin que importen ni la escala vertical ni un posible pedestal. La
medida es el coeficiente de correlación de Pearson entre las dos curvas
muestreadas en los mismos puntos:
\[
   \rho(s_1, s_2) \;=\;
   \frac{\sum_i (s_{1,i}-\bar{s_1})\,(s_{2,i}-\bar{s_2})}
        {\sqrt{\sum_i (s_{1,i}-\bar{s_1})^2 \;\sum_i (s_{2,i}-\bar{s_2})^2}}
   \;\in\; [-1, 1].
\]

\paragraph{Por qué es inmune al offset y a la escala.} Si se reemplaza
$s_2 \to a\,s_2 + b$ con $a>0$, la media se transforma igual
($\bar{s_2} \to a\bar{s_2}+b$), así que las desviaciones respecto de la media
se multiplican por $a$ en el numerador y en el denominador, y el factor se
cancela; el corrimiento $b$ desaparece al restar la media. Es decir
$\rho(s_1, a s_2 + b) = \rho(s_1, s_2)$. Eso es exactamente lo que hace falta
cuando la amplitud del espectro todavía no se conoce: se compara forma y nada
más.

\paragraph{Qué se barre, y por qué un solo parámetro.} La tentación es barrer
el principio y el fin de la ventana por separado, como hace el paper. Sobra:
si el mapa canal-energía pasa por el origen, fijado el canal del borde $B$ la
ventana ya está determinada,
$[\,\texttt{inicio\_rel}\cdot B,\ \texttt{fin\_rel}\cdot B\,]$. Barrer el ancho
aparte agrega un grado de libertad que el modelo no tiene, y sobre datos reales
se lo gasta eligiendo ventanas degeneradas: en la cola vacía del espectro, donde
casi no hay cuentas, el estandarizado amplifica ruido puro hasta parecerse a
cualquier forma y $\rho$ se dispara. Acá se barre sólo $B$, sobre un rango
derivado de dónde el espectro suavizado cae a la mitad de su plateau, y gana el
$B$ de mayor $\rho$. El panel (c) de la figura~\ref{fig:teoria-ncc} es ese
barrido.

\paragraph{Dos detalles que cambian el resultado.} El espectro se suaviza con
un filtro de Savitzky-Golay~\cite{savitzky1964}, que en cada punto ajusta un
polinomio de grado bajo a una ventana de bines vecinos y toma el valor del
polinomio en el centro. A diferencia de un promedio móvil, un polinomio de
tercer grado puede seguir una pendiente y una curvatura, así que el filtro
quita el ruido \emph{sin achatar el flanco}, y el flanco es justamente lo que
se está midiendo. Y las dos curvas se remuestrean a un número fijo de puntos
dentro de la ventana, fijo porque si cambiara con $B$ los $\rho$ de ventanas
distintas dejarían de ser comparables entre sí.

\paragraph{Cuántos puntos.} El paper remuestrea a 100 puntos y no dice por
qué. Acá se usa un punto por bin de la ventana (\EjNptsAuto{} en la corrida de
ejemplo), y la figura~\ref{fig:ncc-detalle}(c) muestra qué pasa al cambiarlo:
entre \NptsMin{} y \NptsMax{} puntos el $B$ ganador se mueve como máximo
\NptsBDesvMax\,\%, o sea que el resultado no depende de esa elección. Lo que
sí cambia es $\rho$: con \NptsMin{} puntos vale \NptsRhoPocos{} y con un punto
por bin, \NptsRhoAuto. Remuestrear a pocos puntos \emph{promedia} ruido y
sube la correlación sin que la forma se parezca más; a más puntos que bines
no se agrega información, sólo puntos correlacionados. Como $\rho$ es el
criterio con que se comparan ventanas, conviene no inflarlo, y por eso el
número de puntos no se inventa: es el que el histograma tiene.

\paragraph{Lo que Pearson ve.} La figura~\ref{fig:ncc-detalle} muestra las
dos curvas exactamente como entran a la ecuación: estandarizadas (media cero,
desviación uno) dentro de la ventana ganadora. En el panel (a) la plantilla
sigue al espectro en el flanco y se separa de él por debajo del borde, donde
el espectro tiene más cuentas que la plantilla; en el panel (b), una contra la
otra, esa separación es el gancho que se aparta de la diagonal en la zona de
bajo canal. Todo lo demás cae sobre la diagonal, y eso es lo que $\rho =
\EjRhoRecalc$ está midiendo. Ese $\rho$ se recalcula acá a partir de la
ventana, y coincide con el del barrido (\EjRho): la figura reconstruye lo que
el barrido hizo, no otra cosa.

\begin{figure}[htbp]
\centering
\includegraphics[width=\linewidth]{teoria_ncc_detalle.pdf}
\caption{El NCC por dentro, sobre \EjCorrida. (a) El espectro suavizado y la
plantilla llevada a canal con el $B$ ganador, las dos estandarizadas dentro
de la ventana: es lo único que Pearson compara. (b) Una contra la otra, con
el color indicando el canal; los puntos sobre la diagonal son donde las dos
formas coinciden. (c) El $\rho$ máximo del barrido y el $B$ ganador contra el
número de puntos de remuestreo; la línea punteada es un punto por bin, que
es lo que se usa.}
\label{fig:ncc-detalle}
\end{figure}

\paragraph{Qué entrega, paso a paso.} El barrido recorre \NBarrido{} valores
del canal de prueba $B$, entre 0.4 y 1.1 veces el canal donde el espectro
suavizado cae a la mitad de su plateau. Para cada $B$ hace cuatro cosas:
(i)~lo convierte en una ganancia de prueba, $D = B/E_{\text{ref}}$, con
$E_{\text{ref}}$ el borde de la primera línea gamma: como el mapa pasa por el
origen, el canal del borde de referencia \emph{es} $D\,E_{\text{ref}}$, así que
fijar $B$ es fijar $D$; (ii)~lleva la plantilla de Geant4 a canal con esa $D$ y
la borronea con un $\sigma$ relativo fijo (\SigmaRelNCC); (iii)~arma la ventana
$[\InicioRel\,B,\ f\,B]$, con $f = E_{\text{último}}/E_{\text{ref}} + 3\sigma$
para que entren el último borde de la fuente y su cola; (iv)~remuestrea el
espectro suavizado y la plantilla en esa ventana a un número fijo de puntos y
calcula $\rho$. Gana el $B$ de mayor $\rho$, y de ese único número salen las
tres cosas que el ajuste necesita:

\begin{enumerate}
\item \textbf{La ventana de ajuste}, $[\InicioRel\,B,\ f\,B]$. Es la que define
      el $\chi^2$: qué bines cuentan y cuáles no. El final de la subida del
      trigger, que se detecta en el propio espectro, es un piso que ninguna
      ventana puede cruzar hacia abajo.
\item \textbf{La semilla}, $D_0 = B/E_{\text{ref}}$, acompañada del $\sigma$
      con que se hizo el barrido. No es una estimación de la resolución (ver
      más abajo); es el punto de referencia contra el que se mide después
      cuánto tuvo que corregir el ajuste.
\item \textbf{Las cotas de la búsqueda}: $D \in [\CotaDMin\,D_0,\ \CotaDMax\,D_0]$
      y $\sigma_{\text{rel}} \in [\CotaSigmaMin,\ \CotaSigmaMax]$. La de $D$ es
      la única que depende de la semilla; la de $\sigma$ es física (una
      resolución entre 2 y 30\,\%).
\end{enumerate}

\paragraph{Cómo entra eso al recocido.} Conviene ser preciso, porque la palabra
\emph{semilla} sugiere que el recocido arranca de $D_0$, y no es así. Los
\NArranques{} recocidos del multiarranque parten de puntos elegidos al azar,
uniformes dentro de las cotas, y eso es a propósito: si todos arrancaran de la
semilla, que terminen juntos no probaría que el mínimo es el global
(sección~\ref{sec:sa}). Lo que $D_0$ fija es otra cosa:
\begin{itemize}
\item \textbf{Dónde se busca.} Las cotas son toda la información sobre $D$ que
      el recocido recibe, y son anchas a propósito ($-40$ a $+60$\,\% de
      $D_0$): la semilla sale de un barrido sin pesos, sobre el espectro
      suavizado y con un $\sigma$ supuesto, así que puede venir sesgada unos
      por ciento, y una cota estrecha dejaría el óptimo afuera sin avisar.
\item \textbf{Con qué paso.} La escala de cada propuesta es el 10\,\% del
      ancho de su cota (el $0.1\,d_i$ del paper); para $D$ eso es $0.1\,D_0$.
      La temperatura inicial se mide alrededor del punto de partida de cada
      arranque con ese mismo paso (sección~\ref{sec:sa}), así que también
      hereda la escala de la semilla.
\item \textbf{Qué se minimiza.} El $\chi^2$ de Pearson dentro de la ventana que
      eligió el NCC, con la amplitud proyectada en forma cerrada. Otra ventana
      es otro $\chi^2$: por eso la elección automática del intervalo es parte
      del resultado y no un preámbulo.
\item \textbf{Contra qué se compara.} La semilla se evalúa en ese mismo
      $\chi^2$ y se guarda junto con el resultado; la distancia entre $D_0$ y
      la $D$ final mide cuánto corrigió el ajuste con pesos a lo que la forma
      sola decía.
\end{itemize}
Cada recocido enfría con $T_{k+1} = \Alfa\,T_k$ y \NMarkov{} propuestas por
temperatura, y corta cuando el mejor costo deja de mejorar. Del mejor de los
\NArranques{} arranques se baja al fondo exacto del valle con un optimizador
local (Nelder--Mead): el recocido sin ese pulido da el tercer decimal, y el
pulido sin recocido, el valle equivocado. Con el exceso de bajo canal modelado
(sección~\ref{sec:decisiones}) el esquema es el mismo con un parámetro más:
la ventana se extiende hacia abajo hasta $\InicioAjusteRel\,B$, y $\lambda$
recibe como semilla el \LambdaSemillaRel\,\% del ancho de la ventana y como
cotas \CotaLambdaMin~\mVns{} y la mitad de ese ancho.

\paragraph{En números.} En la corrida del panel (c) de la
figura~\ref{fig:teoria-ncc} (\EjCorrida) el barrido dio $B = \EjB$~\mVns{} con
$\rho = \EjRho$, la ventana $[\EjVentanaLo,\ \EjVentanaHi]$~\mVns{} y la semilla
$D_0 = \EjSemillaD$; el recocido, buscando entre \EjCotaDMin{} y \EjCotaDMax,
terminó en $D = \EjD$, a \EjDesvSemilla\,\% de la semilla. Sobre las
\NCorridas{} corridas, el intervalo ganador correlaciona entre $\rho = \RhoMin$
y $\rho = \RhoMax$, y la semilla queda a \SemillaMediana\,\% (mediana) de la
$D$ final, \SemillaMax\,\% en la peor.

\paragraph{Qué no puede hacer.} Compara forma contra una plantilla, así que no
tiene manera de saber que a bajo canal hay cuentas que la plantilla no
contiene: si el exceso está ahí, el barrido elegirá una ventana que lo
incluya. Y el $\sigma$ con que se construye la plantilla del barrido no es una
medida de la resolución: el NCC correlaciona sobre el espectro suavizado y sin
pesos, y además el fin de la ventana depende de ese $\sigma$, así que al
barrerlo se comparan $\rho$ de ventanas distintas. La resolución la determina
el ajuste con pesos.

\subsection{Recocido simulado (SA)}
\label{sec:sa}

\paragraph{El problema.} Un optimizador local sólo sabe bajar: desde donde
arranca camina hasta el fondo del valle en el que cayó, y no tiene forma de
saber, desde adentro, si hay otro valle más profundo. El resultado depende del
punto inicial. El panel (a) de la figura~\ref{fig:teoria-sa} muestra el caso.

\paragraph{La idea.} Aceptar, a veces, moverse a un punto peor. Es la única
manera de subir una loma y ver qué hay del otro lado. El nombre, y la idea de
usar la física estadística como optimizador, vienen de Kirkpatrick, Gelatt y
Vecchi~\cite{kirkpatrick1983}, y la metáfora es de la metalurgia: enfriar despacio un metal deja que los átomos encuentren la
configuración de mínima energía, mientras que enfriarlo de golpe los congela en
un vidrio, en un mínimo local.

\paragraph{El criterio de Metropolis.} Es la regla de aceptación de
Metropolis \emph{et al.}~\cite{metropolis1953}, treinta años anterior al
recocido, inventada para muestrear configuraciones de un sistema a
temperatura fija. Desde el punto actual $x$ se propone
$x' = x + \Delta$, con $\Delta$ uniforme en una vecindad, y se acepta con
probabilidad
\[
   P(x \to x') \;=\;
   \min\!\left[\,1,\; e^{-\Delta E / T}\,\right],
   \qquad \Delta E = E(x') - E(x),
\]
o sea siempre que mejore, y con probabilidad $e^{-\Delta E/T}$ si empeora.

\paragraph{Por qué esa regla y no otra.} Porque es la que hace que la cadena,
a temperatura fija, muestree la distribución de Boltzmann
$\pi(x) \propto e^{-E(x)/T}$. La regla satisface balance detallado respecto de
$\pi$: para cualquier par de estados, $\pi(x)P(x\to x') = \pi(x')P(x'\to x)$,
y de ahí que $\pi$ sea la distribución estacionaria de la cadena. Lo que
importa de eso es el límite: cuando $T \to 0$, $\pi$ se concentra en el mínimo
global de $E$. Enfriar es, entonces, apretar la distribución alrededor de la
respuesta.

\paragraph{El enfriamiento, y por qué no hay garantía.} Acá se usa el
enfriamiento geométrico del paper, $T_{k+1} = \alpha\,T_k$ con $\alpha$ cerca
de 1, con una cadena de Markov de varias propuestas a cada temperatura. El
teorema de convergencia del recocido simulado~\cite{geman1984} pide, en
cambio, un enfriamiento \emph{logarítmico}, $T_k \gtrsim c/\log(1+k)$, que es
tan lento que en la
práctica no se usa (panel (c) de la figura~\ref{fig:teoria-sa}: la diferencia
entre los dos es de órdenes de magnitud). O sea que el recocido geométrico
\textbf{no} garantiza el óptimo global; lo que hay es evidencia empírica, y por
eso el multiarranque de más abajo no es un lujo.

\paragraph{La temperatura inicial no se puede copiar.} $T$ se compara contra
$\Delta E$, así que está en las unidades del costo. El paper fija
$T_{\text{init}} = 2000$ sin decir en qué unidades queda su función de costo,
y ese número no es transferible: con un $\chi^2$ de orden mil anda bien, y con
un costo normalizado a 1 acepta absolutamente todo y la cadena se vuelve una
caminata al azar. Acá $T_{\text{init}}$ se \emph{mide}: se tiran propuestas al
azar alrededor del punto inicial, se promedia cuánto empeoran, y se despeja
\[
   T_{\text{init}} \;=\; \frac{-\langle \Delta E_+ \rangle}{\ln p},
\]
con $p$ la fracción de esos empeoramientos que uno quiere aceptar al principio.
Lo mismo con el criterio de corte, que acá es relativo y no absoluto.

\paragraph{La tasa de aceptación es el termómetro.} Si al principio se acepta
casi todo, la cadena está paseando al azar y la temperatura es demasiado alta;
si no se acepta casi nada, ya se comporta como un optimizador local y sobra el
recocido. Lo que se espera ver es una tasa alta al principio que cae al
enfriarse, y eso se puede mirar corrida por corrida.

\paragraph{Los parámetros lineales no se buscan.} La amplitud $A$ entra
linealmente en el modelo, así que para cada valor de $(D, \sigma)$ la amplitud
óptima se despeja en forma cerrada por mínimos cuadrados pesados. Proyectarla
fuera (la proyección de variables de Golub y Pereyra~\cite{golub1973}, VARPRO;
la fórmula está en la sección~\ref{sec:estadistica}) tiene dos efectos: el recocido
busca en dos dimensiones en vez de tres, y desaparece el problema de escala
entre parámetros, porque una amplitud de miles de cuentas y una ganancia de
orden 1 no admiten un mismo paso de propuesta.

\paragraph{El multiarranque es la evidencia.} Como no hay garantía teórica, la
prueba práctica de que el mínimo encontrado es el global es correr
\NArranques{} recocidos independientes desde puntos iniciales dispersos dentro
de las cotas y ver que todos terminan en el mismo lugar. Arrancarlos todos
desde la semilla no probaría nada. En esta campaña la dispersión entre
arranques es de \DispersionMediana\,\% de $D$ en la corrida mediana y no pasa
de \DispersionMax\,\% en la peor.

\paragraph{Lo que se verificó del recocido, y dónde.} Las afirmaciones de
esta sección no se dan por sabidas: \texttt{analisis/verificar\_sa.py}
comprueba sobre nuestro propio $\chi^2$ que la cadena de Metropolis a
temperatura fija muestrea Boltzmann, que $T_{\text{init}} = 2000$ del paper
anda con su escala de costo y falla con otra, que el paso fijo y el
adaptativo dan la misma dispersión, que el objetivo de aceptación no mueve
$D$, y que proyectar $A$ no cambia el óptimo. El detalle, con el código
comentado, está en \texttt{docs/recocido\_simulado.pdf} y
\texttt{docs/recocido\_simulado\_codigo.pdf}.

\begin{figure}[htbp]
\centering
\includegraphics[width=\linewidth]{teoria_sa.pdf}
\caption{(a) Por qué hace falta: una función con varios valles, un optimizador
local que se queda en el primero que encuentra (estrella) y el recorrido del
recocido, coloreado por iteración, que escapa y llega al mínimo global.
(b) El criterio de Metropolis para tres temperaturas: a $T$ alta se acepta casi
cualquier empeoramiento, a $T$ baja sólo los muy chicos. (c) Los dos
enfriamientos: el geométrico que se usa en la práctica y el logarítmico que
exige el teorema de convergencia.}
\label{fig:teoria-sa}
\end{figure}

\subsection{Estadística del ajuste}
\label{sec:estadistica}

El recocido entrega el mínimo del $\chi^2$; lo que convierte ese mínimo en un
resultado científico es lo que viene después: con qué peso entró cada bin,
qué tan hondo es el valle, y si el número que se reporta como error es el
error de verdad. Esta sección recorre eso en orden, y termina con la única
prueba que cierra la pregunta: meterle a la cadena un espectro de respuesta
conocida y ver qué devuelve.

\paragraph{Poisson, bin por bin.} El número de cuentas en un bin es una
variable de Poisson: si en promedio caen $\mu$ eventos, la varianza es también
$\mu$, y la barra de error es $\sqrt{\mu}$. Como $\mu$ no se conoce, lo
habitual es usar el número observado, $\sigma_i = \sqrt{n_i}$. Cuando hay
fondo restado el neto es $k_m n_m - k_f n_f$, con $k$ los factores de tiempo
vivo, y las varianzas se suman con los factores al cuadrado,
$\sigma^2 = k_m^2 n_m + k_f^2 n_f$. Un piso de una cuenta evita que un bin
vacío tenga varianza cero y, por lo tanto, peso infinito.

\paragraph{Por qué pesar.} El ajuste minimiza
\[
   \chi^2(\theta) \;=\; \sum_i \frac{\big[n_i - \mu_i(\theta)\big]^2}{\sigma_i^2},
\]
y el denominador no es un adorno. Sin él (mínimos cuadrados planos) el
plateau, que tiene miles de cuentas por bin, domina la suma, y el flanco, que
es donde vive la información sobre $D$, pesa poco; además no habría manera de
decir si el ajuste es bueno ni de ponerle una barra a $D$. Con las barras de
Poisson, minimizar $\chi^2$ es, para bines con muchas cuentas, lo mismo que
maximizar la verosimilitud~\cite{cowan1998}, y $\chi^2/\text{ndf}$ cerca de
uno dice que el modelo describe los datos dentro de su fluctuación
estadística.

\paragraph{Neyman contra Pearson.} Hay dos maneras de escribir $\sigma_i^2$,
y no dan lo mismo~\cite{baker1984}. La de arriba, con el número
\emph{observado}, es el $\chi^2$ de Neyman; la otra usa el número
\emph{predicho} por el modelo, $\sigma_i^2 = \mu_i(\theta)$, y es el
$\chi^2$ de Pearson. Con Neyman, un bin que fluctuó hacia abajo recibe una
barra más chica y por lo tanto más peso, y el ajuste se deja arrastrar hacia
abajo; el sesgo va como el inverso del número de cuentas por bin. Con Pearson
la varianza sigue al modelo y ese sesgo desaparece, a cambio de que la
amplitud ya no se despeja en un paso (la varianza depende de ella) y hay que
iterar. Esta campaña se ajusta con Neyman, y cuánto cuesta eso se mide más
abajo, con verdad conocida.

\paragraph{La amplitud se despeja.} $A$ multiplica al modelo, así que para
cada $(D, \sigma)$ el $\chi^2$ es una parábola en $A$ y su mínimo tiene
fórmula cerrada: $A = \sum_i w_i f_i n_i / \sum_i w_i f_i^2$, con $f$ la
plantilla convolucionada y $w = 1/\sigma^2$. Sacar los parámetros lineales
del problema no lineal es lo que Golub y Pereyra~\cite{golub1973} llamaron
proyección de variables, y acá compra dos cosas: el recocido busca en dos
dimensiones en vez de tres, y desaparece el problema de escala entre una
amplitud de miles de cuentas y una ganancia de orden uno.

\paragraph{La barra de error: curvatura y perfil.} Si el valle del $\chi^2$
es angosto, moverse poco empeora mucho el ajuste y el parámetro está bien
determinado; si es ancho, no. Hay dos maneras de medir el ancho. La rápida es
la \emph{curvatura}: se aproxima el mínimo por una parábola, la matriz de
segundas derivadas (el hessiano) se invierte, y la raíz de la diagonal de
$2H^{-1}$ es la barra de cada parámetro. La honesta es el \emph{perfil}: para
cada valor de $D$ se reoptimiza todo lo demás y se anota cuánto subió el
$\chi^2$; la barra es donde $\Delta\chi^2 = 1$. Las dos coinciden si el
mínimo es cuadrático, y esa es una hipótesis que hay que mirar, no suponer.
El panel (a) de la figura~\ref{fig:teoria-estadistica} las superpone para la
corrida de ejemplo: la curvatura da $\pm\EjErrCurv$ y el perfil
$\pm\EjErrPerfil$ \DkeV, o sea que acá el valle es cuadrático y la barra
rápida vale. Cuando no coinciden gana el perfil, y la diferencia es
información: dice que el $\chi^2$ es más plano lejos del mínimo de lo que su
curvatura sugiere, típicamente porque el parámetro se confunde con otro.

\paragraph{La correlación.} La matriz de covarianza tiene términos fuera de
la diagonal, y en este ajuste no son chicos: $D$ y $\sigma$ están
correlacionados entre \CorrMin{} y \CorrMax{} según la corrida. El panel (b)
es el contorno de $\Delta\chi^2 = 1$ en el plano $(D, \sigma)$: una elipse
inclinada, no un círculo. Quiere decir que subir $D$ y bajar $\sigma$ (o al
revés) deja el $\chi^2$ casi igual, porque un borde más ancho centrado más
abajo se parece a uno más angosto centrado más arriba. Citar los dos errores
por separado, sin la correlación, sobreestima la precisión de cualquier
cantidad que dependa de los dos.

\paragraph{El escalado por $\chi^2/\text{ndf}$.} El criterio
$\Delta\chi^2 = 1$ supone que el modelo es correcto y que las barras son las
de Poisson. Sobre los datos reales el $\chi^2/\text{ndf}$ es \EjChiNdf{} en
la corrida de ejemplo: el modelo no describe todo lo que hay (el exceso de
bajo canal, sección~\ref{sec:decisiones}), y entonces la barra de
$\Delta\chi^2 = 1$ es demasiado optimista. El parche estándar es multiplicar
la barra por $\sqrt{\chi^2/\text{ndf}}$, o lo que es lo mismo, leer el perfil
en $\Delta\chi^2 = \chi^2/\text{ndf}$, que es la segunda línea horizontal del
panel (a). Con eso la barra reportada pasa de \EjErrCurv{} a
\EjErrEscalado{} \DkeV. Hay que decirlo como lo que es: una corrección
empírica que reparte la falta de ajuste entre todos los bines por igual, no
una medición de qué parte de esa falta afecta a $D$.

\paragraph{Lo que $\chi^2/\text{ndf}$ no dice.} Un $\chi^2/\text{ndf}$ cerca
de uno dice que el modelo describe los datos dentro del ruido. No dice que
sea el modelo correcto: dos modelos distintos pueden describir igual de bien
y dar $D$ distintas, y entonces el $\chi^2$ no puede elegir entre ellos. En
este proyecto pasó con la forma del exceso de bajo canal (exponencial contra
ley de potencia, que dan $D$ distintas con el mismo $\chi^2/\text{ndf}$) y
con la definición del área, que el $\chi^2$ tampoco
distingue. Lo que separa esos casos no es estadística sino física, o una
medición más.

\paragraph{Sistemáticos: comunes y punto a punto.} Todo lo que se decide y no
se ajusta (la compuerta, dónde arranca la ventana, qué plantilla, cómo se
modela el exceso) es una fuente de error que la barra estadística no
contiene. Se mide de una sola manera: cambiando la decisión dentro de lo
razonable y viendo cuánto se mueve $D$. Para el test de linealidad importa
además de qué tipo es cada uno: un sistemático \emph{común} a todas las
corridas (la escala de la plantilla, la compuerta si es la misma en todas)
mueve todas las $D$ juntas, se va en la pendiente de la recta y no ensucia
los residuos; uno \emph{punto a punto} (la ventana, que el NCC elige corrida
por corrida) sí, y es el que hay que comparar contra los residuos de la
sección~\ref{sec:testA}.

\paragraph{El cierre con verdad conocida.} Todo lo anterior son argumentos.
La prueba es otra: tomar el ajuste de una corrida real como verdad, fabricar
el espectro que ese modelo predice, muestrearlo con ruido de Poisson a la
estadística real, y pasar cada realización por la cadena \emph{completa},
NCC, recocido, pulido y curvatura, como si fuera un archivo nuevo. Si la
cadena devuelve la $D$ que se le metió, dentro de la barra que ella misma
reporta, en la fracción de veces que corresponde, entonces el número y su
error son creíbles. Si no, no hay argumento que lo salve. Fuera de la ventana
de ajuste (el exceso, la subida del umbral, la cola de apilamiento) el
espectro sintético lleva la forma del dato real, porque la cadena tiene que
elegir la ventana frente a lo mismo que ve en una corrida real.

\begin{table}[htbp]
\centering
\caption{El cierre con verdad conocida sobre \EjCorrida, con las dos formas
de pesar el $\chi^2$. La cobertura cuenta en cuántas realizaciones la $D$
verdadera quedó dentro de la barra reportada; con barras correctas se espera
el 68\,\%.}
\label{tab:cierre}
\footnotesize
\begin{tabular}{lcccccc}
\toprule
\textbf{barras} & \textbf{sesgo $D$ [\%]} & \textbf{disp. [\%]} & \textbf{error rep. [\%]} & \textbf{cobertura} & \textbf{sesgo $\sigma$ [\%]} & \textbf{$\chi^2$/ndf} \\
\midrule
dato (Neyman) & $-0.04 \pm 0.03$ & 0.17 & 0.15 & 18 / 30 & $-0.4$ & 1.05 \\
modelo (Pearson) & $+0.00 \pm 0.03$ & 0.17 & 0.15 & 21 / 30 & $+0.4$ & 1.05 \\
\bottomrule
\end{tabular}
\par\smallskip
\begin{minipage}{0.95\linewidth}\footnotesize Sesgo: media de $D_{\text{rec}}/D_{\text{real}} - 1$ sobre las realizaciones, con el error de esa media. Disp.: desviación estándar de las $D$ recuperadas. Error rep.: la barra que la cadena reportó, en promedio. Si la barra es correcta, coincide con la dispersión y la cobertura queda cerca de 68\,\%.\end{minipage}

\end{table}

La tabla~\ref{tab:cierre} es el resultado sobre \CierreN{} realizaciones de
\EjCorrida, con una mediana de \CierreCuentasBin{} cuentas por bin en la
ventana. Tres lecturas:
\begin{itemize}
\item \textbf{La cadena no sesga $D$} al nivel que esta campaña puede
      resolver: pesando con el dato el sesgo es
      $\CierreSesgoNeyman \pm \CierreErrSesgoNeyman$\,\%, y pesando con el
      modelo, $\CierreSesgoPearson \pm \CierreErrSesgoPearson$\,\%. La
      diferencia entre los dos es el sesgo de Neyman a esta estadística, y
      está muy por debajo de los sistemáticos de la
      sección~\ref{sec:limites}.
\item \textbf{La barra reportada es del orden de la dispersión real, y algo
      corta}: \CierreErrRepNeyman\,\% reportado contra \CierreDispNeyman\,\%
      de dispersión, con cobertura de \CierreCobNeyman{} sobre \CierreN.
      La barra sale de la curvatura del $\chi^2$ \emph{con la ventana fija};
      la cadena, además, elige la ventana de nuevo en cada realización, y
      ese paso agrega una dispersión que la curvatura no puede ver (el panel
      (b) de la figura~\ref{fig:teoria-estadistica} superpone las dos
      elipses). Sobre datos sintéticos el $\chi^2/\text{ndf}$ queda cerca de
      uno, así que ahí el escalado no interviene; en las corridas reales sí,
      y es lo que agranda la barra de la tabla~\ref{tab:calib}.
\item \textbf{La semilla del NCC queda a \CierreSemillaNeyman\,\% de la
      verdad} en promedio, o sea que la forma sola, sin pesos, ya acierta
      dentro del uno por ciento; y $\sigma$ se recupera con un sesgo de
      \CierreSesgoSigmaNeyman\,\% (Neyman) y \CierreSesgoSigmaPearson\,\%
      (Pearson), del orden de su propia barra.
\end{itemize}

\begin{figure}[htbp]
\centering
\includegraphics[width=\linewidth]{teoria_estadistica.pdf}
\caption{De dónde sale la barra de error, sobre \EjCorrida. (a) El perfil de
$\Delta\chi^2$ en $D$ (puntos; en cada uno se reoptimizó $\sigma$) contra la
parábola que supone la matriz de covarianza (línea gris), en unidades de la
barra de la curvatura; las líneas horizontales son $\Delta\chi^2 = 1$ y
$\Delta\chi^2 = \chi^2/\text{ndf}$, los dos cortes de los que sale la barra
estadística y la barra escalada. (b) El plano $(D, \sigma)$: la elipse de
$\Delta\chi^2 = 1$ del ajuste real (gris, lo que la barra reportada supone),
la elipse de la dispersión empírica del cierre (trazos, lo que la cadena
hace), la verdad del cierre (estrella) y lo que
la cadena recuperó en cada realización, pesando con el dato (círculos) y con
el modelo (cuadrados). (c) La $D$ recuperada en cada realización, con su
barra, respecto de la verdad; las líneas de trazos son el sesgo medio de cada
forma de pesar.}
\label{fig:teoria-estadistica}
\end{figure}

\begin{aviso}
El cierre prueba que la cadena devuelve lo que se le mete \emph{cuando el
modelo es correcto}: el espectro sintético se fabrica con la misma plantilla
con que después se ajusta. Lo que no puede probar es que la plantilla sea la
correcta para el detector real. Eso lo dicen el $\chi^2/\text{ndf}$ de las
corridas reales y el test de dos fuentes, y es por donde entran los
sistemáticos.
\end{aviso}

% ===========================================================================
\section{La simulación de Geant4}
\label{sec:geant4}

El modelo del espectro no es una fórmula: es el histograma de una simulación
Monte Carlo de la geometría del detector con la fuente correspondiente, hecha
con Geant4~\cite{agostinelli2003,allison2016} (versión \SimVersion). Esta
sección dice qué se simuló y qué no, qué se guarda por evento, cómo eso se
convierte en la plantilla que el ajuste usa, y cómo se vuelve a correr. Lo que
importa para el lector que hereda esto: la plantilla \emph{es} la fuente, y
lo que no está en la geometría simulada tampoco está en el modelo.

\subsection{Qué se simula}

\paragraph{Geometría.} Un cilindro de plástico centellador de \SimRadio{} cm de
radio y \SimLargo{} cm de largo (2 por 2 pulgadas), con el eje en $z$. Delante,
hacia la fuente, una tapa de aluminio de \SimTapa{} mm; alrededor, una pared
lateral de aluminio de \SimPared{} mm; detrás, un cilindro de vidrio
borosilicato (\SimVentanaPmt) de \SimPmtSemilargo{} cm de semilargo que hace
las veces de ventana del fotomultiplicador. Todo dentro de un mundo de aire.
La tapa y la pared atenúan y dispersan el haz que entra; el vidrio de atrás es
el principal contribuyente a la retrodispersión desde detrás del detector.
Cada uno de los tres se puede apagar por macro (\texttt{/det/housing},
\texttt{/det/pmt}).

\paragraph{Material.} El centellador es el material NIST
\SimMaterial{} (poliviniltolueno, \SimDensidad{} g/cm$^3$, H
\XcomFraccionH{} y C \XcomFraccionC{} en peso), la base común de la familia
BC-400 (BC-404, BC-408, EJ-200). Los fluors que distinguen un modelo de otro
cambian la luz y su tiempo de decaimiento, no el depósito de energía: la
sección eficaz Compton depende del número de electrones por gramo y ese es el
mismo. Por eso el depósito se simula con el PVT desnudo y la luz sale del
quenching, que es lo que sigue.

\paragraph{Luz y quenching.} La luz producida por unidad de camino no es
proporcional a la energía perdida: a lo largo de una traza muy ionizante las
moléculas excitadas se desexcitan sin emitir, y la luz satura. Es la ley de
Birks~\cite{birks1951},
\[
   \frac{dL}{dx} \;=\; \frac{S\,\dfrac{dE}{dx}}{1 + k_B\,\dfrac{dE}{dx}},
\]
con $S$ la eficiencia de centelleo y $k_B$ la constante de Birks, acá
\SimBirks{} mm/MeV, el valor de literatura para la familia BC-400. Geant4 no
la trae para este material (su tabla interna cubre poliestireno, BGO, argón
líquido y PbWO$_4$, y deja el PVT explícitamente comentado): sin fijarla a
mano queda $k_B = 0$ y la luz sale exactamente proporcional al depósito, con
lo que el efecto que se quiere estudiar desaparece sin ningún aviso. La
simulación la fija en \texttt{DetectorConstruction::ApplyBirks()} y la
aplica \emph{paso a paso}, en cada step de cada partícula cargada, nunca
sobre el depósito total del evento: el quenching es no lineal en $dE/dx$, y
en un evento con dos o tres electrones de baja energía, quenchar la suma da
un resultado distinto de sumar los quench individuales. El detalle de qué
significa esto para las dos escalas de energía está en la
figura~\ref{fig:simulacion}(c) y en el test de dos fuentes
(sección~\ref{sec:testA}).

\paragraph{Fuente.} Gammas monoenergéticos desde un punto sobre el eje, a
\SimDistancia{} cm de la tapa. No se simula el decaimiento del núcleo ni los
electrones $\beta$: se dispara una línea por evento, sorteada con los pesos de
la fuente (para el Co-60, las dos líneas con igual peso; en los eventos que
depositan quedan \SimRazonLineasCo{} de la línea de \EgammaCoA{} keV por cada
uno de la de \EgammaCoB, porque la sección eficaz baja con la energía). El
haz de los archivos que usa este informe es \SimHaz, o sea axial: todos los
fotones entran por el centro de la tapa. Existe la opción de un cono que
cubra la cara o de emisión isótropa (\texttt{/gun/beam}), y también la de
emitir las dos líneas del Co-60 en el mismo evento para reproducir el
apilamiento por coincidencia (\texttt{/gun/cascade}); ninguna de las dos se
usó.

\paragraph{Física.} La lista electromagnética \SimFisica, la más precisa de
las estándar (modelos de baja energía con perfil Doppler para el Compton), y
\texttt{G4DecayPhysics}. El corte de producción de secundarios es de
\SimCorte{} mm, muy por debajo del alcance de un electrón de retroceso
de cientos de keV en el plástico (milímetros a 1 MeV): un secundario que no
llegaría a recorrer esa distancia no se sigue como partícula nueva, su
energía se deposita en el lugar.

\begin{aviso}
\textbf{Hay dos simulaciones y sólo una es esta.} En la carpeta del proyecto
mayor hay otra (\texttt{ComptonSimulation/}, Windows), anterior: sin carcasa
ni vidrio, con iones de Cs-137 decayendo y un solo histograma de salida en
CSV. No produjo los archivos que este informe usa. La fuente de la que sí los
produjo está copiada en \texttt{simulacion/} (el original vive en WSL) y es
de donde este informe lee la geometría y las constantes; \texttt{simulacion/LEEME.md}
dice qué hay y cómo se compila.
\end{aviso}

\subsection{Qué se guarda por evento}

Por cada evento con depósito en el centellador queda una fila en el árbol
\texttt{ComptonSpectrum}, con ocho columnas: el número de evento; la energía
depositada (\texttt{Edep\_keV}); la energía del gamma primario
(\texttt{E0\_keV}, que es lo que permite verificar de qué fuente es el
archivo); cuántas veces dispersó el fotón por Compton dentro del centellador
(\texttt{NCompton}) y cuántas absorbió por fotoeléctrico (\texttt{NPhoto});
si el depósito fue casi total (\texttt{FullAbs}); la luz emitida con el
quenching ya aplicado (\texttt{Lvis\_keVee}); y el cociente de las dos
últimas (\texttt{QF}). Los eventos sin depósito, que son la mayoría, no se
guardan. La tabla~\ref{tab:simulacion} resume los dos archivos.

\begin{table}[htbp]
\centering
\caption{Los dos archivos de simulación de la campaña, leídos de los propios
\texttt{.root}. El número de primarios lo escribe el programa en el nombre
del archivo a partir del \texttt{beamOn}. La fracción que deposita se compara
con la de la sección~\ref{sec:interaccion} (\ProbInteraccionCs\,\% y
\ProbInteraccionCoB\,\% por la sección eficaz total en \LadoDetector{} cm de
plástico): no tienen por qué coincidir exactamente, porque la tapa de Al
atenúa el haz y a la vez la carcasa entera dispersa fotones hacia adentro.}
\label{tab:simulacion}
{\small\begin{tabular}{lrr}
\toprule
\textbf{} & \textbf{Co-60} & \textbf{Cs-137} \\
\midrule
archivo & \texttt{Co60\_3M} & \texttt{Cs137\_10M} \\
gammas primarios & 3 M & 10 M \\
eventos con $E_{dep} > 0$ & 852079 & 3531670 \\
fracción que deposita [\%] & 28.4 & 35.3 \\
\midrule
1 scattering Compton [\%] & 75.1 & 73.0 \\
2 scatterings [\%] & 17.1 & 19.3 \\
3 o más [\%] & 5.7 & 7.0 \\
ningún Compton en el centellador [\%] & 2.1 & 0.7 \\
con absorción fotoeléctrica [\%] & 2.9 & 1.3 \\
absorción total ($E_{dep} \approx E_0$) & 87 & 724 \\
\midrule
$\langle L/E_{dep}\rangle$, 50--100 keV & 0.863 & 0.859 \\
$\langle L/E_{dep}\rangle$, $E_{dep} > 700$ keV & 0.954 & -- \\
$\langle L/E_{dep}\rangle$ en la zona del ajuste & 0.954 & 0.935 \\
bordes en luz [keVee] & 921, 1069 & 448 \\
\bottomrule
\end{tabular}
\par\smallskip
\begin{minipage}{0.95\linewidth}\footnotesize Las fracciones por número de scatterings son sobre los eventos con depósito. «Ningún Compton en el centellador» son casi todos eventos cuyo depósito lo hace un electrón Compton nacido en la tapa de aluminio que entra al plástico (mediana de 275 keV en el Cs-137); el resto son fotoeléctricos directos. La zona del ajuste es $E_{dep} > 0{,}7\,E_{\text{ref}}$, que es donde se promedia el quenching que convierte $D_{\text{luz}}$ en $D_{\text{dep}}$.\end{minipage}
}
\end{table}

Tres cosas que la tabla dice y que conviene tener presentes al leer un
ajuste:

\begin{itemize}
\item \textbf{Un cuarto de los eventos es scattering múltiple}
      (\SimFracMultipleCo\,\% en el Co-60, \SimFracMultipleCs\,\% en el
      Cs-137, con cadenas de hasta \SimNcMaxCs{} dispersiones). Esos eventos
      pueden depositar más que el borde, hasta la energía completa del gamma
      (el máximo de \texttt{Edep} es \SimEdepMaxCs{} keV en el Cs-137, que es
      la línea entera), y son los que llenan la zona por encima de $E_{\max}$
      en la figura~\ref{fig:simulacion}(b). Klein-Nishina describe sólo la
      componente de un scattering; el resto es lo que la plantilla agrega y
      la fórmula no tiene.
\item \textbf{El fotoeléctrico existe, y el fotopico también, pero no se
      ve}: lo sufre el \SimFracFotoCs\,\% de los eventos, casi siempre un
      fotón que ya dispersó por Compton al menos una vez y llegó con una
      fracción de su energía, donde la sección eficaz fotoeléctrica es mucho
      mayor (sección~\ref{sec:interaccion}). Son los eventos que absorben la
      energía completa: el fotopico del plástico son unos pocos cientos de
      eventos en millones, invisible al lado del continuo.
\item \textbf{El quenching no es un factor constante.} $\langle L/E_{dep}\rangle$
      vale \SimQfBajoCo{} entre 50 y 100 keV y \SimQfAltoCo{} por encima de
      700; en la zona que domina el ajuste vale \SimQfAjusteCo{} para el Co-60
      y \SimQfAjusteCs{} para el Cs-137, y esa diferencia es la que aparece en
      el test de dos fuentes (sección~\ref{sec:testA}) como el cociente
      esperado \CocEsperado. Los bordes, en luz, caen en \SimBordeLuzCs{}
      keVee (Cs-137) y \SimBordeLuzCo{} y \SimBordeLuzCoB{} keVee (Co-60), no
      en \BordeCs, \BordeCoA{} y \BordeCoB.
\end{itemize}

\begin{figure}[htbp]
\centering
\includegraphics[width=\linewidth]{simulacion_geant4.pdf}
\caption{Lo que hay en la simulación y lo que le falta. (a) Co-60: energía
depositada (gris) y luz (azul) de los mismos eventos; los bordes teóricos
(trazos) y los mismos bordes en luz (puntos), corridos por el quenching.
(b) Cs-137, en luz, separado por número de dispersiones Compton del fotón:
la componente de una sola dispersión es lo que Klein-Nishina describe y
termina en el borde; las de dos y más lo rellenan por encima. (c) El
quenching de Birks evento a evento: mediana de $L/E_{dep}$ en bines de
$E_{dep}$, con la banda del 16 al 84\,\%, para los dos archivos; las
líneas verticales son los tres bordes. (d) La razón dato/modelo del ajuste
de \EjCorrida{} en todo el rango: dentro de la ventana de ajuste (gris) es
uno; por debajo crece hasta \SimExcesoMax, que es el exceso de bajo canal
que la geometría simulada no contiene; y en el umbral del trigger (línea
punteada) cae a cero, porque eso tampoco está simulado.}
\label{fig:simulacion}
\end{figure}

\subsection{Del árbol a la plantilla}

La plantilla es el histograma de \texttt{\Rama} de todos los eventos con
depósito, en \SimBinsPlantilla{} bines uniformes desde cero hasta el máximo de la rama
(\texttt{analisis/simulacion.py}). Eso es todo lo que hay que decidir, y
resulta que casi no importa:

\begin{itemize}
\item \textbf{A canal.} El mapa es $ch = D\,E$, así que la plantilla en canal
      es la misma curva leída en $ch/D$: se interpola linealmente el
      histograma en $ch/D$ para cada bin del espectro medido
      (\texttt{ncc.plantilla\_en\_canal}). No se rebinnea; por eso no aparece
      el peine de \emph{aliasing} que da rebinnear un histograma sobre una
      grilla que no le es conmensurable.
\item \textbf{La resolución.} La curva en canal se convoluciona con una
      gaussiana de ancho $\sigma_{ch} = \sigma_{\text{rel}}\,D\,E_{\text{ref}}$,
      normalizada a suma uno para que la convolución conserve las cuentas
      (sección~\ref{sec:lineshape}). La simulación no lleva resolución
      adentro a propósito: así se puede volver a borronear con otro $\sigma$
      sin volver a simular, y $\sigma$ queda como parámetro del ajuste.
\item \textbf{El binning es irrelevante, medido.} Reajustando \EjCorrida{}
      con la misma cadena y la plantilla en \BinsMin{} a \BinsMax{} bines
      (de \BinsAnchoMax{} a \BinsAnchoMin{} keV por bin), $D$ se mueve como
      máximo \BinsDesvMax\,\%, muy por debajo de su barra estadística. La
      razón es que el kernel de la resolución mide $\sigma_E \approx
      \BinsSigmaE$ keV, unas \BinsSigmaSobreBin{} veces el bin más fino:
      cualquier estructura del tamaño de un bin queda borrada por la
      convolución antes de compararse con el dato. Lo que sí depende un poco
      del binning es la \emph{barra} de $D$ por curvatura, porque la plantilla
      interpolada es lineal a trozos y el paso de derivación la ve; es el
      mismo efecto de la sección~\ref{sec:estadistica} y otra razón para
      preferir el perfil de $\Delta\chi^2$.
\end{itemize}

\paragraph{La plantilla es la fuente, y hay que verificarlo.} Usar el archivo
de Co-60 contra un espectro de Cs-137 no da error en ningún lado: da una $D$
perfectamente formateada y sin sentido, porque la plantilla tiene sus bordes
en \BordeCoA{} y \BordeCoB{} keV y los datos uno solo en \BordeCs. Por eso la
cadena lee la columna \texttt{E0\_keV} del archivo
(\texttt{simulacion.lineas\_primarias}), compara las líneas con las de la
fuente declarada en el catálogo y se niega a seguir si no coinciden. Es la
misma clase de falla silenciosa que un canal del digitalizador cableado a
mano (sección~\ref{sec:cadena}): lo que no se verifica, se propaga con buena
cara.

\subsection{Qué no está simulado}

Es la lista de por dónde entran los sistemáticos, y la figura~\ref{fig:simulacion}(d)
muestra los dos primeros directamente:

\begin{itemize}
\item \textbf{El entorno.} No hay portafuente, mesa, soporte ni paredes:
      sólo el detector con su carcasa, en aire. Los fotones que dispersan en
      lo que rodea al detector y entran después con menos energía son el
      candidato principal para el \textbf{exceso de bajo canal}: el dato
      supera al modelo por debajo de la ventana de ajuste, y crece hacia
      abajo. No es el umbral (eso sería un déficit) ni el fondo (está
      restado en las corridas que lo tienen). Sin medirlo, la única defensa
      es dejarlo fuera de la ventana (sección~\ref{sec:decisiones}).
\item \textbf{El umbral del trigger.} La simulación registra todo depósito
      mayor que cero; el discriminador de la placa no. Por eso el espectro
      medido sube desde cero y el simulado no, y por eso ninguna ventana
      puede empezar antes del fin de esa subida.
\item \textbf{El fotomultiplicador y la electrónica.} El vidrio de atrás es
      geometría, no fotodetección: no hay fotones ópticos, ni fotocátodo, ni
      ganancia, ni pulso. Toda la cadena que va de la luz al área en \mVns{}
      está resumida en dos números del ajuste, $D$ y $\sigma$.
\item \textbf{La resolución}, que como se dijo se pone en el análisis.
\item \textbf{El fondo, los cósmicos y el apilamiento.} Un evento simulado
      es un gamma; en la corrida real hay muones, radiactividad ambiente y
      dos pulsos que se superponen en la compuerta. Es lo que hay por encima
      del borde en el espectro medido y que la plantilla no tiene, y otra
      razón para que la ventana termine donde termina.
\item \textbf{Si la geometría es la del detector real.} El tamaño del
      plástico, el grosor de la carcasa y la distancia a la fuente son los
      valores por defecto de la simulación, no medidas de los cuatro
      detectores de la campaña (sección~\ref{sec:limites}). La cantidad de
      scattering múltiple depende del tamaño, y con ella la forma por
      encima del borde. En la implementación anterior del análisis se dejó
      libre la fracción de un solo scattering y el ajuste la puso a menos de
      una desviación estándar de la de la simulación, que es el único control
      de este punto que se hizo; con el detector medido, hay que repetirlo.
\end{itemize}

\subsection{Cómo volver a correrla}

La fuente está en \texttt{simulacion/} y se compila en WSL contra Geant4
\SimVersion{} (\texttt{cmake} y \texttt{make}, ver \texttt{LEEME.md}). Cada
fuente tiene su macro (\texttt{co60.mac}, \texttt{cs137.mac},
\texttt{na22.mac}): elige el preset, inicializa y dispara. El programa
escribe el archivo con el nombre \texttt{ComptonSpectrum\_<fuente>\_<N>.root},
donde \texttt{N} es el número de primarios; los de \texttt{datos/} se corrieron
con \SimPrimariosCo{} y \SimPrimariosCs{} millones. Para otra fuente basta
agregar sus líneas y pesos a \texttt{SetPreset} (o pasar una energía suelta
con \texttt{/gun/energy}), copiar el archivo a \texttt{datos/} y declararlo
en el catálogo (\texttt{analisis/rutas.py}); la verificación de líneas hace el
resto. Para otro detector, los comandos \texttt{/det/radius},
\texttt{/det/halfZ}, \texttt{/det/housingWall} y \texttt{/det/housingFront}
van antes de \texttt{/run/initialize}, y conviene que salgan de una medición
y no de un valor por defecto.

% ===========================================================================
\section{La cadena de análisis, paso a paso}
\label{sec:cadena}

Esta sección es el recorrido completo, del archivo crudo al número. Cada paso
dice qué hace, por qué está ahí y qué se rompería si no estuviera.

\subsection{Del archivo del DRS4 al área de cada pulso}

El detalle, con los pulsos reales y los números de cada corrida, está en la
sección~\ref{sec:pulsos}; acá va el resumen para que la cadena se lea de
corrido. Cada corrida es un archivo con decenas de miles de eventos, cada uno
con la forma de onda digitalizada. El formato lo decide el contenido, no el
nombre, y el canal también se detecta leyendo el archivo. De cada evento se
calcula el área así:

\begin{enumerate}
\item \textbf{Se localiza el onset}, el instante en que arranca el pulso,
      definido como el cruce del \FracOnset\,\% de la amplitud retrocediendo
      desde el pico. Anclar todo al onset de \emph{cada} pulso, y no a un
      tiempo fijo del registro, es lo que hace inmune la medición al jitter y
      al \emph{walk} del trigger.
\item \textbf{Se estima la línea base} como la mediana de las muestras
      anteriores al onset, con un margen de \GuardaBase{} ns. Si la ventana de
      línea base se comiera la subida del propio pulso, el error resultante se
      multiplicaría al integrar y cada evento perdería una fracción distinta
      de su carga.
\item \textbf{Se integra} entre \emph{pre} ns antes del onset y \PostFijo{} ns
      después, con la línea base ya restada. La integral usa el paso real
      entre muestras: el DRS4 no muestrea uniformemente, el intervalo entre
      celdas varía, y multiplicar por un $\Delta t$ promedio introduce un
      error que depende de dónde cayó el pulso.
\item \textbf{Se marcan las banderas de calidad}: saturación del
      digitalizador, pulso truncado porque la compuerta se sale del registro,
      falta de pre-pulso para la línea base, ausencia de pulso. Los eventos
      con problemas que afectan el área se descartan, y la fracción descartada
      queda registrada (tabla~\ref{tab:pulsos}).
\end{enumerate}

\subsection{Del área al espectro neto}

Las áreas de todos los eventos se histograman
(sección~\ref{sec:areas-espectro}). Tres cosas que parecen detalles y no lo
son:

\begin{itemize}
\item \textbf{Binning común} para muestra y fondo. Sin el mismo rango y el
      mismo número de bines la resta no significa nada, porque cada bin
      abarcaría un intervalo distinto de \mVns.
\item \textbf{Normalización por tiempo vivo, no por tiempo de pared.} El
      programa de adquisición no guarda todos los eventos que la placa dispara;
      el scaler de la placa da la tasa real de trigger, y de ahí sale cuánto
      tiempo estuvo viva la corrida. Usar el tiempo de reloj en lugar del
      tiempo vivo puede sobre-restar el fondo por un factor grande.
\item \textbf{Sin recortar los negativos.} Después de restar el fondo hay bines
      con valor negativo por fluctuación estadística. Recortarlos a cero sesga
      el ajuste hacia arriba; se conservan.
\end{itemize}

El resultado es el espectro neto con su barra de error por bin, propagada de
la estadística de Poisson de la muestra y del fondo. Ese error es el que pesa
el ajuste, y sin él no habría $\chi^2$, ni bondad de ajuste, ni barra de error
sobre $D$.

\subsection{La plantilla de Geant4}

El modelo del espectro es el histograma de energía depositada de la
simulación de la sección~\ref{sec:geant4}: contiene el continuo Compton, los
dos bordes del Co-60 con su intensidad relativa correcta, el scattering
múltiple y el quenching, todo junto y sin parámetros libres. Se lleva a canal
interpolándolo en $ch/D$ y se convoluciona con la resolución; el binning del
histograma no cambia $D$ (medido: sección~\ref{sec:geant4}).

La cadena verifica que la plantilla sea de la fuente declarada leyendo los
gammas primarios del archivo de simulación, y se niega a seguir si no
coinciden. Una plantilla de la fuente equivocada no da error en ningún lado:
da un número equivocado bien formateado.

\subsection{La ventana y la semilla: el NCC en la práctica}

El formalismo, y el detalle de qué entrega y cómo lo usa el recocido, está en
la sección~\ref{sec:ncc}. Acá, lo que se corre: el barrido decide la ventana
de ajuste de cada corrida y entrega la semilla $D_0 = B/E_{\text{ref}}$, sin
intervención; de la semilla salen las cotas ($\CotaDMin$ a $\CotaDMax\,D_0$) y
el paso del recocido, que arranca \NArranques{} veces desde puntos al azar
dentro de esas cotas. Sobre las \NCorridas{} corridas el intervalo ganador
correlaciona entre $\rho = \RhoMin$ y $\rho = \RhoMax$, y entre la semilla y la
$D$ con la que termina el ajuste hay una diferencia mediana de
\SemillaMediana\,\% (máxima \SemillaMax\,\%, sobre \NSemillas{} corridas). O
sea que la forma sola, sin pesos, ya deja a $D$ a menos del uno por ciento; lo
que el ajuste agrega es la barra de error, la resolución y la evidencia de que
ese mínimo es el global.


\subsection{El ajuste}
\label{sec:ajuste}

El modelo es la plantilla llevada al eje medido dividiendo por $D$ y
convolucionada con una gaussiana de ancho $\sigma$:
\[
   S(ch) \;=\; A \cdot \big[\,P(ch/D) \otimes G(\sigma)\,\big].
\]
Nada más: dos parámetros no lineales, $\{D, \sigma_{\text{rel}}\}$, y la
amplitud $A$.

\paragraph{La amplitud no se busca.} $A$ entra linealmente, así que fijados $D$
y $\sigma$ la amplitud que minimiza el $\chi^2$ se despeja en forma cerrada.
Eso se llama proyección y compra dos cosas: el optimizador busca en dos
dimensiones en vez de tres, y desaparece el problema de escala entre
parámetros (una amplitud de miles de cuentas y una ganancia de orden 1 no
admiten un paso común).

\paragraph{Por qué recocido simulado y no un optimizador corriente.} Un
optimizador local sólo sabe bajar: desde donde arranca, camina hasta el fondo
del valle en el que cayó, y no tiene manera de saber si hay otro más hondo. El
recocido acepta a veces moverse a un punto peor, con probabilidad
$e^{-\Delta E/T}$, y esa es la única forma de cruzar una loma. La temperatura
$T$ arranca alta (la cadena pasea y explora) y se enfría geométricamente hasta
comportarse como un optimizador local. Acá se corren \NArranques{} recocidos
desde puntos iniciales dispersos dentro de las cotas, no desde la semilla: si
todos arrancaran del mismo punto, que todos terminen igual no probaría nada.
El detalle completo, con el código comentado línea por línea, está en
\texttt{docs/recocido\_simulado\_codigo.pdf}.

\paragraph{Qué tan estable es el resultado.} La dispersión entre los
\NArranques{} arranques es de \DispersionMediana\,\% de $D$ en la corrida
mediana y no pasa de \DispersionMax\,\% en la peor: el optimizador no aporta
incerteza apreciable. Lo que limita la precisión es la estadística y, sobre
todo, los sistemáticos de la sección~\ref{sec:limites}.

\subsection{De dónde sale la barra de error}

Del ancho del mínimo (el fundamento, con el perfil del $\chi^2$ y el cierre
con verdad conocida, está en la sección~\ref{sec:estadistica}; acá va lo que
se corre). Si el valle del $\chi^2$ es angosto, moverse poco empeora
mucho el ajuste y el parámetro está bien determinado; si es ancho y plano, no.
Numéricamente se evalúa la curvatura (el hessiano) alrededor del óptimo y se la
invierte para obtener la matriz de covarianza; la raíz de la diagonal es la
barra de error, que corresponde al contorno $\chi^2_{\min} + 1$.

Dos advertencias sobre esa barra:

\begin{itemize}
\item \textbf{Está escalada por $\sqrt{\chi^2/\text{ndf}}$.} Cuando el ajuste
      no describe perfectamente los datos, el criterio $\Delta\chi^2 = 1$
      subestima la incerteza. Escalar es el parche estándar y hay que
      reportarlo como tal, no como una medición del error.
\item \textbf{$D$ y $\sigma$ están correlacionados}, entre \CorrMin{} y
      \CorrMax{} según la corrida. Mover uno se compensa parcialmente con el
      otro, así que citar sus errores por separado sobreestima la precisión de
      cualquier cosa que dependa de los dos.
\end{itemize}

\subsection{Cuatro decisiones que valen para todas las corridas}
\label{sec:decisiones}

\paragraph{Plantilla de la fuente que corresponde.} El Co-60 se ajusta contra
la simulación de Co-60 y el Cs-137 contra la de Cs-137, verificado leyendo el
archivo de simulación.

\paragraph{Compuerta común de \PostFijo{} ns.} La ventana de integración se
fija igual para todas las corridas, contada desde el onset de cada pulso
(cómo se mide y qué recoge: sección~\ref{sec:compuerta}). Que
sea la \emph{misma} hace falta porque la ventana que cada archivo mediría por
su cuenta depende de su propio pulso, y dos corridas que integran distinto
recogen fracciones distintas de la carga, lo que entra derecho en $D$. Que sea
\emph{corta} es una decisión de análisis: pasado el pulso, seguir integrando no
agrega señal, agrega línea base y su ruido.

El precio hay que decirlo con todas las letras, porque es grande: con
\PostFijo{} ns la compuerta recoge entre el \CargaMin\,\% y el \CargaMax\,\% de
la carga del pulso registrado (columna \emph{Carga} de la
tabla~\ref{tab:calib}), y la resolución que sale del ajuste queda entre
\FwhmMin\,\% y \FwhmMax\,\% FWHM. Dos consecuencias:

\begin{itemize}
\item Estas $D$ son un observable común y comparable, \textbf{no} la ganancia
      de pulso entero. Para comparar corridas entre sí da igual, mientras la
      ventana sea la misma; para citar una $D$ como ganancia absoluta del
      detector, no.
\item Integrar una fracción del pulso empeora la relación señal-ruido por
      evento, y eso se ve directamente en la resolución. Es el costo de la
      ventana corta y se paga en todas las corridas por igual.
\end{itemize}

\paragraph{Modelo sin exceso de bajo canal.} El ajuste es sólo la plantilla
convolucionada. Una versión anterior de esta cadena agregaba una exponencial
para describir el exceso de cuentas que aparece muy por debajo del borde, con
dos parámetros más ($\lambda$ y su amplitud); se sacó, y con eso el recocido
busca en dos dimensiones en lugar de tres.

Lo que se pierde es la única defensa contra ese exceso, que sigue estando en
los datos aunque el modelo ya no lo tenga. Con la exponencial adentro el ajuste
podía arrancar bien abajo del borde; sin ella hay que dejar el exceso
\emph{afuera} de la ventana, y dónde empieza la ventana lo decide el parámetro
\texttt{inicio\_rel} (aquí, \InicioRel{} veces el canal del borde). O sea que
ese número pasa a ser un sistemático elegido a mano: bajarlo mete zona sucia en
el ajuste y sube el $\chi^2$; subirlo achica la palanca del flanco.

\paragraph{Sin fondo prestado.} Sólo dos corridas tienen fondo propio (el
Co-60 de junio y el Cs-137 del detector A a 1600 V). Las demás se calibran con
el espectro crudo. Prestarles el fondo de otra sesión sería peor que no restar
nada: el fondo depende del punto de operación, así que introduciría una forma
ajena en el espectro. La tabla~\ref{tab:calib} dice corrida por corrida cuál
lleva fondo restado.

% ===========================================================================
\section{Las \NCorridas{} calibraciones}
\label{sec:calibraciones}

\subsection{La tabla}

La tabla~\ref{tab:calib} es el resultado crudo de la campaña: $\chi^2$/ndf
entre \ChiMin{} y \ChiMax{}, resolución entre \FwhmMin\,\% y \FwhmMax\,\% FWHM
en el borde de referencia, y error estadístico mediano de $D$ del
\ErrRelMediano\,\%. Los dos primeros números hay que leerlos junto con la
compuerta corta y el modelo sin exceso de la sección~\ref{sec:decisiones}: son
las dos decisiones que los fijan, y las dos se pueden cambiar en la interfaz
sin tocar código.

\begin{center}
\captionof{table}{Las \NCorridas{} calibraciones, en escala de energía
depositada y con compuerta común de \PostFijo{} ns.}
\label{tab:calib}
% La tabla tiene diez columnas: sin apretar `tabcolsep` no entra a lo ancho.
{\footnotesize\setlength{\tabcolsep}{4pt}\begin{longtable}{llllrrrrrc}
\toprule
\textbf{Det.} & \textbf{HV} & \textbf{Fuente} & \textbf{Medida} & \textbf{$D$} & \textbf{$1/D$} & \textbf{$\chi^2$/ndf} & \textbf{FWHM} & \textbf{Carga} & \textbf{Fondo} \\
\midrule
\endfirsthead
\toprule
\textbf{Det.} & \textbf{HV} & \textbf{Fuente} & \textbf{Medida} & \textbf{$D$} & \textbf{$1/D$} & \textbf{$\chi^2$/ndf} & \textbf{FWHM} & \textbf{Carga} & \textbf{Fondo} \\
\midrule
\endhead
A & 1550 & Co-60 & 2026-06-23 & $0.9558 \pm 0.0022$ & 1.0 & 1.80 & 31.3 & 82 & si \\
A & 1600 & Co-60 & 2026-09-01 & $1.2617 \pm 0.0037$ & 0.8 & 2.46 & 33.2 & 77 & no \\
A & 1600 & Cs-137 & 2026-08-26 & $1.3563 \pm 0.0022$ & 0.7 & 2.05 & 25.1 & 72 & si \\
A & 1650 & Co-60 & 2026-09-01 & $1.6958 \pm 0.0022$ & 0.6 & 2.20 & 29.7 & 77 & no \\
A & 1650 & Cs-137 & 2026-08-26 & $1.8710 \pm 0.0035$ & 0.5 & 1.92 & 26.7 & 77 & no \\
A & 1700 & Co-60 & 2026-09-01 & $2.2694 \pm 0.0044$ & 0.4 & 2.08 & 33.4 & 76 & no \\
A & 1700 & Cs-137 & 2026-08-26 & $2.5502 \pm 0.0046$ & 0.4 & 1.78 & 27.4 & 74 & no \\
\midrule
B & 1600 & Co-60 & 2026-09-22 & $1.0033 \pm 0.0017$ & 1.0 & 2.59 & 32.4 & 77 & no \\
B & 1600 & Cs-137 & 2026-09-01 & $1.0717 \pm 0.0021$ & 0.9 & 2.01 & 25.6 & 76 & no \\
B & 1650 & Co-60 & 2026-09-22 & $1.3472 \pm 0.0021$ & 0.7 & 2.35 & 31.1 & 78 & no \\
B & 1650 & Cs-137 & 2026-09-01 & $1.4364 \pm 0.0032$ & 0.7 & 2.15 & 24.6 & 77 & no \\
B & 1700 & Co-60 & 2026-09-22 & $1.7780 \pm 0.0033$ & 0.6 & 2.04 & 32.0 & 78 & no \\
B & 1700 & Cs-137 & 2026-09-01 & $1.9140 \pm 0.0047$ & 0.5 & 1.75 & 25.5 & 79 & no \\
\midrule
C & 1600 & Co-60 & 2026-09-02 & $1.3868 \pm 0.0021$ & 0.7 & 2.70 & 31.3 & 75 & no \\
C & 1600 & Cs-137 & 2026-09-01 & $1.2055 \pm 0.0024$ & 0.8 & 1.84 & 23.4 & 76 & no \\
C & 1650 & Co-60 & 2026-09-02 & $1.7735 \pm 0.0029$ & 0.6 & 2.20 & 30.9 & 72 & no \\
C & 1650 & Cs-137 & 2026-09-01 & $1.5989 \pm 0.0032$ & 0.6 & 2.06 & 24.3 & 77 & no \\
C & 1700 & Co-60 & 2026-09-02 & $2.3850 \pm 0.0065$ & 0.4 & 2.10 & 32.1 & 73 & no \\
C & 1700 & Cs-137 & 2026-09-01 & $2.1246 \pm 0.0049$ & 0.5 & 1.78 & 24.6 & 77 & no \\
\midrule
D & 1600 & Co-60 & 2026-09-08 & $0.4222 \pm 0.0006$ & 2.4 & 2.60 & 34.6 & 81 & no \\
D & 1600 & Cs-137 & 2026-09-08 & $0.4934 \pm 0.0011$ & 2.0 & 1.96 & 32.8 & 81 & no \\
D & 1600 & Cs-137 (descartada) & 2026-09-01 & $0.3599 \pm 0.0004$ & 2.8 & 1.18 & 30.7 & 73 & no \\
D & 1650 & Co-60 & 2026-09-08 & $0.5538 \pm 0.0009$ & 1.8 & 2.25 & 37.3 & 81 & no \\
D & 1650 & Cs-137 & 2026-09-08 & $0.6323 \pm 0.0010$ & 1.6 & 1.54 & 31.1 & 81 & no \\
D & 1650 & Cs-137 (descartada) & 2026-09-01 & $0.4668 \pm 0.0010$ & 2.1 & 1.75 & 32.9 & 75 & no \\
D & 1700 & Co-60 & 2026-09-08 & $0.7863 \pm 0.0015$ & 1.3 & 2.39 & 36.8 & 81 & no \\
D & 1700 & Cs-137 & 2026-09-08 & $0.7877 \pm 0.0019$ & 1.3 & 1.81 & 32.5 & 81 & no \\
D & 1700 & Cs-137 (descartada) & 2026-09-01 & $0.6171 \pm 0.0010$ & 1.6 & 1.53 & 33.1 & 74 & no \\
\bottomrule
\end{longtable}
\par\smallskip
\begin{minipage}{0.95\linewidth}\footnotesize Unidades: HV en V, $D$ en mV$\cdot$ns/keV, $1/D$ en keV/mV$\cdot$ns, FWHM y carga en \%. $1/D$ es lo que se usa para convertir: energia depositada por unidad de area. FWHM esta evaluada en el borde de referencia de cada fuente (477.3 keV el Cs-137, 963.4 keV el Co-60). «Carga» es la fraccion del pulso registrado que entra en la compuerta comun. «Fondo» dice si la corrida tiene fondo propio restado.\end{minipage}
}
\end{center}

\subsection{Columna por columna}

\begin{description}
\item[Det., HV, Fuente] El punto de operación. Una calibración vale para un
      detector \emph{a un voltaje}, no para un detector: la ganancia de un
      fotomultiplicador cambia con el alto voltaje como una potencia de
      exponente cercano a 10, así que 50 V de diferencia son un factor 1.3.
\item[Medida] La fecha del archivo. Sirve para saber qué corridas comparten
      sesión, que es la variable escondida de la sección~\ref{sec:testA}.
\item[$D$] La calibración, en \DkeV. El error es el estadístico del ajuste,
      escalado por $\sqrt{\chi^2/\text{ndf}}$.
\item[$1/D$] Lo mismo dado vuelta, en keV por \mVns, que es la forma en que se
      usa: multiplicar el área medida por este número da la energía.
\item[$\chi^2$/ndf] Bondad del ajuste. Un valor cercano a 1 significa que el
      modelo describe los datos dentro de las barras de error. Por encima de 2,
      como pasa en varias corridas de esta campaña, significa que hay
      estructura en los datos que el modelo no tiene, y el candidato conocido
      es el exceso de bajo canal.
\item[FWHM] El ancho a media altura de la resolución en el borde de
      referencia, en porcentaje de esa energía. Es
      $2.35\,\sigma_{\text{rel}}$.
\item[Carga] Qué fracción del pulso registrado entra en la compuerta de
      \PostFijo{} ns. Es el precio de la ventana corta, y explica por qué estas
      $D$ no son la ganancia de pulso entero.
\item[Fondo] Si la corrida tiene fondo propio restado o se calibró con el
      espectro crudo.
\end{description}

\subsection{Cómo mirar un ajuste}
\label{sec:mirarajuste}

Cada corrida deja seis páginas de diagnóstico
(\NFigurasDiag{} figuras en total, apéndice~\ref{ap:archivos}). La
figura~\ref{fig:ajuste} es la primera de ellas. Lo que hay que mirar, en orden:

\begin{enumerate}
\item \textbf{Que el modelo siga al dato en el flanco}, que es la zona que
      determina $D$. Un modelo que va sistemáticamente por encima o por debajo
      del plateau pero sigue bien el flanco todavía da un $D$ razonable; al
      revés, no.
\item \textbf{Dónde cae la línea del borde} respecto de la caída visible. Si la
      línea roja está claramente corrida del punto de inflexión, algo salió
      mal.
\item \textbf{El espectro calibrado} (panel derecho): con el eje ya en keV, el
      borde tiene que caer sobre su valor teórico. Para la fuente que se
      ajustó, eso pasa por construcción; el valor está ahí para comparar con la
      otra fuente del mismo detector.
\item \textbf{La zona de bajo canal}: el pico que sube a la izquierda es el
      exceso más el umbral del trigger, y no está modelado. Si la ventana de
      ajuste lo toca, el $\chi^2$ lo va a acusar.
\end{enumerate}

\begin{figure}[htbp]
\centering
\includegraphics[width=\linewidth]{ajuste_A.pdf}
\caption{Un ajuste típico (detector A a 1600 V). Izquierda: el espectro neto
con el modelo ajustado y los bordes Compton donde caen. Derecha: el mismo
espectro con el eje ya calibrado en energía. Los \NCorridas{} ajustes están en
\texttt{analisis/figuras/calibraciones/}, ordenados por detector, voltaje y
fuente.}
\label{fig:ajuste}
\end{figure}

% ===========================================================================
\section{Mismo detector al mismo voltaje: ¿coinciden las dos fuentes?}
\label{sec:testA}

Ésta es la única pregunta que valida la calibración. Con una sola fuente el
mapa por el origen tiene un solo parámetro, así que el punto medido cae exacto
sobre su propia recta y no hay nada que testear. Con dos fuentes en el
\emph{mismo} punto de operación sí lo hay: las dos tienen que dar la misma
$D$.

\subsection{Qué habría que esperar}

En escala de energía depositada, dos fuentes distintas \textbf{no} dan
exactamente la misma $D$ aunque la cadena sea perfectamente lineal. El área
mide luz, y la luz por keV depositado vale \QuenchCs{} en el borde del Cs-137
(\BordeCs{} keV) y \QuenchCo{} en el del Co-60 (\BordeCoA{} keV), según la
propia simulación. El cociente que corresponde esperar es entonces
\[
   \frac{D(\text{Cs-137})}{D(\text{Co-60})}
   \;=\; \frac{\QuenchCs}{\QuenchCo} \;=\; \CocEsperado,
\]
y no 1. La columna \emph{Desac.} de la tabla~\ref{tab:testA} ya tiene ese
factor descontado: es el desacuerdo que queda después de la física conocida.

\subsection{Lo que se mide}

\begin{table}[htbp]
\centering
\caption{Test de linealidad en cada punto de operación.}
\label{tab:testA}
{\small\setlength{\tabcolsep}{4.5pt}\begin{tabular}{lrrrrrrc}
\toprule
\textbf{Punto} & \textbf{$D$(Cs-137)} & \textbf{$D$(Co-60)} & \textbf{$D$ conjunta} & \textbf{Cs/Co} & \textbf{Desac. [\%]} & \textbf{$c_0$ [mV$\cdot$ns]} & \textbf{Sesiones} \\
\midrule
A @ 1550 V & \multicolumn{7}{l}{\textit{una sola fuente: el test no se puede correr}} \\
A @ 1600 V & 1.3563 & 1.2617 & 1.3312 & 1.075 & $+9.7$ & $+90 \pm 4$ & distintas \\
A @ 1650 V & 1.8710 & 1.6958 & 1.7445 & 1.103 & $+12.6$ & $+166 \pm 4$ & distintas \\
A @ 1700 V & 2.5502 & 2.2694 & 2.4037 & 1.124 & $+14.7$ & $+266 \pm 6$ & distintas \\
B @ 1600 V & 1.0717 & 1.0033 & 1.0306 & 1.068 & $+9.0$ & $+65 \pm 3$ & distintas \\
B @ 1650 V & 1.4364 & 1.3472 & 1.3738 & 1.066 & $+8.8$ & $+84 \pm 4$ & distintas \\
B @ 1700 V & 1.9140 & 1.7780 & 1.8234 & 1.076 & $+9.9$ & $+129 \pm 5$ & distintas \\
C @ 1600 V & 1.2055 & 1.3868 & 1.3076 & 0.869 & $-11.3$ & $-172 \pm 3$ & distintas \\
C @ 1650 V & 1.5989 & 1.7735 & 1.6969 & 0.902 & $-8.0$ & $-165 \pm 4$ & distintas \\
C @ 1700 V & 2.1246 & 2.3850 & 2.2185 & 0.891 & $-9.1$ & $-246 \pm 8$ & distintas \\
D @ 1600 V & 0.4934 & 0.4222 & 0.4392 & 1.169 & $+19.3$ & $+67 \pm 1$ & misma \\
D @ 1650 V & 0.6323 & 0.5538 & 0.5882 & 1.142 & $+16.5$ & $+74 \pm 1$ & misma \\
D @ 1700 V & 0.7877 & 0.7863 & 0.7868 & 1.002 & $+2.2$ & $+1 \pm 2$ & misma \\
\bottomrule
\end{tabular}
\par\smallskip
\begin{minipage}{0.95\linewidth}\footnotesize «Desac.» es el desacuerdo entre las dos fuentes una vez descontado el quenching de Birks: en escala de energia depositada una cadena perfectamente lineal da Cs/Co $= 0.980$, no 1. $c_0$ es la ordenada de un ajuste de dos puntos y dos parametros, asi que pasa exacto: lo unico que dice algo es su error.\end{minipage}
}
\end{table}

El desacuerdo va de \DesacMin\,\% a \DesacMax\,\%, con el error estadístico de
cada $D$ en el \ErrRelMediano\,\%: es un efecto real, no ruido. Tres cosas que
conviene mirar juntas:

\begin{itemize}
\item \textbf{El signo cambia de detector a detector.} En A y D el Cs-137 da
      una ganancia mayor que el Co-60; en C, menor. Una no linealidad de la
      cadena de lectura, o un sesgo del método de ajuste, tendrían el mismo
      signo en todos.
\item \textbf{Crece con el voltaje en el detector A.} Ahí el desacuerdo es
      monótono con el alto voltaje (tabla~\ref{tab:testA}), mientras que en C
      se mantiene parejo. Un efecto que crece con la ganancia apunta a la
      cadena de lectura; uno que no, al punto de operación.
\item \textbf{El detector D desacuerda aun midiendo el mismo día.} Es el único
      punto donde las dos fuentes se tomaron en la misma sesión, con 45 minutos
      de diferencia, y aun así difieren en dos de los tres voltajes. Ahí no
      hay deriva entre sesiones que culpar.
\end{itemize}

La figura~\ref{fig:cociente} resume las tres columnas de la derecha de la
tabla en un solo gráfico, y la figura~\ref{fig:porpunto} muestra de dónde sale
ese número sin pasar por el ajuste: los dos espectros de cada punto,
calibrados con la \emph{misma} $D$, y las marcas de dónde deberían caer los
bordes.

\begin{figure}[htbp]
\centering
\includegraphics[width=0.72\linewidth]{cociente_fuentes.pdf}
\caption{El cociente entre las dos calibraciones de cada punto. La línea
cortada es lo que daría una cadena lineal en escala de energía depositada
($\CocEsperado$, el quenching de Birks), no el 1 de la línea punteada.}
\label{fig:cociente}
\end{figure}

\begin{figure}[htbp]
\centering
\includegraphics[width=\linewidth]{comparacion_por_punto.pdf}
\caption{Mismo detector, mismo voltaje. Cada espectro está dividido por su
propio plateau Compton (la mediana de sus cuentas en la mitad inferior de su
ventana de ajuste, la zona plana de antes del borde), así que esa zona vale 1
en las dos curvas y las alturas no comparan nada: sólo importan las posiciones.
Los dos espectros se llevan al eje de energía con una única $D$ (la del ajuste
conjunto por el origen). La línea punteada gris marca dónde debería caer cada
borde ($E$); la llena, del color de la fuente, dónde cae el borde que esa
fuente midió, leído con la escala común: $E\,D_{\mathrm{fuente}}/D_{\mathrm{conjunta}}$.
Coinciden sólo si esa fuente comparte la $D$ conjunta, y como la conjunta es
un compromiso entre las dos, cuando las fuentes difieren una línea se corre
hacia un lado y la otra hacia el otro; ese corrimiento es el desacuerdo
(sección~\ref{sec:leer-porpunto}). Si cada espectro se calibrara con su propia
$D$, los dos bordes caerían sobre sus marcas por construcción y la figura no
diría nada. Un panel con fondo crema sería ese caso: un punto de operación con
una sola fuente, cuyo borde cae sobre la marca por construcción y no es un test.
Hasta el 22 de septiembre de 2026 la fila del detector B era así; con su Co-60
medido, las cuatro filas tienen ahora las dos fuentes y las cuatro son tests.}
\label{fig:porpunto}
\end{figure}

\subsection{Cómo se lee la figura~\ref{fig:porpunto}}
\label{sec:leer-porpunto}

Es la figura menos obvia del informe y la que más dice, así que conviene
explicarla despacio:

\begin{itemize}
\item Cada panel es un punto de operación (un detector a un voltaje). Las dos
      curvas son los espectros medidos con Cs-137 (naranja) y Co-60 (azul).
\item \textbf{Qué es el eje vertical.} Cada curva está dividida por su propio
      plateau Compton: la mediana de sus cuentas en la mitad inferior de su
      ventana de ajuste, que es la zona plana que precede al borde. Así esa
      zona vale 1 en las dos curvas, y las alturas no comparan ni actividades
      ni tiempos de medida ni eficiencias: sólo importan las posiciones. Sin
      esa división el Co-60, que reparte cuentas parecidas en el doble de
      energía, quedaría aplastado contra el eje y no se podría comparar dónde
      cae cada borde. El pico de bajo canal (el exceso, fuera de la ventana de
      ajuste) se sale del eje a propósito.
\item Los dos espectros se llevan al eje de energía con la \emph{misma} $D$,
      la del ajuste conjunto de ese punto: la recta por el origen que mejor
      pasa por los bordes de las dos fuentes a la vez. Ésa es la clave: si cada
      uno se calibrara con la suya, los dos bordes caerían sobre sus marcas
      por construcción y no habría nada que ver.
\item Las líneas punteadas grises marcan dónde \emph{deberían} caer los bordes
      (\BordeCs, \BordeCoA{} y \BordeCoB{} keV). Las líneas llenas, del color
      de cada fuente, marcan dónde cae el borde que esa fuente midió: el canal
      del borde ajustado, $D_{\mathrm{fuente}}\,E$, leído con la escala común,
      o sea en $E\,D_{\mathrm{fuente}}/D_{\mathrm{conjunta}}$.
\item \textbf{Por qué la mayoría no coincide con su marca, y por qué se
      corren en direcciones opuestas.} Una línea llena cae sobre su punteada
      sólo si $D_{\mathrm{fuente}} = D_{\mathrm{conjunta}}$. Pero la conjunta
      es un compromiso entre las dos fuentes: si difieren, una queda por
      encima y la otra por debajo, y el corrimiento relativo de cada línea es
      $D_{\mathrm{fuente}}/D_{\mathrm{conjunta}} - 1$. Cuando el Cs-137 pide
      más \mVns{} por keV que el Co-60 (detectores A, B y D) su borde queda a
      la derecha de \BordeCs{} keV y los dos del Co-60 a la izquierda de
      \BordeCoA{} y \BordeCoB; en el detector C es al revés. Los dos
      corrimientos, uno para cada lado, componen el cociente Cs/Co que muestra
      el panel, y el segundo número (\emph{desac.}) es ese cociente con el
      quenching de Birks descontado, la misma columna de la
      tabla~\ref{tab:testA}. Las líneas coincidirían sólo si las dos fuentes
      dieran la misma $D$, que es la hipótesis que el test A rechaza. Con una
      cadena perfectamente lineal tampoco coincidirían \emph{exactamente}: en
      escala de energía depositada el quenching pone Cs/Co en \CocEsperado, o
      sea un corrimiento de $\sim 1$\,\% por línea (5 keV en el Cs-137, 10 en
      el Co-60), invisible a esta escala. Lo que la figura muestra son
      desacuerdos de \DesacMin{} a \DesacMax\,\%. El panel D @ 1700~V, donde
      Cs/Co vale casi 1, es el único en que las llenas caen sobre las
      punteadas.
\end{itemize}

\subsection{Control: no es la compuerta}

La compuerta común es la única decisión de análisis que se le impone a todas
las corridas por igual, así que es el primer sospechoso. Y es un buen control
porque la palanca es larga: 75 ns contra \PostFijo{} ns es multiplicar por seis
la ventana, y pasar de recoger tres cuartos de la carga a recogerla casi
entera. Aun así, donde el registro alcanza para las dos ventanas el cociente se
mueve menos de \ControlMax\,\% (tabla~\ref{tab:control}), un orden de magnitud
menos que el desacuerdo que hay que explicar. En el detector D el control no se
puede correr: con 75 ns entra el 2\,\% de los eventos.

\begin{table}[htbp]
\centering
\caption{El mismo test con otra compuerta, donde el registro lo permite.}
\label{tab:control}
{\small\begin{tabular}{lrrr}
\toprule
\textbf{Punto} & \textbf{Cs/Co con 12 ns} & \textbf{Cs/Co con 75 ns} & \textbf{Cambio [\%]} \\
\midrule
A @ 1600 V & 1.0750 & 1.0884 & $+1.24$ \\
A @ 1650 V & 1.1033 & 1.1092 & $+0.53$ \\
A @ 1700 V & 1.1238 & 1.1337 & $+0.88$ \\
B @ 1600 V & 1.0682 & 1.0796 & $+1.07$ \\
B @ 1650 V & 1.0663 & 1.0748 & $+0.80$ \\
B @ 1700 V & 1.0765 & 1.0756 & $-0.08$ \\
C @ 1600 V & 0.8693 & 0.8762 & $+0.79$ \\
C @ 1650 V & 0.9015 & 0.9100 & $+0.93$ \\
C @ 1700 V & 0.8908 & 0.9013 & $+1.18$ \\
\bottomrule
\end{tabular}
\par\smallskip
\begin{minipage}{0.95\linewidth}\footnotesize Donde el registro alcanza para las dos compuertas, el cociente entre fuentes se mueve menos del 1\%: la eleccion de la compuerta no explica el desacuerdo del test A. En el detector D el control no se puede correr, porque con 75 ns entra el 2\% de los eventos.\end{minipage}
}
\end{table}

\subsection{Lectura}

Con estos datos \textbf{el eje canal--energía no se puede dar por lineal}. El
caso más claro es el del detector D: sus dos fuentes se midieron en la misma
sesión y con el detector ya en buenas condiciones (sección~\ref{sec:repro}), así
que ahí no hay ni deriva entre sesiones ni falla de montaje que culpar. En A,
en B y en C las dos fuentes vienen de sesiones distintas, así que a una posible
no linealidad se le suma una posible deriva del punto de operación, y con estos
datos las dos causas no se pueden separar. El detector B agrega un dato que
acota la discusión: su desacuerdo es del mismo tamaño que el del A (cerca del
9\,\%) pero, a diferencia del A, \emph{no crece con el voltaje}
(figura~\ref{fig:cociente}). Una no linealidad del eje no tiene por qué
depender del alto voltaje, pero una diferencia de ganancia entre las dos
sesiones tampoco la tendría si es fraccional y fija, así que el B es
compatible con esa segunda explicación de la manera más simple posible: un
factor constante, equivalente a menos de un 1\,\% de HV. Mientras eso siga abierto, la
calibración de una sola fuente vale para esa fuente y ese día, y extrapolarla
a otra energía arrastra un error del orden del 10\,\%.

\begin{aviso}
Conviene ser explícito sobre qué \emph{no} dice este resultado. No dice que el
método esté mal: el mismo método, sobre las mismas corridas, da $D$ estables
al \ErrRelMediano\,\% y una ley de ganancia limpia (sección~\ref{sec:detectores}).
Tampoco dice cuál de las dos fuentes tiene razón: el test es simétrico, mide
que discrepan y no cuál está corrida. Para resolverlo hace falta medir, no
analizar mejor.
\end{aviso}

% ===========================================================================
\section{Comparación entre detectores}
\label{sec:detectores}

\subsection{La ganancia y su ley con el voltaje}

La ganancia crece con el alto voltaje como una potencia, que es lo que se
espera de un fotomultiplicador, y no como una recta: el $\chi^2$ de la recta es
entre uno y dos órdenes de magnitud peor en los \NGruposTestB{} grupos
(tabla~\ref{tab:testB}). Los exponentes caen entre \BMin{} y \BMax{}, dentro o
al borde del rango 6 a 10 típico de un PMT.

Por qué una potencia y no una recta: cada dínodo multiplica los electrones por
un factor que depende del voltaje entre etapas, y con $n$ dínodos en cascada la
ganancia total va como ese factor a la $n$. De ahí sale $D \propto V^{b}$ con
$b$ del orden del número de etapas. Que el exponente medido caiga en ese rango
es una verificación independiente de que la cadena se comporta como se espera.

\begin{table}[htbp]
\centering
\caption{Ganancia contra alto voltaje, detector por detector y fuente por
fuente.}
\label{tab:testB}
{\small\begin{tabular}{llrrrrr}
\toprule
\textbf{Det.} & \textbf{Fuente} & \textbf{$n_V$} & \textbf{$b$} & \textbf{$\chi^2$ potencia} & \textbf{$\chi^2$ recta} & \textbf{Peor residuo [\%]} \\
\midrule
A & Co-60 & 4 & $9.37 \pm 0.03$ & 53 & 1707 & 1.1 \\
A & Cs-137 & 3 & $10.42 \pm 0.04$ & 0 & 358 & 0.1 \\
B & Co-60 & 3 & $9.44 \pm 0.04$ & 5 & 244 & 0.3 \\
B & Cs-137 & 3 & $9.56 \pm 0.05$ & 0 & 192 & 0.1 \\
C & Co-60 & 3 & $8.70 \pm 0.05$ & 168 & 638 & 2.1 \\
C & Cs-137 & 3 & $9.34 \pm 0.05$ & 4 & 245 & 0.3 \\
D & Co-60 & 3 & $10.10 \pm 0.04$ & 493 & 1749 & 2.9 \\
D & Cs-137 & 3 & $7.73 \pm 0.05$ & 22 & 31 & 0.6 \\
\bottomrule
\end{tabular}
\par\smallskip
\begin{minipage}{0.95\linewidth}\footnotesize Ajuste de $D(V)$ con las dos hipotesis. El $\chi^2$ de la recta es una a dos ordenes de magnitud mayor en todos los casos, o sea que la ley de potencia gana sin ambiguedad. Un PMT tipico tiene $b$ entre 6 y 10.\end{minipage}
}
\end{table}

\begin{figure}[htbp]
\centering
\includegraphics[width=\linewidth]{ganancia_dv.pdf}
\caption{$D(V)$ de los cuatro detectores en escala doble logarítmica, con la
ley de potencia ajustada, y los residuos de ese ajuste. En log-log una ley de
potencia es una recta y su pendiente es $b$. Las cruces del detector D son las
corridas repetidas de la sección~\ref{sec:repro}, que no entran en ningún
ajuste. La figura completa, con el exponente de cada grupo, está en
\texttt{figuras/resumen\_ganancia.png}.}
\label{fig:ganancia}
\end{figure}

\subsection{Cuánta ganancia tiene cada uno}

Los cuatro detectores no están ni cerca de la misma ganancia: a 1600 V con
Cs-137 hay un factor \RazonGanancias{} entre el mayor y el menor
($D = \DmilseisA$ el A, $\DmilseisB$ el B, $\DmilseisC$ el C y $\DmilseisD$ el
D, en \DkeV). La última columna de la tabla~\ref{tab:detectores} traduce eso a
lo único que se puede accionar en el laboratorio: a qué voltaje habría que
poner cada detector para que iguale la ganancia que el A tiene a 1600 V, usando
su propia ley $D \propto V^{b}$.

\begin{table}[htbp]
\centering
\caption{Comparación entre detectores y voltaje equivalente.}
\label{tab:detectores}
{\small\begin{tabular}{lrrrrrr}
\toprule
\textbf{Det.} & \textbf{1550 V} & \textbf{1600 V} & \textbf{1650 V} & \textbf{1700 V} & \textbf{$D/D_A$} & \textbf{$V_{eq}$} \\
\midrule
\midrule
\multicolumn{7}{l}{\textbf{Cs-137}} \\
A & -- & 1.3563 & 1.8710 & 2.5502 & 1.000 & 1600 \\
B & -- & 1.0717 & 1.4364 & 1.9140 & 0.790 & 1640 \\
C & -- & 1.2055 & 1.5989 & 2.1246 & 0.889 & 1620 \\
D & -- & 0.4934 & 0.6323 & 0.7877 & 0.364 & 1824 \\
\midrule
\multicolumn{7}{l}{\textbf{Co-60}} \\
A & 0.9558 & 1.2617 & 1.6958 & 2.2694 & 1.000 & 1600 \\
B & -- & 1.0033 & 1.3472 & 1.7780 & 0.795 & 1639 \\
C & -- & 1.3868 & 1.7735 & 2.3850 & 1.099 & 1583 \\
D & -- & 0.4222 & 0.5538 & 0.7863 & 0.335 & 1783 \\
\bottomrule
\end{tabular}
\par\smallskip
\begin{minipage}{0.95\linewidth}\footnotesize Las columnas de voltaje son $D$ en mV$\cdot$ns/keV a ese alto voltaje; $D/D_A$ compara con el detector A a 1600 V. $V_{eq}$ responde la pregunta practica: a que voltaje hay que poner cada detector para que tenga la ganancia que el detector A tiene a 1600 V, usando su propia ley $D \propto V^{b}$. Extrapolar fuera de 1600--1700 V es justamente eso, extrapolar.\end{minipage}
}
\end{table}

\begin{figure}[htbp]
\centering
\includegraphics[width=\linewidth]{comparacion_detectores.pdf}
\caption{Comparación entre detectores. Arriba: la ganancia de cada uno voltaje
por voltaje, y la misma ganancia relativa al detector A. Abajo: la resolución
en el borde de referencia y la calidad de los \NCorridas{} ajustes.}
\label{fig:detectores}
\end{figure}

\subsection{Resolución}

La resolución en el borde de referencia va de \FwhmMin\,\% a \FwhmMax\,\%
FWHM (panel inferior izquierdo de la figura~\ref{fig:detectores}). Dos cosas
que mirar ahí: si depende del alto voltaje y si separa a los detectores. Lo
primero no debería, si la resolución está dominada por la estadística de
fotoelectrones y no por el ruido electrónico; lo segundo es la comparación que
importa, y es limpia porque las cuatro se midieron con la misma cadena y el
mismo análisis.

Hay que leerla, eso sí, junto con la compuerta: integrar sólo \PostFijo{} ns
deja afuera buena parte de la carga (sección~\ref{sec:decisiones}), y menos
señal por evento es peor resolución. El número que sale de acá no es la
resolución intrínseca del centellador, es la de este observable.

% ===========================================================================
\section{El detector D con la tierra mal conectada}
\label{sec:repro}

El detector D tiene, en cada voltaje, dos corridas de Cs-137 tomadas con una
semana de diferencia (1 y 8 de septiembre), y no son dos intentos equivalentes:
\textbf{en la sesión del 1 de septiembre el detector estaba mal, con la tierra
mal conectada}. Las corridas válidas del detector D son las del 8 de
septiembre, y son las que se usan en todo el resto del informe. Las del 1-sep
se calibran igual, y aparecen acá, porque miden el tamaño de esa falla.

\begin{table}[htbp]
\centering
\caption{El mismo punto y la misma fuente, con y sin la falla de tierra.}
\label{tab:repro}
{\small\begin{tabular}{llrlrlr}
\toprule
\textbf{Punto} & \textbf{Fuente} & \textbf{$D$ (sesion valida)} & \textbf{Medida} & \textbf{$D$ (tierra mal)} & \textbf{Medida} & \textbf{Cambio [\%]} \\
\midrule
D @ 1600 V & Cs-137 & $0.4934 \pm 0.0011$ & 2026-09-08 & $0.3599 \pm 0.0004$ & 2026-09-01 & $-27.1$ \\
D @ 1650 V & Cs-137 & $0.6323 \pm 0.0010$ & 2026-09-08 & $0.4668 \pm 0.0010$ & 2026-09-01 & $-26.2$ \\
D @ 1700 V & Cs-137 & $0.7877 \pm 0.0019$ & 2026-09-08 & $0.6171 \pm 0.0010$ & 2026-09-01 & $-21.7$ \\
\bottomrule
\end{tabular}
\par\smallskip
\begin{minipage}{0.95\linewidth}\footnotesize Mismo detector, mismo voltaje nominal y misma fuente. La sesion del 1-sep se tomo con la tierra del detector mal conectada y esta descartada; la del 8-sep es la valida. La columna del cambio mide, entonces, cuanto movio la ganancia esa falla, no la reproducibilidad del montaje.\end{minipage}
}
\end{table}

La ganancia se corrió entre \RepMin\,\% y \RepMax\,\% (tabla~\ref{tab:repro}),
cuando el error estadístico de cada $D$ es del orden del \ErrRelMediano\,\%. O
sea que una tierra mal conectada mueve la calibración cien veces más que la
estadística de la medición, y lo hace sin romper nada visible: las corridas del
1-sep ajustan, dan un $\chi^2$ razonable y entregan un número perfectamente
formateado. Es el mejor argumento del informe a favor de anotar el estado del
montaje en cada sesión: el análisis no tiene forma de detectar esto solo.

\begin{aviso}
Dos lecturas que \textbf{no} corresponde hacer con estos números. Primero, no
son una medida de la reproducibilidad del montaje: una de las dos mediciones
está mal por una causa identificada, así que lo que comparan es «con falla»
contra «sin falla». Segundo, ese \RepMax\,\% no es una barra de error que haya
que sumarle a las demás corridas; es el tamaño de un problema puntual, ya
resuelto, del detector D.
\end{aviso}

Hay una consecuencia sobre el test de linealidad que conviene explicitar. El
detector D es el único punto de operación donde las dos fuentes se midieron en
la misma sesión, y ahora se sabe que esa sesión (la del 8-sep) es la buena. O
sea que su desacuerdo entre fuentes no se puede atribuir ni a deriva entre
sesiones ni a la falla de tierra: es la evidencia más limpia que tiene la
campaña de que el problema de la sección~\ref{sec:testA} es real. La otra cara
es que el detector D es el de menor ganancia y peor resolución, y aquel en que
el borde queda más cerca del umbral del trigger, así que también es donde más
puede pesar el exceso de bajo canal que el modelo no tiene.

% ===========================================================================
\section{Comparación con el método publicado}
\label{sec:paper}

El método que usa este informe es el de Kong \emph{et al.}~\cite{kong2025}:
selección automática del intervalo por correlación cruzada
normalizada y búsqueda del óptimo por recocido simulado. Esta sección compara
lo que obtenemos con lo que ellos publican, para saber dónde estamos parados.
Todos los números del paper que aparecen acá salen de una transcripción de sus
figuras que vive en \texttt{replica\_paper/auditoria.py}, con la referencia a la
figura de la que se leyó cada uno.

\subsection{Qué no es comparable}

Dos cosas, y conviene sacarlas del medio antes de comparar nada:

\begin{itemize}
\item \textbf{El eje.} Ellos calibran canal de ADC contra energía; nosotros
      área de pulso en \mVns{} integrada en una compuerta de \PostFijo{} ns.
      Sus pendientes y las nuestras son cantidades distintas de electrónicas
      distintas; que se parezcan numéricamente es casualidad.
\item \textbf{El detector.} Su validación de linealidad (su Fig. 3) usa
      LaBr$_3$:Ce y NaI:Tl, que son inorgánicos, de alto $Z$ y \textbf{tienen
      fotopico}. Eso les permite comparar la calibración por borde Compton
      contra la calibración por fotopico \emph{en el mismo espectro}, que es
      una referencia interna que nosotros no tenemos: en un plástico no hay
      fotopico, y por eso nuestro único control es medir dos fuentes distintas.
      Los centelladores plástico y líquido (BC408 y EJ309) aparecen en sus
      figuras 4 y 5, sin tabla publicada.
\end{itemize}

\subsection{Robustez frente al intervalo de ajuste}

Éste sí es el mismo test: tomar una corrida y ajustarla en seis intervalos
anidados, desde uno que apenas cubre el borde hasta uno que abarca casi todo el
espectro, y mirar cuánto se mueve el resultado. Los seis intervalos, en
unidades del canal del borde, se leyeron de los paneles de su Fig. 5 y se
aplican tal cual a nuestras corridas.

\begin{table}[htbp]
\centering
\caption{El test de la Fig. 5 del paper, sobre nuestras corridas y sobre la
suya.}
\label{tab:robustez}
{\small\begin{tabular}{llrrr}
\toprule
\textbf{Punto} & \textbf{Fuente} & \textbf{Intervalos} & \textbf{Dispersion [\%]} & \textbf{Sin el mas angosto [\%]} \\
\midrule
A @ 1600 V & Cs-137 & 5 de 6 & 1.0 & 1.0 \\
B @ 1600 V & Cs-137 & 5 de 6 & 49.5 & 1.4 \\
C @ 1600 V & Cs-137 & 5 de 6 & 1.6 & 1.6 \\
D @ 1600 V & Cs-137 & 4 de 6 & 0.8 & 0.5 \\
\midrule
\textit{paper, Fig. 5} & \textit{Cs-137} & 6 de 6 & 1.2 & 0.3 \\
\bottomrule
\end{tabular}
\par\smallskip
\begin{minipage}{0.95\linewidth}\footnotesize El mismo test de la Fig. 5 del paper: la misma corrida ajustada en seis intervalos anidados, definidos en unidades del canal del borde y tomados de sus propios paneles. Un intervalo con menos de 30 bines se saltea, igual que en la replica. La fila del paper sale de los seis valores que rotula su figura.\end{minipage}
}
\end{table}

Sacando el intervalo más angosto, que en varias corridas ni siquiera contiene
el flanco entero, nuestra dispersión va de \RobustezNuestraMin\,\% a
\RobustezNuestraMax\,\% contra el \RobustezPaperSinAngosta\,\% que sale de sus
propios paneles. Mismo orden de magnitud, un poco peor. Incluyendo el intervalo
más angosto la comparación se ensucia en los dos casos: a ellos la dispersión
les sube a \RobustezPaper\,\% y a nosotros mucho más en la corrida donde ese
intervalo es evaluable. La conclusión es la misma para los dos: el método es
estable mientras la ventana contenga el flanco completo, y deja de serlo cuando
no.

\subsection{Las métricas con las que el paper valida}

Acá aparece la diferencia importante, y no es de resultados sino de vara. El
paper valida su calibración con dos números: $R^2$ (entre \RdosPaperMin{} y
\RdosPaperMax{} en sus cuatro ajustes) y una similitud coseno entre las dos
relaciones energía-canal (\CosenoPaperMin{} y \CosenoPaperMax). La
tabla~\ref{tab:paper} aplica esas mismas métricas a nuestros puntos.

\begin{table}[htbp]
\centering
\caption{Nuestros puntos, medidos con las métricas de validación del paper.}
\label{tab:paper}
{\small\begin{tabular}{lrrrr}
\toprule
\textbf{Punto} & \textbf{Desacuerdo [\%]} & \textbf{SD [keV]} & \textbf{Similitud coseno} & \textbf{$R^2$} \\
\midrule
A @ 1600 V & $+9.7$ & 36.1 & 1.00000000 & 0.9713 \\
A @ 1650 V & $+12.6$ & 31.0 & 1.00000000 & 0.9787 \\
A @ 1700 V & $+14.7$ & 43.3 & 1.00000000 & 0.9539 \\
B @ 1600 V & $+9.0$ & 22.5 & 1.00000000 & 0.9896 \\
B @ 1650 V & $+8.8$ & 20.3 & 1.00000000 & 0.9917 \\
B @ 1700 V & $+9.9$ & 23.9 & 1.00000000 & 0.9882 \\
C @ 1600 V & $-11.3$ & 49.0 & 1.00000000 & 0.9716 \\
C @ 1650 V & $-8.0$ & 36.4 & 1.00000000 & 0.9829 \\
C @ 1700 V & $-9.1$ & 53.1 & 1.00000000 & 0.9663 \\
D @ 1600 V & $+19.3$ & 49.3 & 1.00000000 & 0.9360 \\
D @ 1650 V & $+16.5$ & 47.2 & 1.00000000 & 0.9427 \\
D @ 1700 V & $+2.2$ & 0.6 & 1.00000000 & 1.0000 \\
\bottomrule
\end{tabular}
\par\smallskip
\begin{minipage}{0.95\linewidth}\footnotesize Las dos ultimas columnas son las metricas con las que el paper valida su calibracion (su Fig. 3). La similitud coseno da 1 EXACTAMENTE en todos los puntos, incluido el que discrepa mas: entre dos rectas por el origen es 1 sean cuales sean las pendientes, asi que no puede ver este modo de falla. El $R^2$ se calcula sobre dos puntos y con la recta forzada por el origen; con ordenada libre y dos puntos daria 1 exacto.\end{minipage}
}
\end{table}

\begin{aviso}
La similitud coseno da \textbf{\CosenoNuestro{} exactamente} en los nueve
puntos, incluido aquel en que las dos fuentes discrepan un \DesacMax\,\%. No es
falta de precisión numérica: dos rectas que pasan por el origen son vectores
paralelos, así que su coseno vale 1 cualquiera sea la diferencia de pendientes.
La métrica es ciega, por construcción, al modo de falla que este informe
reporta. Con esa vara nuestra campaña se vería impecable.
\end{aviso}

El $R^2$ es algo más sensible (baja hasta \RdosNuestroMin{} en el peor punto)
pero tampoco alcanza: leído sin contexto, un $R^2$ de 0.94 sobre dos puntos no
parece un problema, y sin embargo corresponde a una discrepancia del
\DesacMax\,\% entre fuentes. Por eso este informe reporta residuos en keV y en
por ciento, que es lo que hace visible el desacuerdo.

En esa escala: el paper cita desviaciones de \SDPaperMin{} a \SDPaperMax{} keV
para sus rectas, con 5 o 6 puntos de cuatro fuentes; nuestros puntos de
operación dan entre \SDNuestroMin{} y \SDNuestroMax{} keV con dos puntos. La
comparación es desfavorable para nosotros y hay que tomarla con la salvedad de
que no es simétrica: con dos puntos y una recta forzada por el origen, un
desacuerdo de pendiente aparece entero como residuo, mientras que con seis
puntos y ordenada libre se reparte.

\subsection{Qué encontramos al auditar los números publicados}

Reproducir sus resultados con sus propios datos (los puntos de la Fig. 3 están
rotulados en la figura, así que se leen sin digitalizar píxeles) deja cuatro
observaciones. Están en \texttt{replica\_paper/} con el detalle:

\begin{enumerate}
\item En el recuadro de la Fig. 3, el SD y el $R^2$ están intercambiados entre
      las dos rectas. Al corregirlo, la afirmación de su Sección III de que el
      SD del NCC-SA se mantiene por debajo del ajuste de fotopico \textbf{no se
      cumple} en ninguno de los dos detectores.
\item La Ec. 1 impresa da la energía del fotón retrodispersado, no la del borde
      Compton: está mal por un factor 2.6 a 5.2. Sus propios rótulos siguen la
      versión correcta, así que es una errata de tipografía y no de cálculo.
\item La Ec. 12 define $\Delta E$ como la diferencia entre líneas gamma; la
      separación entre los dos bordes del Co-60 que muestra su Fig. 3
      corresponde a la diferencia entre \emph{bordes}. Con la fórmula impresa,
      el segundo borde queda 4.2 keV corrido, contra un SD de ajuste de 2.2 keV.
\item Ninguna lectura razonable de la similitud coseno reproduce los dos
      valores publicados, y el paper no define entre qué vectores la calcula.
\end{enumerate}

Nada de esto invalida el método, que es el que usamos; sí quiere decir que sus
números publicados no sirven como vara contra la cual medir los nuestros.

\subsection{Dónde queda nuestra metodología}

\begin{itemize}
\item \textbf{Lo que hacemos igual o mejor.} La selección automática del
      intervalo y el optimizador son los mismos (el paper no publica sus $\rho$,
      así que ahí no hay comparación posible; los nuestros van de \RhoMin{} a
      \RhoMax). La plantilla, en cambio, es una simulación de Geant4 de la
      geometría real y no la lineshape analítica de Klein-Nishina: eso
      incorpora el scattering múltiple y el quenching de Birks, que en un
      plástico sesgan la posición del borde varios por ciento si se los ignora.
      También acá se verificó que el óptimo del ajuste es global
      (sección~\ref{sec:ajuste}), cosa que el paper afirma pero no muestra.
\item \textbf{Lo que hacemos con más cuidado.} Reportamos residuos en keV en
      lugar de similitud coseno, y decimos qué fracción de la carga entra en la
      compuerta. Las dos cosas hacen que nuestros números se vean peor y que
      digan más.
\item \textbf{Lo que nos falta.} Fuentes. Ellos validan con cuatro (Na-22,
      Cs-137, Mn-54 y Co-60) y nosotros con dos, y en la mayoría de los puntos
      medidas en sesiones distintas. Mientras eso siga así, nuestro test de
      linealidad detecta que hay un problema pero no puede decir de qué es.
\end{itemize}

% ===========================================================================
\section{Cómo usar estos números}
\label{sec:uso}

\subsection{La receta}

Para convertir un pulso medido en energía:

\begin{enumerate}
\item Integrar el pulso \textbf{de la misma manera}: onset al \FracOnset\,\%
      de la amplitud, línea base de las muestras anteriores al onset, integral
      desde unos \PreMin{} ns antes del onset hasta \PostFijo{} ns después
      (sección~\ref{sec:area}). Una $D$ de este informe sólo vale para áreas
      medidas así.
\item Buscar en la tabla~\ref{tab:calib} la fila del detector y el voltaje
      correspondientes.
\item Multiplicar: energía $=$ área $\times\ (1/D)$, con el área en \mVns{} y
      la energía en \keVdep.
\end{enumerate}

Ejemplo con el detector A a 1600 V, fuente Cs-137: $D = \DmilseisA$, así que un
pulso de 1000 \mVns{} corresponde a \EjemploKeV{} \keVdep. La barra de error de
$D$ (\ErrRelMediano\,\% típico) se propaga directamente a la energía.

\subsection{Cuándo NO vale}

\begin{itemize}
\item \textbf{Otro voltaje.} La ganancia va como $V^{b}$ con $b$ entre \BMin{}
      y \BMax: 50 V arriba son un \SubidaCincuenta\,\% más de ganancia. Se
      puede interpolar con la ley de la tabla~\ref{tab:testB}, pero es una
      extrapolación del modelo, no una medición.
\item \textbf{Otra ventana de integración.} Una compuerta más larga recoge más
      carga y da un $D$ mayor; la columna \emph{Carga} dice cuánto falta en
      cada caso.
\item \textbf{Otra energía muy distinta de la del borde usado.} Es justamente
      lo que el test de linealidad no pudo confirmar; hoy hay que contar con un
      error del orden del 10\,\% al extrapolar de una fuente a otra.
\item \textbf{Otra escala.} Estos números dan energía \emph{depositada}. Si lo
      que se quiere es luz (keVee), hay que recalibrar con la plantilla
      \texttt{Lvis\_keVee}; el cociente entre las dos escalas es el quenching
      de Birks, del orden del 5\,\%.
\item \textbf{Otra atenuación.} Las 25 corridas se tomaron con 20 dB en línea y
      el área se mide a la entrada del digitalizador, así que estos $D$ valen
      para esa cadena. Sacar el atenuador multiplica el área por diez.
\end{itemize}

% ===========================================================================
\section{Limitaciones y qué falta medir}
\label{sec:limites}

\begin{enumerate}
\item \textbf{Una sola fuente en el detector B: medido.} Era la medición más
      barata que faltaba y se hizo el 22 de septiembre de 2026: tres corridas
      de Co-60 a 1600, 1650 y 1700 V. Con ellas el detector B pasa a tener las
      dos fuentes y es el único de los cuatro cuyas dos fuentes dan el mismo
      exponente de ganancia ($b = 9{,}44 \pm 0{,}04$ con Co-60 y
      $9{,}56 \pm 0{,}05$ con Cs-137, compatibles a $1{,}9\,\sigma$; en los
      otros tres la diferencia va de 9 a 37 $\sigma$). Lo que sigue abierto es
      el punto siguiente: ese Co-60 es de otra sesión que su Cs-137, así que el
      desacuerdo entre fuentes no se puede atribuir todavía a una sola causa.

      \begin{aviso}
      El 10 de septiembre aparecieron tres archivos \texttt{Co60.dat} en las
      carpetas del detector B, y \textbf{no son del detector B}: son copias
      byte a byte de los del detector C (mismo tamaño, misma fecha y mismo
      contenido), y en el disco original el detector B sólo tiene Cs-137.
      Quedaron fuera del análisis. Calibrarlos habría dado tres números
      perfectamente formateados del detector equivocado, que es la clase de
      error que no falla en ningún lado. \texttt{rutas.duplicados()} ahora
      revisa esto y avisa antes de calibrar.
      \end{aviso}

\item \textbf{Las dos fuentes en la misma sesión, en los cuatro detectores.}
      Es la única forma de separar la deriva entre sesiones del comportamiento
      del detector, y hoy sólo el detector D la tiene. En A, en B y en C el
      Co-60 y el Cs-137 se midieron con días o semanas de diferencia, así que
      su desacuerdo no se puede atribuir a una sola causa. El detector B lo
      deja ver con particular claridad: su desacuerdo es del mismo tamaño que
      el del A pero \emph{no depende del voltaje} (tabla~\ref{tab:testA}),
      que es la firma de una diferencia de ganancia fija entre las dos
      sesiones y no de una no linealidad del eje.

\item \textbf{El atenuador: cerrado.} Era una limitación de este informe
      —no había registro de la atenuación de cada corrida— y dejó de serlo el
      22 de septiembre de 2026: son \textbf{20 dB en las 25 corridas},
      confirmado por quien midió, y están puestos para que el pulso del
      fotomultiplicador no pase el rango de entrada del digitalizador. Con las
      amplitudes medidas del Co-60 a 1550 V (mediana 45 mV, máximo en el techo
      de 500 mV) eso no es una precaución nominal: sin atenuar, casi la mitad
      de los eventos habría saturado. Dos consecuencias para leer las tablas.
      Primera, el área se integra a la \emph{entrada} del digitalizador y
      ningún paso del análisis corrige por la atenuación, así que todo $D$ de
      este informe es por \mVns{} \emph{en esa entrada}, y la carga a la salida
      del tubo es diez veces mayor. Segunda, lo que antes era un argumento
      indirecto —que las cuatro sesiones cerraran razonablemente sugería que la
      atenuación no había cambiado— pasa a ser una comprobación cruzada que
      coincide con el registro: el punto de junio cae sobre la ley
      $D \propto V^{b}$ a menos del 0,1 \%, y un cambio de atenuador lo habría
      corrido un factor diez.

\item \textbf{El retardo del trigger, más a la izquierda.} En la sesión del 8
      de septiembre el pulso completo no entra en el registro: después del
      onset sobran unos 46 ns. Con la compuerta corta de este informe eso no
      molesta, pero deja sin opción a cualquier análisis que quiera el pulso
      entero, y sin el pulso entero no hay con qué medir qué fracción de la
      carga se está integrando de verdad. Adelantando el trigger 40 a 50 ns se
      recupera esa posibilidad.

\item \textbf{El exceso de bajo canal no está modelado.} Es la causa conocida
      de los $\chi^2$/ndf por encima de 2. Hay dos caminos: volver a modelarlo
      (lo que agrega dos parámetros) o subir \texttt{inicio\_rel} para dejarlo
      afuera de la ventana. Hoy está en \InicioRel{} y el efecto sobre $D$ de
      moverlo es del orden del 1\,\%.

\item \textbf{El $\chi^2$ usa las barras del dato.} Pesar con la incerteza
      medida en lugar de la predicha por el modelo sesga $D$ hacia abajo, y el
      sesgo va como el inverso del número de cuentas por bin. Con la
      estadística de esta campaña (mediana de \CuentasBin{} cuentas por bin
      dentro de la ventana) el efecto está muy por debajo de los demás
      sistemáticos, pero está: medido con verdad conocida
      (sección~\ref{sec:estadistica}, tabla~\ref{tab:cierre}), pesar con el
      dato da un sesgo de $\CierreSesgoNeyman \pm \CierreErrSesgoNeyman$\,\%
      contra $\CierreSesgoPearson \pm \CierreErrSesgoPearson$\,\% pesando
      con el modelo.

\item \textbf{Na-22 del detector C.} Está medido (1600 V) y no se analizó: no
      hay simulación de Geant4 de esa fuente con esta geometría, y la única
      ruta sería la plantilla analítica, que sesga $D$ un 4\,\% respecto de las
      demás. Un punto con ese sesgo metido en la recta de linealidad mediría la
      elección de plantilla, no el detector. Con una simulación de Na-22, el
      detector C pasaría a tener cuatro bordes y sería el mejor test de
      linealidad de la campaña.

\item \textbf{El fondo.} \NSinFondo{} de las \NCorridas{} corridas se calibran
      con el espectro crudo. Para la posición del borde el efecto es chico,
      pero una corrida de fondo por punto de operación cerraría la duda.
\end{enumerate}

% ===========================================================================
% Las referencias: se citan con \cite{clave} y viven en referencias.bib.
% `unsrt` las numera en orden de aparición.
\bibliographystyle{unsrt}
\bibliography{referencias}

% ===========================================================================
\appendix
\section{Glosario de símbolos}
\label{ap:glosario}

\begin{center}
\begin{tabular}{lll}
\toprule
\textbf{Símbolo} & \textbf{Unidad} & \textbf{Qué es} \\
\midrule
área & \mVns & la integral del pulso; el observable que se calibra \\
onset & ns & cruce del \FracOnset\,\% de la amplitud; ancla de la compuerta \\
compuerta & ns & de \emph{pre} antes del onset a \PostFijo{} después; se integra ahí \\
$D$ & \DkeV & la calibración: área por unidad de energía depositada \\
$1/D$ & keV/\mVns & lo mismo dado vuelta; es lo que se usa para convertir \\
$E_\gamma$ & keV & energía de la línea gamma de la fuente \\
$E_{\max}$ & keV & energía del borde Compton de esa línea \\
$E_{\text{ref}}$ & keV & el borde de la primera línea; ancla del mapa \\
$B$ & \mVns & canal donde cae el borde de referencia; $B = D\,E_{\text{ref}}$ \\
$\sigma_{\text{rel}}$ & --- & resolución, como fracción de $B$ \\
FWHM & \% & ancho a media altura, $2.35\,\sigma_{\text{rel}}$ \\
$A$ & cuentas & amplitud del modelo; se despeja, no se ajusta \\
$\rho$ & --- & correlación de Pearson entre dato y plantilla (NCC) \\
$\chi^2$/ndf & --- & bondad del ajuste; cerca de 1 es bueno \\
ndf & --- & grados de libertad: bines de la ventana menos parámetros \\
$b$ & --- & exponente de la ley de ganancia, $D \propto V^{b}$ \\
$c_0$ & \mVns & ordenada del mapa cuando se la libera; debería ser cero \\
\keVdep & keV & energía depositada en el centellador \\
keVee & keV & energía electrón-equivalente (luz); \emph{no} se usa acá \\
$\chi^2$ de Neyman & --- & pesado con la varianza \emph{observada}; el de la campaña \\
$\chi^2$ de Pearson & --- & pesado con la varianza \emph{predicha}; sin sesgo \\
perfil & --- & $\Delta\chi^2$ contra un parámetro, reoptimizando el resto \\
cobertura & --- & cuántas veces la verdad cae dentro de la barra (68\,\%) \\
VARPRO & --- & la amplitud $A$ se despeja, no se busca \\
\bottomrule
\end{tabular}
\end{center}

\section{Qué archivo es cada cosa}
\label{ap:archivos}

\begin{center}
\begin{tabular}{ll}
\toprule
\textbf{Archivo} & \textbf{Qué tiene} \\
\midrule
\multicolumn{2}{l}{\textit{Resultados}} \\
\texttt{figuras/calibraciones/resultados.csv} &
   una fila por corrida, para abrir en cualquier lado \\
\texttt{figuras/calibraciones/resultados.json} &
   lo mismo, completo, es lo que lee este informe \\
\texttt{figuras/calibraciones/RESUMEN.txt} &
   el informe en texto plano, con los dos tests \\
\midrule
\multicolumn{2}{l}{\textit{Figuras de resumen}} \\
\texttt{RESUMEN\_ganancia.png} & $D(V)$, residuos, exponentes y tabla \\
\texttt{RESUMEN\_linealidad.png} & test A: puntos, residuos, cocientes \\
\midrule
\multicolumn{2}{l}{\textit{Por corrida}, en
   \texttt{Detector X/<voltaje>/<fuente>/}} \\
\texttt{\ldots\_1\_calibracion.png} & el ajuste y el espectro calibrado \\
\texttt{\ldots\_2\_parametros.png} & $\Theta$, errores y correlaciones \\
\texttt{\ldots\_3\_espectro\_y\_simulacion.png} & qué entra al ajuste \\
\texttt{\ldots\_4\_ncc\_e\_intervalo.png} & el barrido del NCC y la ventana \\
\texttt{\ldots\_5\_recocido.png} & los \NArranques{} arranques y el camino \\
\texttt{\ldots\_6\_ventana\_de\_integracion.png} &
   el pulso medio y la compuerta \\
\bottomrule
\end{tabular}
\end{center}

\section{Cómo reproducir estos números}
\label{ap:reproducir}

Todo el informe sale de tres comandos, y ninguno de los números de este
documento está tipeado a mano:

\begin{quote}\ttfamily\small
python analisis/calibrar\_todo.py \\
python analisis/calibrar\_todo.py -{}-resumen -{}-control \\
python analisis/informe.py
\end{quote}

\noindent
El primero recorre el catálogo de corridas, calibra cada una y deja las seis
páginas de diagnóstico de cada ajuste en
\texttt{analisis/figuras/calibraciones/Detector X/<voltaje>/<fuente>/}, más
\texttt{resultados.json}, \texttt{resultados.csv} y un informe de texto. El
segundo agrega el control de compuerta de la tabla~\ref{tab:control}. El
tercero regenera las figuras y las tablas de este documento y lo compila.

Los parámetros que definen el análisis (la compuerta, el modelo con o sin
exceso, el arranque de la ventana, la escala) están en un solo lugar y se
pueden cambiar desde la interfaz: \texttt{python analisis/gui.py}. El lote los
importa de ahí, así que no pueden quedar dos definiciones distintas del mismo
observable.

Lo que se edita a mano es \texttt{informe/informe.tex} y
\texttt{informe/preambulo.tex}. Lo que hay en \texttt{informe/generado/} se
reescribe en cada corrida: si un número del texto tiene que cambiar cuando
cambien los datos, va como macro (la lista está al principio de
\texttt{generado/numeros.tex}) y no escrito en el párrafo.

\section{Ajustes de ejemplo}

Los tres detectores que faltan, todos a 1600 V. Los \NCorridas{} ajustes
completos, con sus seis páginas de diagnóstico cada uno, están en
\texttt{analisis/figuras/calibraciones/}.

\begin{figure}[H]
\centering
\includegraphics[width=\linewidth]{ajuste_B.pdf}\\[2mm]
\includegraphics[width=\linewidth]{ajuste_C.pdf}\\[2mm]
\includegraphics[width=\linewidth]{ajuste_D.pdf}
\caption{Ajustes de los detectores B, C y D a 1600 V. El D es el de menor
ganancia y peor resolución de los cuatro, y el que da los ajustes menos
estables.}
\end{figure}

\end{document}
